\chapter{自动驾驶场景}

\paragraph*{本章阅读锚点}
自动驾驶场景把 3DGS 推到一个很特殊的工作区：输入来自车载环视相机、鱼眼相机、LiDAR、雷达和连续行驶轨迹；输出不仅要能合成新视角，还要支持变道、停车、闭环仿真、传感器重仿真和下游感知训练。原始 3DGS 的短板在这里会被放大：训练视角沿车辆轨迹高度偏置，外推到旁车道或俯视角会崩；动态车辆和行人需要实例轨迹或自监督分解；夜间、鱼眼、LiDAR 和雷达又有完全不同的成像模型。本章按路线阅读：街景重建看静态背景和动态对象如何组合；动态驾驶看运动场、轨迹和实例分解；仿真与数据生成看 3DGS 如何变成闭环或数据机器；传感器融合看相机、LiDAR、雷达各自怎么渲染；占用感知看高斯如何从渲染表示变成 BEV 或 occupancy 的训练桥梁。

\section{街景重建与渲染}

街景重建的基础锚点是把驾驶场景分解为静态背景、动态前景和天空。与室内 3DGS 不同，车载数据沿道路单向采集，左右和下方视角严重缺失；动态车辆又常遮挡道路和建筑。代表性工作通常先解决“场景怎么分解”，后续方法再用 LiDAR、对象先验、扩散先验、鱼眼相机、夜间材质或停车位语义补足困难区域。



% METHOD: 2312.07920
\paragraph{DrivingGaussian（复合静动态高斯锚点）（2023.12）}
DrivingGaussian~\cite{2312_07920} 是动态自动驾驶 3DGS 的早期锚点之一。原始 3DGS 假设一个静态场；直接用于车载环视序列时，长距离背景、遮挡车辆和多相机一致性会互相干扰。DrivingGaussian 的核心改变是把场景拆成两个可组合部分：静态背景用增量式 3D Gaussians 逐段扩展，动态车辆和行人用独立对象高斯图建模，再按照 3D 边界框姿态合回世界坐标系。这样一来，背景不会被移动车辆反复污染，前景对象也能保留自己的局部几何和遮挡关系。

具体做法中，静态部分先利用 LiDAR 深度初始化第一个时间 bin 的高斯，再按行驶时间顺序处理后续 bin。新的场景片段不是从零训练，而是继承上一阶段高斯作为位置先验，并通过重叠区域对齐，把道路、建筑和远景逐步并入一个全局背景。动态部分为每个检测到的对象建立独立高斯模型，训练时用语义先验、边界框位姿和多相机图像监督对象外观；渲染时再把对象高斯按当前姿态放回静态背景。LiDAR 既用于初始化，也作为深度约束改善全景一致性。相对 Mega-NeRF 或单场 3DGS，DrivingGaussian 的关键不是换一个渲染器，而是把自动驾驶场景组织成“增量背景 + 可组合 actor”的图结构；后来的 Street Gaussians、AutoSplat 和多传感器方法基本都沿着这个分解思路继续加约束。

\textbf{实现细节：} Street Gaussians 虽然奠定了动态街景显式高斯表示的基础，但其一次性建模策略在远距离场景容易产生尺度混淆。DrivingGaussian（2023.12） 提出渐进式静态背景建模——按时间顺序逐步累加静态 3D 高斯——并引入复合动态高斯图处理多个运动物体，解决了这一局限。渐进式静态建模。直觉上，这如同画家从近到远逐层绘制风景——先画近处确定的物体建立锚点，再逐步扩展到远处模糊的背景。静态背景根据 LiDAR 深度范围均匀划分为 $N$ 个深度区间（bin），按时间顺序依次构建。第一个区间使用 LiDAR 先验初始化高斯参数（位置 $P(x,y,z)$、协方差 $\Sigma$、球谐颜色 $C(r,g,b)$、不透明度 $o$），后续区间以上一个区间的高斯位置为先验，基于重叠区域对齐。这种渐进策略避免了远距离场景一次性引入导致的尺度混淆和模糊伪影。复合动态高斯图。动态物体通过复合图 $H = \langle O, G_d, M, P, A, T \rangle$ 建模，其中每个节点存储一个实例物体 $o \in O$、对应动态高斯 $g_i \in G_d$、变换矩阵 $m_o \in M$、边界框中心和朝向 $(p_o, a_o)$。多个动态物体之间的遮挡通过基于距离的透明度调整处理：距离相机更近的物体获得更高的透明度，这符合光线传播原理。LiDAR 先验与训练。使用 LiDAR 点云替代 SfM 初始化，投影到多相机图像获取颜色，再通过密集束调整（Dense Bundle Adjustment, DBA）扩展到多相机设置来优化点位置。在 Waymo 上实现高保真环绕视角合成。
\textbf{锚点细化：} Street Gaussians 虽然奠定了动态街景显式高斯表示的基础，但其一次性建模策略在远距离场景容易产生尺度混淆。DrivingGaussian（2023.12）~\cite{2312_07920} 提出渐进式静态背景建模——按时间顺序逐步累加静态 3D 高斯——并引入复合动态高斯图处理多个运动物体，解决了这一局限。
\textbf{方法信息：} Google Research 参与了这篇论文；项目页为 \url{https://github.com/VDIGPKU/DrivingGaussian}，仓库约 402 stars，但当前公开内容以 README、权重和申请文档为主，未确认完整训练代码；Semantic Scholar 约 462 次引用。



% METHOD: 2403.20079
\paragraph{SGD（扩散 SDS 正则街景外推）（2024）}
SGD~\cite{2403_20079} 是早期用扩散先验增强街景外推的工作。它先在相邻帧图像上微调 Stable Diffusion，让模型学习场景内容；再通过 ControlNet 式模块加入 LiDAR depth，让扩散理解物体空间关系。训练 3DGS 时，对未见视角渲染施加 SDS 损失，引导未见区域生成合理内容，同时保留已见区域的多视角光度损失。相对 VEGS 的 LiDAR 初始化、法线监督和扩散蒸馏组合，SGD 的核心更简单：用场景特定扩散作为未见视角正则。
\textbf{实现细节：} MagicDrive3D（2024.05） 首先训练多视角视频生成模型再构建 3DGS，支持多条件控制。SGD（2024.03） 利用扩散先验和 LiDAR 深度增强 3DGS 的街景外推能力。StreetUnveiler（2024.05） 从拥挤观测中恢复空街道。FaithFusion（2025.11） 通过像素级期望信息增益融合 3DGS 和扩散模型。Drone-assisted Road GS（2024.08） 引入跨视图不确定性使无人机数据辅助地面视角合成。LR-SGS（2026.03） 利用 LiDAR 反射率引导显著高斯表示增强复杂光照场景。Omran et al.（2025.10） 提出高斯泼溅与扩散模型混合的两阶段新视角合成方案。Hybrid Gaussian（2025.03） 和 Hybrid Rendering（2025.03） 分别探索了混合表示和混合渲染策略。小结上，街景重建与渲染的核心突破在于鱼眼相机高斯变换和多模态统一渲染，使 3DGS 能同时建模针孔相机、鱼眼相机和 LiDAR 等多种传感器，为自动驾驶仿真提供完整的场景表示基础设施。
\textbf{方法信息：} Baidu 参与了这篇论文。
\textbf{设计动机与边界：} 相对原始 3DGS 以独立高斯和图像重建误差为主的设定，SGD 把问题落在以下可实现的改动上：SGD方法包含两个主要组件：（1）扩散模型微调：在驾驶场景数据上微调Stable Diffusion模型。采用两阶段训练策略：第一阶段，将相邻帧图像通过交叉注意力引入U-Net，使模型学习场景内容（「场景中有什么」）；第二阶段，通过ControlNet式模块引入LiDAR深度信息，使模型学习空间关系（「物体在哪里」）。（2）扩散模型引导的3DGS训练：将微调后的扩散模型作为正则化项，在训练中对3DGS在未见视角的渲染结果施加SDS（Score Distillation Sampling）损失。通过引导3DGS在未见区域生成合理内容，避免因训练视角稀疏导致的渲染退化。同时保持标准多视角光度损失用于已见区域的优化。LiDAR深度提供了精确的空间先验，帮助扩散模型理解场景的3D结构。其设计目的不是单独增加一个模块，而是用上述表示、监督或渲染约束消除原流程在目标任务中的具体失配。提出利用扩散模型先验增强3DGS在未见视角的渲染能力。通过在相邻帧图像和LiDAR深度数据上微调扩散模型，将其作为3DGS训练中未见视角的正则化。方法显著改善了偏离训练视角的渲染质量。适用时还应保留论文给出的边界：扩散模型推理增加训练开销；SDS损失可能引入扩散模型的偏差；推理时不需要扩散模型，但训练时间显著增加；对扩散模型微调质量敏感。
% METHOD: 2405.14475
\paragraph{MagicDrive3D（可控街景生成到 3DGS）（2024）}

MagicDrive3D~\cite{2405_14475} 走生成路线。第一阶段训练可控多视角视频生成模型，条件包括道路地图、3D 物体和文本描述，通过相对位姿控制相机运动；第二阶段把生成的多视角视频提升为 3DGS。由于生成视图会有微小几何错误，方法提出 Fault-Tolerant Gaussian Splatting，用单目深度初始化、外观建模和容错损失处理曝光与几何不一致。相对 RGE-GS 或 FaithFusion 对已有重建做修复，MagicDrive3D 更像从可控视频生成器出发构建可任意视角渲染的街景 3D 表示。
\textbf{设计动机与边界：} 相对原始 3DGS 以独立高斯和图像重建误差为主的设定，MagicDrive3D 把问题落在以下可实现的改动上：MagicDrive3D的核心流程分为两阶段：1）视频视角合成：训练一个可控多视角视频生成模型，以场景描述S\_t和相机位姿P\_c,t为条件，生成多视角连续图像序列。使用相对位姿控制实现灵活的相机运动。该模型在常规自动驾驶数据（如nuScenes）上训练，将动态场景转化为静态场景。2）3DGS生成：利用生成的多视角图像通过FTGS构建3DGS表示。FTGS包含三项改进：使用预训练单目深度模型提供更好的高斯初始化；引入外观建模处理不同视角间的曝光差异；设计容错损失函数处理生成图像中的微小几何误差。整体框架将控制信号的处理交给视频生成器（处理动态物体和语义），3DGS生成步骤保证任意视角渲染的泛化能力。支持BEV分割等下游任务增强。其设计目的不是单独增加一个模块，而是用上述表示、监督或渲染约束消除原流程在目标任务中的具体失配。提出MagicDrive3D框架，结合视频视角合成与3DGS生成实现可控街景3D生成。首先训练多视角视频生成模型合成多样街景视角，再通过Fault-Tolerant Gaussian Splatting（FTGS）处理生成视图中的微小误差。使用单目深度初始化和外观建模处理跨视角曝光差异。支持道路地图、3D物体、文本描述等多条件控制，可使用常规自动驾驶数据训练。适用时还应保留论文给出的边界：视频生成质量影响最终3DGS质量；对于极端天气或罕见场景生成效果有限；FTGS的容错能力有上限。
\textbf{方法信息：} Huawei 参与了这篇论文；引用很多，Semantic Scholar 约 74 次。



% METHOD: 2405.20323
\paragraph{S3Gaussian（4D 自监督街景高斯）（2024）}
$S^3$Gaussian~\cite{2405_20323} 是自监督街景高斯的另一条锚点路线。它继承 Street Gaussians 的显式高斯表示，但不再使用 3D 车辆框来分离动态对象。核心想法是：如果一个高斯只属于静态背景，它的解释主要由空间位置决定；如果它属于运动车辆，就必须依赖时间维度才能解释不同帧中的位置和外观。因此方法在标准 3D Gaussians 之外增加 spatial-temporal field network，用 4D 一致性来学习动静分解。

具体结构由 multi-resolution HexPlane encoder 和 multi-head Gaussian decoder 组成。HexPlane 把 4D 网格拆成三个空间平面和三个时空平面，在不同分辨率上聚合特征；高斯的位置和时间作为输入，经编码后由解码头预测规范空间位移、语义或动态相关属性。训练中仍然用可微 splatting 做图像重建，同时用位移和颜色正则让“没有运动”成为默认解释，避免所有高斯都无意义地变形；可渲染语义图则帮助观察和约束分解结果。公开代码中的训练脚本还显示了长序列处理方式：先做一段静态 warm-up，再以 50 帧左右的 clip 逐段训练，后一段可以加载前一段 checkpoint 作为 prior，以维持跨片段的时空连续性。相对 DeSiRe-GS 的“先找运动 mask 再蒸馏”，$S^3$Gaussian 更像端到端学习一个时空场，让动态区域自然落到时空平面上。

\textbf{实现细节：} $S^{}3$Gaussian（2024.05） 首次提出无需 3D 标注的自监督街景高斯方法，从 4D 一致性分解动静元素。Baseline 与问题。此前的动态街景方法（如 Street Gaussians、DrivingGaussian）依赖 3D 边界框标注来分离动态物体和静态背景，获取这些标注成本高昂。如何在完全没有物体边界框、分割掩码和光流标注的情况下实现高质量的动静分解是一个开放问题。核心改动——时空场网络。$S^{}3$Gaussian 维护一组规范（canonical）3D 高斯，配备时空场网络预测每帧的变形。网络核心是 Multi-resolution HexPlane 编码器，将 4D 时空网格分解为六个多分辨率可学习特征平面：三个空间平面 $P_{xy}, P_{xz}, P_{yz}$ 和三个时空平面 $P_{xt}, P_{yt}, P_{zt}$。这种分解天然有利于动静分离——动态物体在时空平面上呈现明显特征，静态物体仅在空间平面上可见。多解码头分别预测位移 $\Delta X$、SH 系数残差 $\Delta C$ 和语义特征 $f_s$，变形后的 4D 高斯为 $G' = \{X + \Delta X, C + \Delta C, s, r, o, f_s\}$。自监督优化。使用 LiDAR 点云替代 SfM 初始化以提供更好的几何先验。设计动机。HexPlane 分解的核心洞察是：静态物体在时间维度上不变，因此其特征完全由空间平面 $P_{xy}, P_{xz}, P_{yz}$ 编码，时空平面 $P_{xt}, P_{yt}, P_{zt}$ 的贡献趋近于零；动态物体则必须在时空平面上产生显著特征才能准确重建——这种天然的特征分离使得变形场自动区分动静区域。位移和颜色正则化进一步强化这一先验，最小化运动假设的复杂度，使"无需运动"成为默认选择。关键公式。
\textbf{锚点细化：} $S^{}3$Gaussian：针对此前动态街景方法（Street Gaussians、DrivingGaussian）依赖 3D 边界框标注来分离动静、获取标注成本高昂的局限，用\textbf{Multi-resolution HexPlane 时空场网络}替代边界框驱动的动静分离，核心思想是\textbf{将 4D 时空网格分解为三个空间平面（$P_{xy}, P_{xz}, P_{yz}$）和三个时空平面（$P_{xt}, P_{yt}, P_{zt}$），利用天然的特征分离自动区分动静区域}——静态物体仅在空间平面上可见、动态物体必须在时空平面上产生显著特征——配合位移和颜色正则化最小化运动假设的复杂度、使"无需运动"成为默认选择，首次实现无需 3D 标注的自监督街景高斯方法。
\textbf{方法信息：} 代码已公开：\url{https://github.com/nnanhuang/S3Gaussian}，仓库约 544 stars；Semantic Scholar 约 111 次引用。



% METHOD: 2406.18198
\paragraph{VDG（视觉里程计初始化的无姿态动态高斯）（2024.06）}
VDG~\cite{2406_18198} 去掉 SfM、LiDAR 和预计算姿态依赖。它把 self-supervised visual odometry 接入动态高斯：相邻图像估计相对位姿、单目深度和 motion mask，由深度和绝对姿态初始化单视图高斯。每个高斯的均值和不透明度改成围绕生命峰值的时间函数，实现静动态自监督分解；运动 mask 进一步监督分解。相对 SplatFlow 的 LiDAR 与相机融合，VDG 的关键是仅用 RGB 输入启动自动驾驶动态场景重建。
\textbf{实现细节：} DENSER（2024.09） 引入小波动态估计 SH 基增强远距离动态物体外观，跨多帧致密化点云加速收敛。VDG（2024.06） 首次集成自监督视觉里程计实现仅 RGB 输入的无位姿动态高斯。ArmGS（2025.07） 通过多粒度外观精化（局部高斯级、全局图像级、动态角色级）提升帧间和相机间精细变化建模。2D 先验与 SDF 引导（2025.10） 结合 SDF 与 3DGS 实现无 LiDAR 城市场景渲染和场景编辑。SRIF（2024.09） 利用扩散图像变形和流估计实现语义形状注册。
\textbf{方法信息：} Baidu 参与了这篇论文。
\textbf{设计动机与边界：} 相对原始 3DGS 以独立高斯和图像重建误差为主的设定，VDG 把问题落在以下可实现的改动上：VDG解决了动态高斯splats方法依赖SfM预计算姿态和初始化的问题：1）自监督VO集成：利用自监督视觉里程计从相邻图像估计相对姿态、运动掩码和单目深度图。自监督VO的优势是深度和姿态在同一尺度上对齐，这对高斯初始化至关重要。2）单视图高斯初始化：使用预测的深度图和绝对姿态从单目深度提升点来初始化3DGS，替代广泛使用的SfM点云。3）静动态分解与建模：借鉴PVG的思路，将传统高斯的均值和不透明度修改为围绕生命峰值的时间依赖函数。每个高斯的均值做周期性运动，不透明度呈高斯衰减，实现静动态的自监督分解。4）运动监督：引入运动监督机制，利用运动掩码更好地分解静动态高斯。5）大规模场景优化：提出训练策略在姿态优化过程中维持几何表示的稳定性，适用于自动驾驶等大规模前向运动场景。其设计目的不是单独增加一个模块，而是用上述表示、监督或渲染约束消除原流程在目标任务中的具体失配。首次将自监督视觉里程计（VO）集成到无姿态动态高斯方法中，实现姿态和深度的无监督初始化以及静动态分解。仅需RGB图像输入即可构建动态场景，无需SfM、LiDAR或额外传感器。支持更大规模场景和更快构建速度。适用时还应保留论文给出的边界：自监督VO的精度受限于训练数据。单目深度的尺度恢复可能不够准确。对快速运动或纹理缺失区域的处理可能受限。
% METHOD: 2407.02598
\paragraph{AutoSplat（道路天空约束与反射前景）（2024.07）}
AutoSplat~\cite{2407_02598} 继承 DrivingGaussian/Street Gaussians 的背景--前景分解，但把自动驾驶场景里最稳定的结构先验写进训练过程。普通 3DGS 在变道视角下容易把道路拉成漂浮纹理，天空也会因为无深度约束而不稳定；前景车辆又常只有可见面被 LiDAR 或 SfM 覆盖。AutoSplat 的改变是分别约束道路、天空和车辆：道路高斯被压成接近平面的各向异性椭球，天空由单独分支或远景表示处理，车辆则用 3D 模板初始化，并用反射一致性把可见侧的信息传给不可见侧。

流程上，方法先用语义分割得到道路、天空和其他背景区域，并分阶段优化背景高斯；道路/天空区域的尺度自由度被限制，避免它们在稀疏视角下向错误深度发散。前景车辆不完全依赖稀疏点云，而是从模板几何生成初始高斯；训练时把可见面的高斯沿车辆对称平面反射，形成对不可见面的约束。颜色和反射随时间变化的问题则交给 residual spherical harmonics，让车辆在不同时刻保留动态外观。相对 DrivingGaussian 的“先分对象再合成”，AutoSplat 的增量是把“道路应当近似平、天空应当远、车辆通常左右近似对称”这些场景知识转化为高斯参数约束，尤其面向车道横向移动和新轨迹模拟。

\textbf{实现细节：} AutoSplat（2024.07） 对道路和天空区域的高斯施加几何约束实现多视角一致的车道变换模拟。道路高斯被约束为扁平的各向异性椭球（垂直缩放接近零），天空高斯被推到无穷远并通过环境贴图渲染。利用 3D 模板引入反射高斯一致性约束监督前景物体可见和不可见面——可见面的高斯通过光度损失优化，不可见面的高斯通过模板的对称性约束间接正则化。并估计每个前景高斯的残差球谐以建模动态外观。
\textbf{设计动机与边界：} 相对原始 3DGS 以独立高斯和图像重建误差为主的设定，AutoSplat 把问题落在以下可实现的改动上：AutoSplat针对自动驾驶场景的三个核心问题提出解决方案：(1) 背景重建：将道路和天空区域从背景中分离，约束其高斯变为扁平形状（减少一个维度的缩放），确保多视角一致性。使用语义分割模型获取区域掩码，分两阶段训练背景；(2) 前景重建：利用3D模板作为先验初始化前景物体的高斯（因为SfM和LiDAR点云稀疏），引入反射高斯一致性约束：将可见面的高斯沿对称面反射来监督不可见面，使用分组的球谐系数确保颜色一致性；(3) 动态外观：为每个前景高斯估计残差球谐系数，建模随时间变化的外观变化。最终融合前景和背景高斯生成统一表示。其设计目的不是单独增加一个模块，而是用上述表示、监督或渲染约束消除原流程在目标任务中的具体失配。提出AutoSplat，针对自动驾驶场景的特殊需求设计约束：(1) 道路和天空区域高斯的几何约束确保多视角一致性（尤其变道场景）；(2) 基于对称平面的反射高斯一致性约束监督前景物体不可见面；(3) 残差球谐函数建模前景动态外观。在Pandaset和KITTI上超越SOTA。适用时还应保留论文给出的边界：依赖3D模板和语义分割的准确性；对称假设可能不适用于所有物体类型；大规模场景的内存消耗较高
\textbf{方法信息：} Noah's Ark Lab, Huawei Technologies 参与了这篇论文；项目页为 \url{https://autosplat.github.io/}，当前未确认完整公开代码；Semantic Scholar 约 55 次引用。



% METHOD: 2407.02945
\paragraph{VEGS（外推视角街景合成）（2024）}

VEGS~\cite{2407_02945} 较早明确提出驾驶城市场景的 extrapolated view synthesis 问题。车载前向相机只覆盖窄轨迹，左看、右看或向下看时普通 3DGS 会退化。VEGS 的改变有三步：多帧 LiDAR 拼成密集静态地图和实例动态点云，给高斯更完整初始化；用法线先验监督协方差渲染，避免高斯形状过拟合前向视角；再微调扩散模型并蒸馏外推视角。它提供代码，GitHub 关注度也较高。相对 DrivingForward 的前馈输入灵活性，VEGS 专门补训练轨迹外视角。
\textbf{实现细节：} VEGS（2024.07） 首次系统研究城市场景中的外推视角合成（EVS）问题，揭示了前向相机优化导致高斯协方差"懒惰优化"的关键瓶颈。Baseline 与问题。标准 3DGS 在训练轨迹附近渲染质量良好，但偏离到新车道或方向时质量急剧下降：高斯基元的协方差和密度在仅有前向相机监督时会过拟合到训练视角分布，从侧面观察时暴露几何缺陷。此前缺乏系统研究这一外推视角合成（EVS）问题的工作。核心发现——懒惰协方差优化。VEGS 发现标准 3DGS 在仅有前向相机时，高斯协方差的形状和方向会过拟合到特定视角——协方差被训练为以最小优化代价覆盖训练像素的视锥。当从非训练方向观察时，这些过拟合的协方差在表面下方产生空洞（cavity），导致外推视角渲染质量急剧下降。三项技术贡献。
\textbf{设计动机与边界：} 相对原始 3DGS 以独立高斯和图像重建误差为主的设定，VEGS 把问题落在以下可实现的改动上：VEGS针对城市驾驶场景前向相机训练分布下的外推视角合成问题提出三个核心策略：(1) 密集LiDAR地图初始化：将多帧LiDAR点云拼接为密集静态地图和实例级动态点云，用于初始化高斯位置，比稀疏LiDAR监督更有效；(2) 法线先验约束协方差：发现前向相机优化导致高斯协方差过拟合到特定视角，利用表面法线估计器提取训练图像的法线图，通过新设计的协方差渲染器和法线监督约束高斯的形状和方向；(3) 扩散模型蒸馏外推视角：微调大规模扩散模型的子集参数（平衡场景特定知识和泛化能力），直接监督外推视角的渲染结果。动态场景建模分离静态模型和多个实例级动态模型。其设计目的不是单独增加一个模块，而是用上述表示、监督或渲染约束消除原流程在目标任务中的具体失配。首次提出城市场景中的外推视角合成(EVS)问题，并设计VEGS框架通过密集LiDAR初始化、法线先验约束高斯协方差和扩散模型蒸馏三种策略提升外推视角渲染质量。适用时还应保留论文给出的边界：依赖LiDAR传感器；扩散模型微调增加训练时间；外推角度过大时仍有质量下降
\textbf{方法信息：} 代码：\url{https://github.com/deepshwang/vegs}；GitHub 关注度较高，约 128 stars。



% METHOD: 2408.16760
\paragraph{OmniRe（全要素城市 Gaussian 场景图）（2024）}

OmniRe~\cite{2408_16760} 的目标是全要素城市数字孪生与交互仿真。它构建 Gaussian scene graph，节点包括天空、静态背景、刚性车辆和非刚性行人或骑行者。近距行人用 SMPL 节点绑定高斯并通过 LBS 驱动，远距行人用可变形节点；车辆在局部坐标系中定义高斯并随姿态变换。所有节点拼接后渲染。相对 Street Gaussians 主要建车辆前景，OmniRe 把车辆、行人、骑行者和天空都纳入统一场景图，适合交互仿真。
\textbf{实现细节：} OmniRe（2024.08） 扩展动态建模从车辆到行人、自行车等所有动态前景，在 3DGS 上构建场景图并在规范空间构建多种高斯表示。Baseline 与问题。此前动态街景重建方法（如 Street Gaussians）仅处理车辆这一类动态物体，采用刚体运动模型。但城市场景中行人、自行车等非刚性动态参与者的运动模式与车辆截然不同——行人有体态变化（摆臂、迈步），自行车有倾斜转向——刚体模型无法捕捉这些非刚性变形。如何在统一框架中同时建模刚性和非刚性动态物体，是实现完整街景重建的关键挑战。改了什么。OmniRe 将 Street Gaussians 的"单一刚体动态模型"扩展为"分类动态表示"——根据动态参与者的运动特性将其分为三类：静态背景（不变的高斯基元）、刚性物体（车辆，用刚体 SE(3) 变换建模）和非刚性物体（行人、自行车，在规范空间额外配备变形场）。实现原理。场景图为每种动态参与者分配独立的高斯簇和可优化的 SE(3) 位姿序列。刚性物体的建模沿用 Street Gaussians 的局部坐标系方法：车辆高斯定义在物体坐标系中，通过可优化位姿 $(R_t, T_t)$ 变换到世界坐标。
\textbf{实现补充：} 实现上还要核对输入、输出、参数、特征、初始化、训练损失、正则、渲染或推理接口，以及优化中的剪枝、分裂、采样、编码或解码步骤；这些抓手决定该方法到底改了 3DGS 管线的哪一环。
\textbf{设计动机与边界：} 相对原始 3DGS 以独立高斯和图像重建误差为主的设定，OmniRe 把问题落在以下可实现的改动上：OmniRe构建Gaussian场景图，包含四类节点：(1)天空节点：可优化环境纹理图；(2)背景节点：静态Gaussian表示建筑、道路等；(3)刚性节点：车辆等刚体，Gaussian在物体局部空间定义，通过姿态变换到世界空间；(4)非刚性节点：行人、骑行者等，近距行人用SMPL节点建模（通过SMPL参数驱动Gaussian变形），远距行人用可变形节点建模。Gaussian SMPL节点将SMPL网格顶点与Gaussian绑定，通过LBS权重实现姿态驱动的变形。所有节点Gaussian拼接后通过3DGS光栅化器渲染。支持人车交互等高级仿真场景，可利用GPT-4o对个体进行自动分类。其设计目的不是单独增加一个模块，而是用上述表示、监督或渲染约束消除原流程在目标任务中的具体失配。提出OmniRe，首个全方位城市场景重建系统，通过场景图上的多种Gaussian表示（背景、刚性节点、非刚性节点、天空）实现所有动态参与者的完整重建，支持约60Hz的交互式仿真。适用时还应保留论文给出的边界：依赖边界框标注；行人建模质量受SMPL拟合精度影响；数据集特定标签可能限制泛化。
\textbf{方法信息：} NVIDIA 参与了这篇论文；引用很多，Semantic Scholar 约 166 次。



% METHOD: 2409.10041
\paragraph{DENSER（小波动态外观与形状密集化）（2024.09）}
DENSER~\cite{2409_10041} 继承场景图动态街景表示，但针对远距离动态物体细节和外观建模不足进行改造。动态物体外观不再使用静态 SH 系数，而是通过小波变换动态估计 SH basis，使外观在空间和时间上更灵活。形状方面，多帧点云 densification 增强动态物体几何，加速收敛。相对 ArmGS 的多粒度外观网络，DENSER 更明确地用小波替代固定球谐建模动态外观。
\textbf{实现细节：} DENSER（2024.09） 引入小波动态估计 SH 基增强远距离动态物体外观，跨多帧致密化点云加速收敛。VDG（2024.06） 首次集成自监督视觉里程计实现仅 RGB 输入的无位姿动态高斯。ArmGS（2025.07） 通过多粒度外观精化（局部高斯级、全局图像级、动态角色级）提升帧间和相机间精细变化建模。2D 先验与 SDF 引导（2025.10） 结合 SDF 与 3DGS 实现无 LiDAR 城市场景渲染和场景编辑。SRIF（2024.09） 利用扩散图像变形和流估计实现语义形状注册。
\textbf{实现补充：} 实现上还要核对输入、输出、参数、特征、初始化、训练损失、正则、渲染或推理接口，以及优化中的剪枝、分裂、采样、编码或解码步骤；这些抓手决定该方法到底改了 3DGS 管线的哪一环。
\textbf{方法信息：} 代码：\url{https://github.com/sntubix/denser}。
\textbf{设计动机与边界：} 相对原始 3DGS 以独立高斯和图像重建误差为主的设定，DENSER 把问题落在以下可实现的改动上：DENSER框架基于场景图表示，将场景分解为静态背景节点和动态物体节点：(1)动态外观建模：不直接使用SH系数表示动态物体外观，而是引入小波变换动态估计SH基函数，使外观在空间和时间上更灵活地变化；(2)物体形状密集化：通过多帧场景中的点云密集化增强动态物体的形状表示，加速模型训练收敛；(3)场景图表示：背景节点在世界坐标系中直接优化，物体节点在各自参考系中优化后通过轨迹变换到世界坐标系进行渲染。所有高斯拼接后按标准3DGS方式渲染。其设计目的不是单独增加一个模块，而是用上述表示、监督或渲染约束消除原流程在目标任务中的具体失配。DENSER首次将小波变换引入3DGS动态物体外观建模，通过动态估计SH基函数替代静态SH系数，在KITTI上大幅超越SOTA方法。适用时还应保留论文给出的边界：场景图依赖物体检测和跟踪的准确性；对极端天气条件（雨雪雾）的建模有限；小波变换增加了计算复杂度。
% METHOD: 2409.12753
\paragraph{DrivingForward（环视输入的前馈高斯重建）（2024.09）}
DrivingForward~\cite{2409_12753} 把前面这些 per-scene 优化路线改成前馈环视重建。DrivingGaussian、Street Gaussians 和 AutoSplat 都要针对每段驾驶日志优化一个高斯场；如果目标是快速从任意环视输入生成场景，这种逐场训练就成了瓶颈。DrivingForward 的改变是联合训练 pose network、depth network 和 Gaussian network，让模型在推理时从灵活数量的 surround-view frames 直接预测高斯基元。

它不依赖大重叠多视角匹配，因为车载环视相机之间往往只有很小重叠区域。方法先用位姿网络自监督估计车辆运动，再用深度网络从时间上下文、空间相邻相机和时空联合约束中学习尺度感知深度；高斯网络则对每张输入图像独立预测协方差、不透明度和球谐系数。每个像素生成一个候选高斯，位置由预测深度反投影得到；训练时把这些来自多相机、多帧的高斯合成到 3D 空间，再通过可微渲染监督目标视角。实现中显式区分 single-frame 和 multi-frame 两种 novel-view 模式，训练入口只切换配置即可改变输入帧数；推理时主要保留 depth network 和 Gaussian network，因此可以在六路相机或不同时间间隔输入下直接前馈生成场景。相对 pixelSplat/MVSplat 这类依赖图像对几何的前馈方法，它把尺度学习放到驾驶时序和环视结构中；相对 DrivingGaussian，它牺牲了一部分逐场精修能力，换来无需长期优化的快速重建。

\textbf{锚点细化：} DrivingForward：针对前馈 3DGS 方法依赖大重叠配对视角通过极线几何预测高斯位置、但驾驶场景车载相机间重叠极小极线匹配几乎不可行的局限，用\textbf{自监督深度驱动的单帧高斯预测}替代多视角匹配范式，核心思想是\textbf{从单帧图像预测深度并通过三类光度损失学习尺度信息}——时间上下文（同一相机相邻帧）、空间上下文（同帧相邻相机重叠区域）和时空上下文联合约束，利用固定相机到车体变换矩阵而非世界坐标外参、实现无需深度真值和位姿标注的前馈重建。
\textbf{方法信息：} 代码已公开：\url{https://github.com/fangzhou2000/DrivingForward}，仓库约 232 stars；Semantic Scholar 约 55 次引用。



% METHOD: 2411.11921
\paragraph{DeSiRe-GS（自监督动静分解与表面重建）（2024）}
DeSiRe-GS~\cite{2411_11921} 是自监督动态街景的重要锚点。DrivingGaussian 和 Street Gaussians 用 3D 框把动态对象分出来，效果清楚但标注或检测链路很重；$S^3$Gaussian 则把所有高斯交给统一时空场，容易把静态区域也误当成可运动。DeSiRe-GS 的改变是利用一个反直觉但有用的现象：不含时间建模的 3DGS 往往能重建静态区域，却会在动态车辆处留下残差和鬼影。因此，静态重建本身可以作为运动检测器。

它采用两阶段优化。第一阶段先训练静态 3DGS，再比较渲染图和真实图像的预训练特征差异，由 MLP 解码出 2D motion mask；动态区域不是靠 3D 框给出，而是由“静态模型解释不了哪里”得到。更细地说，特征差异用余弦距离形成动态性提示，掩码损失约束 MLP 在高差异区域输出动态响应，masked rendering 又反过来让静态模型更专注于可解释区域。第二阶段把这些 2D motion priors 以可微形式蒸馏进 Gaussian 空间，约束 PVG 式速度或动态属性，减少静态区域的非零速度噪声。为了避免自动驾驶稀疏视角下的高斯漂浮，它还加入受 2DGS 启发的法线一致性、尺度等几何正则，并用 temporal cross-view consistency 让不同时间和视角看到的表面更一致。相对 $S^3$Gaussian，它不是只让 HexPlane 自己学分解，而是先用静态残差定位运动，再把运动信息送回 3D 高斯优化。

\textbf{实现细节：} DeSiRe-GS（2024.11） 修改了 3DGS 的动态场景处理环节——用两阶段自监督分解管线替代 $S^{}3$Gaussian 的统一变形场方案——静态重建能力本身被用作运动检测器。核心观察极为简洁：原始 3DGS 不包含时间信息，因此能准确重建静态区域（停放的车辆、建筑物），但在动态区域会产生伪影（漂浮的"鬼影"点），导致渲染图像与真值之间出现显著差异。第一阶段：2D 运动掩码提取。问题：如何不依赖任何标注来识别动态区域？方案：直接利用 3DGS 的重建失败作为动态性信号。具体而言，用预训练基础模型分别提取渲染图像 $\hat{I}$ 和真值图像 $I$ 的特征 $\hat{F}$ 和 $F$，计算逐像素不相似度： 不相似度 $D$ 在静态区域趋近 0（重建成功），在动态区域趋近 1（重建失败），且该计算不涉及任何可学习参数。但 DeSiRe-GS 并未对 $D$ 直接阈值化，而是引入一个 MLP 解码器预测动态概率 $\hat{\mathcal{D}} \in \mathbb{R}^{H \times W}$，以特征中的丰富语义信息为输入、以 $D$ 为监督信号： 即鼓励解码器在 $D$ 高的区域输出高值。最终二值掩码 $M = \mathbb{I}(\hat{\mathcal{D}} > \tau)$（$\tau$ 为固定阈值）。关键设计：联合优化图像渲染与掩码预测——用掩码遮蔽动态区域后的渲染损失 $L_{\mathrm{masked-render}} = M \odot \|\hat{I} - I\|$ 反过来让动态区域的差异更显著，形成互惠增强。第一阶段训练共 30,000 次迭代，运动解码器从第 6,000 次迭代开始训练。第二阶段：掩码蒸馏至高斯空间。问题：第一阶段的掩码仅存在于 2D 图像空间且依赖真值图像，无法用于新视角合成。方案：以可微分方式将 2D 运动信息蒸馏到 3D 高斯空间。DeSiRe-GS 采用 PVG（Periodic Vibration Gaussian）这一时间变化高斯表示——每个高斯附加速度 $v$，渲染时生成速度图 $V \in \mathbb{R}^{H \times W}$。
\textbf{锚点细化：} DeSiRe-GS：针对 $S^{}3$Gaussian 将所有高斯视为潜在动态元素通过 HexPlane 统一建模、运动信息学习从间接损失正则化到直接掩码蒸馏的局限，用\textbf{渲染残差驱动的两阶段自监督分解管线}替代统一变形场方案，核心思想是\textbf{将静态重建能力本身用作运动检测器}——原始 3DGS 不含时间信息能准确重建静态区域但在动态区域产生"鬼影"伪影，第一阶段用渲染图像与真值特征的逐像素不相似度提取 2D 运动掩码，第二阶段以可微分方式将 2D 运动信息蒸馏到 3D 高斯空间（用掩码正则化 PVG 速度图消除静态区域的非零速度噪声），在 Waymo 和 KITTI 上实现超越 PVG 和 $S^{}3$Gaussian 的自监督分解质量。
\textbf{方法信息：} 代码已公开：\url{https://github.com/chengweialan/DeSiRe-GS}，仓库约 241 stars；Semantic Scholar 约 54 次引用。



% METHOD: 2411.15355
\paragraph{UniGaussian（鱼眼与针孔统一高斯）（2024）}
UniGaussian~\cite{2411_15355} 解决自动驾驶中鱼眼相机无法直接适配标准 3DGS 的问题。它分析 Kannala-Brandt 和 MEI 鱼眼模型，发现光线扭曲可转化为相机-物体向量旋转后的辐射场变形，于是用仿射变换让标准 3DGS 光栅化流程渲染鱼眼图像，避免先矫正图像造成画质损失。统一框架同时监督针孔和鱼眼相机，并融合深度、语义、法线和 LiDAR。相对 ParkGaussian 的停车鱼眼应用，UniGaussian 提供的是相机模型层面的可微渲染基础。
\textbf{实现细节：} UniGaussian（2024.11） 与 ParkGaussian 针对鱼眼相机使用无迹变换（Unscented Transform）不同，UniGaussian 提出基于仿射变换序列的可微分鱼眼渲染方法，统一建模多种传感器（针孔/鱼眼相机）和模态（深度、语义、法向、LiDAR 点云）。鱼眼相机的高斯变换。核心观察是鱼眼相机模型（Kannala-Brandt 和 MEI）的数学本质是将相机-点向量向光轴旋转后再投影——这与 3DGS 的渲染流程结构相似，只需将 3D 点替换为 3D 高斯并引入变形。具体实现分三步：（1）位置和姿态调整——将高斯中心沿旋转轴 $\mathbf{r}_{rot}$ 旋转角度 $\Delta\theta = \theta_d - \theta$（$\theta_d$ 由鱼眼模型计算），四元数同步更新 $q_g' = q \otimes q_g$；（2）极向和切向压缩——由于鱼眼压缩视场角，高斯需要相应压缩以避免泼溅平面上过度重叠，极向压缩比 $k_\theta = d\theta_d / d\theta$ 通过泰勒展开从鱼眼模型导数计算，切向压缩比 $k_\phi = \theta_d / \theta$；（3）缩放因子和旋转更新——压缩后的协方差 $\Sigma' = S_\theta \Sigma S_\phi^T S_\phi \Sigma S_\theta^T$ 通过特征值分解提取新的缩放因子和四元数。整个过程完全可微分，且不需额外 GPU 显存，适合处理驾驶场景重建中的百万级高斯。多模态统一渲染。
\textbf{方法信息：} Huawei 参与了这篇论文。
\textbf{设计动机与边界：} 相对原始 3DGS 以独立高斯和图像重建误差为主的设定，UniGaussian 把问题落在以下可实现的改动上：方法分为两步：(1) 鱼眼3DGS渲染：分析Kannala-Brandt和MEI两种鱼眼模型的数学原理，发现其本质都是通过旋转相机-物体向量将3D点移到新位置再投影。据此设计仿射变换方法：将鱼眼光线扭曲转化为3D高斯体的辐射场变形，使标准3DGS光栅化流程适用于鱼眼图像渲染。相比直接矫正鱼眼图像或拼接立方体图的方法，避免了画质损失和拼接伪影。(2) 统一框架：在可微分渲染基础上构建多相机多模态框架，不同相机模型通过相应仿射变换适配；共享高斯体通过来自不同模态（深度、语义、法线、LiDAR）的监督进行正则化，学习统一的3D高斯表征。支持近场感知和自动泊车等鱼眼相机关键应用场景。其设计目的不是单独增加一个模块，而是用上述表示、监督或渲染约束消除原流程在目标任务中的具体失配。提出UniGaussian，首个支持鱼眼相机的3DGS自动驾驶场景重建方法。核心创新：(1) 新的3DGS鱼眼相机可微分渲染方法，将光线扭曲转化为辐射场变形，通过仿射变换实现，保持实时渲染和可微分性；(2) 统一多相机多模态框架，同时建模针孔和鱼眼相机，融合深度、语义、法线和LiDAR点云等多种模态。实现高质量驾驶场景模拟。适用时还应保留论文给出的边界：鱼眼相机模型的数学近似可能对极端畸变不够精确；多模态融合增加训练复杂度；依赖准确的相机标定参数。
% METHOD: 2411.15482
\paragraph{SplatFlow（神经运动流场）（2024）}
SplatFlow~\cite{2411_15482} 也是无 3D 框动态街景路线，但它不把时间建模主要交给 HexPlane，而是显式学习 Neural Motion Flow Field（NMFF）。问题在于，渲染损失能告诉模型“这一帧像不像”，却不一定能稳定约束动态物体在相邻时间之间如何对应。SplatFlow 的改变是先用时序 LiDAR 点学习 3D motion flow prior，再把这个连续流场接入 4D Gaussians，让每个动态高斯跨时间传播时有可学习的运动对应关系。

方法先用 NMFF 处理连续 LiDAR 点，生成动态/静态分解和场景流先验；静态点初始化为 3D Gaussians，动态部分初始化为与时间相关的 4D Gaussians。随后在图像优化阶段，NMFF 对动态高斯的中心和状态做逐步传播，建立相邻时间之间的对应；这些对应被用于聚合时间特征，使动态车辆在不同视角和时刻下更一致。为了让运动分解不只依赖点云，方法还把 2D 光流和基础模型特征蒸馏到 4D 时空表示中，并联合优化图像重建、光流/特征蒸馏和正则项。相对 $S^3$Gaussian 的时空平面编码，SplatFlow 的技术重点是“先学运动流，再让高斯跟着运动流走”。

\textbf{实现细节：} SplatFlow（2024.11） 设计神经运动流场（Neural Motion Flow Field, NMFF）学习无追踪 3D 边界框的 4D 时空表示。关键设计在于动静分解的几何驱动方式：对每个 LiDAR 点计算 NMFF 预测的运动流大小，通过阈值化自然分离静态和动态点，分别合并后初始化静态 3D 高斯和动态 4D 高斯。动态 4D 高斯在中间时间戳的参考帧中聚合所有观测增加密度和完整性，推理时通过 NMFF 的逐步传播将聚合高斯变形到目标时刻。同时将 SEA-RAFT 光流知识蒸馏到 4D 表示中——利用 NMFF 的 3D 运动流渲染 2D 光流作为监督信号，实现无需人工标注的端到端自监督动态场景重建。
\textbf{方法信息：} Purdue University 与 Microsoft 参与了这篇论文；当前未找到与 arXiv:2411.15482 对应的完整公开代码。
\textbf{设计动机与边界：} 相对原始 3DGS 以独立高斯和图像重建误差为主的设定，SplatFlow 把问题落在以下可实现的改动上：方法包含三个核心设计：(1) 神经运动流场（NMFF）：一组隐函数，将LiDAR点和3DGS的时序运动建模为连续运动流场。首先用NMFF将时序LiDAR点分解为静态和动态点集，分别初始化静态3D高斯和动态4D高斯。NMFF学习每个4D高斯跨时间的对应关系，在优化4D高斯属性的同时建模时序关联。(2) 4D高斯表征：每个4D高斯基元具有时变属性（3D中心位置mu(t)和协方差Sigma(t)）和时不变属性（不透明度和颜色）。动态物体用与NMFF关联的聚合4D高斯表征，静态背景用3D高斯表征。(3) 2D特征蒸馏：从2D基础模型蒸馏光流知识到4D时空表征，增强动态物体识别能力。整体框架从周围相机和LiDAR的带时间戳数据中学习。其设计目的不是单独增加一个模块，而是用上述表示、监督或渲染约束消除原流程在目标任务中的具体失配。提出SplatFlow，首个完全自监督的动态3DGS自动驾驶场景重建方法。核心创新：(1) 神经运动流场（NMFF）将LiDAR点和3DGS的时序运动建模为连续运动流场，自监督分解动静元素；(2) NMFF建模每个4D高斯跨时间的对应关系，聚合时序特征增强动态组件的跨视角一致性；(3) 从2D基础模型蒸馏特征到4D时空表征增强动态物体识别。无需3D边界框标注。适用时还应保留论文给出的边界：依赖LiDAR点云输入；NMFF的运动建模精度受时序分辨率影响；对快速运动物体的处理仍有提升空间。
% METHOD: 2411.15582
\paragraph{EMD（显式运动嵌入和双尺度变形）（2024.11）}
EMD~\cite{2411_15582} 是可插拔 motion module。它为每个高斯引入可学习 motion embedding，并把空间特征、时间特征和高斯特有 latent 组合为运动感知编码。双尺度变形框架分别处理车辆这类大幅快速运动和行人这类微幅慢速运动；自适应时间嵌入从粗时间采样逐步细化到更密样本。相对 S3Gaussian 或 DeSiRe-GS 的整体方法，EMD 更像可接入多种街景高斯系统的运动增强层。
\textbf{实现细节：} EMD（2024.11） 通过引入可学习运动嵌入到高斯基元中增强动静分解，作为即插即用模块补偿自监督方法中缺失的运动建模。与 $S^{}3$Gaussian 用全局 HexPlane 统一建模所有运动不同，EMD 的核心改动是为每个高斯分配可学习的个体嵌入 $z_k \in \mathbb{R}^{32}$，并通过双尺度变形框架分别建模大尺度运动（车辆平移，$D_{coarse}$）和精细变形（行人体态变化，$D_{fine}$）。设计动机是街景中车辆和行人的运动模式差异显著——车辆以刚体平移为主、速度较快，行人则是非刚体运动、速度较慢——全局变形场无法有效捕捉这种异质性。
\textbf{实现补充：} 实现上还要核对输入、输出、参数、特征、初始化、训练损失、正则、渲染或推理接口，以及优化中的剪枝、分裂、采样、编码或解码步骤；这些抓手决定该方法到底改了 3DGS 管线的哪一环。
\textbf{方法信息：} Intel 参与了这篇论文。
\textbf{设计动机与边界：} 相对原始 3DGS 以独立高斯和图像重建误差为主的设定，EMD 把问题落在以下可实现的改动上：EMD是一个即插即用模块，核心包含两个组件：(1) 运动感知特征编码：将高斯体参数编码到综合特征空间，结合三个维度——空间特征（多频位置编码）、时间特征（自适应时间嵌入矩阵，从30到150个时间样本渐进细化）、高斯特有特征（每个高斯体的可学习潜嵌入z\_k）。(2) 双尺度变形框架：针对街道场景中不同运动模式设计，大尺度分支处理车辆等大幅快速运动，小尺度分支处理行人等微幅慢速运动。时间嵌入矩阵W在训练过程中通过渐进采样细化，初始使用少量时间样本捕获粗粒度运动，逐步增加样本量捕获更精细的时序动态。可无缝集成到S3Gaussian、DeSiRe-GS等自监督方法，也可通过定制训练策略扩展到StreetGaussian等监督方法。其设计目的不是单独增加一个模块，而是用上述表示、监督或渲染约束消除原流程在目标任务中的具体失配。提出EMD（Explicit Motion Decomposition），一个即插即用模块，通过为高斯体引入可学习运动嵌入来增强街道场景的动静分解。核心创新：(1) 运动感知特征编码，结合空间、时间和高斯特有信息；(2) 双尺度变形框架，分别处理大幅运动（车辆）和微幅运动（行人）；(3) 自适应时间嵌入的渐进采样细化策略。在自监督设置下达到SOTA新视角合成性能，也可扩展到监督方法。适用时还应保留论文给出的边界：运动嵌入的维度需要调优；双尺度分支的选择基于先验知识（大/小运动）；对非常规运动模式（如自行车）的分类可能不够精细。
% METHOD: 2411.19548
\paragraph{ReconDreamer（在线修复扩展新轨迹）（2024）}
ReconDreamer~\cite{2411_19548} 解决多车道偏移等大幅机动轨迹渲染失败。它用 DriveRestorer 修复从训练不足模型渲染出的退化帧，并用 PDUS 渐进式数据更新策略逐步扩大新颖轨迹范围。每次扩展后，DriveRestorer 修复伪影，修复结果再进入训练集优化 3DGS。结构条件如 3D boxes 和 HD maps 通过投影与新轨迹对齐，保证时空一致。相对 FaithFusion 的像素信息增益，ReconDreamer 更像 curriculum：从原轨迹逐步走向更远轨迹。
\textbf{实现细节：} 多个工作从不同角度丰富了街景重建工具箱。MagicRoad（2025.07） 通过语义引导颜色增强和障碍物修复实现鲁棒道路表面重建。GS-RoadPatching（2025.09） 通过 3D 空间结构匹配实现驾驶场景的替代式修复。MOGS（2025.09） 用单目深度替代 LiDAR 降低传感器依赖。LT-Gaussian（2025.08） 提出高斯地图长期更新方法。LS-Gaussian（2025.07） 实现轻量级流式渲染的算法-硬件协同设计。CE-NPBG（2025.04） 通过连接关系图改善大规模驾驶场景中点云渲染的可扩展性。MuDG（2025.03） 将多模态扩散模型与高斯泼溅结合增强极端视角变化下的渲染鲁棒性。Zheng et al.（2025.03） 提出自适应透明度机制消除运动物体。ReconDreamer（2024.11） 通过在线修复和渐进数据更新实现大偏移渲染。
\textbf{方法信息：} 引用很多，Semantic Scholar 约 89 次。
\textbf{设计动机与边界：} 相对原始 3DGS 以独立高斯和图像重建误差为主的设定，ReconDreamer 把问题落在以下可实现的改动上：ReconDreamer框架包含三个核心组件：(1)DriveRestorer——基于世界模型微调的修复网络，通过利用训练不足的重建模型渲染原始轨迹视频来构建修复训练数据对（退化帧-真实帧），并使用masking策略强调天空和远距离区域的修复。(2)PDUS渐进式数据更新策略——从原始轨迹训练开始，逐步扩展新颖轨迹范围，每次扩展后用DriveRestorer修复渲染伪影，更新训练数据集，迭代优化直至收敛。(3)结构条件（3D框和HDMap）通过投影变换对齐新颖轨迹，确保空时一致性。其设计目的不是单独增加一个模块，而是用上述表示、监督或渲染约束消除原流程在目标任务中的具体失配。提出首个能有效渲染大幅机动（如多车道偏移）的方法，通过DriveRestorer在线修复和PDUS渐进式数据更新策略，在NTA-IoU、NTL-IoU和FID上分别相对提升24.87%、6.72%和29.97%超越Street Gaussians。适用时还应保留论文给出的边界：依赖世界模型的生成质量；DriveRestorer的训练需要构建修复数据对；PDUS的迭代过程增加了总训练时间。
% METHOD: 2401.01339
\paragraph{Street Gaussians（场景图式前景车辆建模）（2024.01）}
Street Gaussians~\cite{2401_01339} 是动态街景显式高斯表示的基础锚点。它面对的问题是：NeRF 式街景方法可以借助车辆轨迹建动态对象，但训练和渲染都慢；原始 3DGS 又把所有内容塞进一个静态点集，无法干净地区分背景和前景车辆。Street Gaussians 的改变是用一个显式场景图组织高斯：背景是世界坐标系中的点云和高斯，前景每辆车都有自己的局部点云、高斯和可优化跟踪位姿，渲染时再按位姿变换回世界坐标。

具体来说，背景点云中的每个点携带 3D Gaussian 参数和语义 logits，用于道路、建筑、天空等静态内容。前景车辆在对象局部坐标中建模，位姿来自 tracker 但在训练中继续优化，避免检测/跟踪误差直接变成重建误差。车辆外观不是固定的 SH 系数，而是用 4D spherical harmonics 表示随时间变化的反射和光照；这样同一辆车在不同帧、不同视角下可以有不同外观。训练时将背景高斯和各车辆高斯合并后进行标准可微 splatting，以图像重建损失优化；实现配置中还把 sky mask、object mask、LiDAR depth、DSSIM 和正则项接入同一优化过程，并为背景与对象设置不同的 densification/pruning 阈值，避免大背景和小车辆被同一密度规则处理。初始化通常结合 LiDAR 点云与 SfM/Colmap 点云，车辆点再按跟踪框分配到对象局部空间；渲染或编辑时，只要改变某个对象的局部到世界位姿，就能移除、平移、重放或重组车辆而不必重训背景。相对 DrivingGaussian 的增量是把“对象局部高斯 + 可优化轨迹 + 时变外观 + 分组件密度控制”做成更清晰的街景表示，因此成为后续自监督分解和可编辑仿真的常用基线。

\textbf{锚点细化：} Street Gaussians：针对标准 3DGS 假设场景完全静态、无法处理城市街道中移动车辆的局限，用\textbf{动静分解的显式高斯点云}替代无结构高斯集合，核心思想是\textbf{将背景用世界坐标高斯建模、将前景车辆用物体局部坐标高斯建模再通过可优化位姿变换到世界坐标}——前景高斯定义在物体坐标系中、通过追踪位姿 $(R_t, T_t)$ 变换到世界坐标，4D 球谐函数模型用傅里叶变换系数替代每帧独立的 SH 系数以 $O(k)$ 参数编码时间信息（直觉上如同音频信号用几个主频分量重建，任何周期性变化的外观如阴影移动、反光闪烁都能用少数几个余弦波叠加近似），使前景与背景可自由组合支持场景编辑、以 135 FPS 实现渲染。
\textbf{实现细节：} Street Gaussians（2024.01） 是该方向的里程碑工作，奠定了动态街景显式高斯表示的基本范式。Baseline 与问题。在 Street Gaussians 之前，动态街景重建主要依赖 NeRF 类隐式方法（如 UniSim、Mars），这些方法将场景建模为连续神经辐射场，训练和渲染速度慢，且难以对前景车辆进行单独编辑。3DGS 虽然在静态场景中实现了实时渲染，但无法直接处理包含移动车辆的动态街景。核心改动。Street Gaussians 将动态城市街景表示为一组带语义 logits 的 3D 高斯点云，分为静态背景和 $N$ 个前景车辆两组独立建模。背景模型使用世界坐标系下的标准 3D 高斯（位置 $\boldsymbol{\mu}_b$、协方差 $\Sigma_b$、不透明度 $o_b$、球谐系数 $z_b$）。前景车辆的每个高斯定义在物体局部坐标系中，通过可优化的追踪位姿 $(R_t, T_t)$ 变换到世界坐标系： \begin{equation*} \boldsymbol{\mu}_w = R_t \boldsymbol{\mu}_o + T_t, \quad R_w = R_t R_o, \quad \Sigma_w = R_w S_o S_o^T R_w^T \end{equation*} 其中下标 $w$ 表示世界坐标，$o$ 表示物体局部坐标。关键创新在于 4D 球谐函数模型：将每个 SH 系数 $z_{m,l}$ 替换为一组傅里叶变换系数 $\{f_i \in \mathbb{R}^k\}$，通过逆离散傅里叶变换恢复时间 $t$ 的外观：$z_{m,l}(t) = \sum_{i=0}^{k-1} f_i \cos(2\pi i t / T)$。直觉上，任何随时间周期性变化的外观（如阴影移动、反光闪烁）都可以用少数几个不同频率的余弦波叠加近似——正如音频信号可以用几个主频分量重建。这样以极少的额外参数编码了时间信息，避免了为每帧单独存储 SH 系数的高存储开销。渲染与初始化。渲染时，将背景和变换后的前景点云拼接后统一投影到 2D 图像空间进行 $\alpha$-混合。天空区域用高分辨率立方体贴图单独建模。初始化方面，用聚合 LiDAR 点云替代 SfM 点云以获得更好的几何先验，物体模型用 3D 边界框内的 LiDAR 点或随机采样点初始化。设计动机。显式点云表示使得前景与背景可自由组合、支持场景编辑操作；可学习的位姿残差补偿追踪器的噪声；4D 傅里叶 SH 以 $O(k)$ 参数（$k$ 为傅里叶阶数，通常取 3--4）替代 $O(T \times C_{SH})$ 参数，兼顾时间建模能力和存储效率。性能。
\textbf{设计动机与边界：} 相对原始 3DGS 以独立高斯和图像重建误差为主的设定，Street Gaussians 把问题落在以下可实现的改动上：Street Gaussians的显式场景表示：(1)背景模型：在世界坐标系中用点云表示，每个点配备3D高斯（位置、协方差、不透明度、球谐系数）和语义逻辑值。协方差由缩放和旋转参数恢复。(2)前景物体模型：每辆车用独立点云表示，位于物体坐标系中，配备可优化的跟踪位姿。为处理时变外观（如车身反射变化），引入4D球谐模型：用时间序列函数预测任意时间步的球谐系数，而非静态系数。(3)场景合成：通过跟踪位姿将前景物体变换到世界坐标系，与背景高斯合并后进行可微光栅化。支持前景物体的分解和编辑操作。增强策略包括跟踪位姿优化、点云初始化和天空建模。其设计目的不是单独增加一个模块，而是用上述表示、监督或渲染约束消除原流程在目标任务中的具体失配。Street Gaussians提出新的显式场景表示，将动态城市场景分解为静态背景和前景车辆，分别用配备语义信息和高斯参数的点云建模。前景车辆使用可优化跟踪位姿和4D球谐模型表示动态外观。半小时训练即可在135FPS下实时渲染（1066x1600），支持场景编辑，在KITTI和Waymo上持续超越SOTA。适用时还应保留论文给出的边界：依赖离线跟踪器提供车辆位姿；假设刚性车辆模型，不支持可变形物体；天空建模较简化；对极端天气和光照变化的处理有限
\textbf{方法信息：} Li Auto 参与了这篇论文；官方代码已公开：\url{https://github.com/zju3dv/street_gaussians}，仓库约 1364 stars；Semantic Scholar 约 246 次引用。



% METHOD: 2502.15635
\paragraph{Para-Lane（跨车道 NVS 基准）（2025）}
Para-Lane~\cite{2502_15635} 是数据集与基准。它通过多遍采集记录不同车道的关联序列，提供前视、环视、LiDAR 和动态物体标签，并设计多传感器位姿对齐与质量评估方案。其意义在于，自动驾驶 NVS 的真正难点不是原轨迹重建，而是跨车道和不同距离的视角迁移。相对 ReconDreamer、VEGS 这类算法，Para-Lane 提供评估场景：方法能否在真实多车道扫描中保持几何和外观一致。

落到使用方式上，研究者可以把某一车道采集的图像和 LiDAR 作为训练输入，再在平行车道或更远距离的相机轨迹上测试渲染质量；动态物体标签用于区分可重建静态背景和需要单独处理的车辆、行人。它并不提出新的高斯优化损失，而是把跨车道位姿求解、传感器对齐和场景划分做成基准条件，让 3DGS、NeRF 或混合方法在同一难度下比较。
\textbf{实现细节：} DeSiRe-GS（2024.11） 修改了 3DGS 的动态场景处理环节——用两阶段自监督分解管线替代 $S^{}3$Gaussian 的统一变形场方案——静态重建能力本身被用作运动检测器。核心观察极为简洁：原始 3DGS 不包含时间信息，因此能准确重建静态区域（停放的车辆、建筑物），但在动态区域会产生伪影（漂浮的"鬼影"点），导致渲染图像与真值之间出现显著差异。第一阶段：2D 运动掩码提取。问题：如何不依赖任何标注来识别动态区域？方案：直接利用 3DGS 的重建失败作为动态性信号。具体而言，用预训练基础模型分别提取渲染图像 $\hat{I}$ 和真值图像 $I$ 的特征 $\hat{F}$ 和 $F$，计算逐像素不相似度： 不相似度 $D$ 在静态区域趋近 0（重建成功），在动态区域趋近 1（重建失败），且该计算不涉及任何可学习参数。但 DeSiRe-GS 并未对 $D$ 直接阈值化，而是引入一个 MLP 解码器预测动态概率 $\hat{\mathcal{D}} \in \mathbb{R}^{H \times W}$，以特征中的丰富语义信息为输入、以 $D$ 为监督信号： 即鼓励解码器在 $D$ 高的区域输出高值。最终二值掩码 $M = \mathbb{I}(\hat{\mathcal{D}} > \tau)$（$\tau$ 为固定阈值）。关键设计：联合优化图像渲染与掩码预测——用掩码遮蔽动态区域后的渲染损失 $L_{\mathrm{masked-render}} = M \odot \|\hat{I} - I\|$ 反过来让动态区域的差异更显著，形成互惠增强。第一阶段训练共 30,000 次迭代，运动解码器从第 6,000 次迭代开始训练。第二阶段：掩码蒸馏至高斯空间。问题：第一阶段的掩码仅存在于 2D 图像空间且依赖真值图像，无法用于新视角合成。方案：以可微分方式将 2D 运动信息蒸馏到 3D 高斯空间。DeSiRe-GS 采用 PVG（Periodic Vibration Gaussian）这一时间变化高斯表示——每个高斯附加速度 $v$，渲染时生成速度图 $V \in \mathbb{R}^{H \times W}$。
\textbf{方法信息：} Alibaba、Baidu 参与了这篇论文。
% METHOD: 2503.10604
\paragraph{MuDG（多模态扩散监督 3DGS）（2025）}
MuDG~\cite{2503_10604} 把可控多模态扩散模型接到城市 3DGS。它先用 LiDAR 点云融合得到稀疏 RGB-depth 条件，再从虚拟视角投影这些条件，由扩散模型生成 RGB、depth 和 semantic 多模态输出。3DGS 模块用这些虚拟视角监督，最小化渲染和扩散重建之间的光度与几何误差。动态和静态通过 bounding box tracking 分离。相对 SGD 只用扩散图像或深度引导，MuDG 同时把颜色、深度和语义作为生成监督。
\textbf{实现细节：} 多个工作从不同角度丰富了街景重建工具箱。MagicRoad（2025.07） 通过语义引导颜色增强和障碍物修复实现鲁棒道路表面重建。GS-RoadPatching（2025.09） 通过 3D 空间结构匹配实现驾驶场景的替代式修复。MOGS（2025.09） 用单目深度替代 LiDAR 降低传感器依赖。LT-Gaussian（2025.08） 提出高斯地图长期更新方法。LS-Gaussian（2025.07） 实现轻量级流式渲染的算法-硬件协同设计。CE-NPBG（2025.04） 通过连接关系图改善大规模驾驶场景中点云渲染的可扩展性。MuDG（2025.03） 将多模态扩散模型与高斯泼溅结合增强极端视角变化下的渲染鲁棒性。Zheng et al.（2025.03） 提出自适应透明度机制消除运动物体。ReconDreamer（2024.11） 通过在线修复和渐进数据更新实现大偏移渲染。
\textbf{方法信息：} 代码：\url{https://github.com/heiheishuang/mudg}。
\textbf{设计动机与边界：} 相对原始 3DGS 以独立高斯和图像重建误差为主的设定，MuDG 把问题落在以下可实现的改动上：MuDG包含两个核心组件：(1) 可控多模态扩散模型（MDM）：利用LiDAR点云融合得到的稀疏RGB-深度对作为条件输入，通过虚拟视角投影生成任意视角的条件输入，扩散模型通过迭代去噪合成RGB、深度和语义的多模态输出；(2) 3DGS模块：使用扩散模型重建的虚拟视角作为监督，最小化渲染和重建视角间的光度与几何差异。通过bbox跟踪分离动态和静态物体，创建融合LiDAR点云。其设计目的不是单独增加一个模块，而是用上述表示、监督或渲染约束消除原流程在目标任务中的具体失配。提出MuDG框架，首次将可控多模态扩散模型与3DGS结合用于城市场景重建，通过扩散模型生成的新视角多模态监督解决极端视角变化下的渲染鲁棒性问题。适用时还应保留论文给出的边界：依赖LiDAR点云和bbox跟踪质量；扩散模型推理增加计算开销；多模态联合训练复杂度高。
% METHOD: 2504.00763
\paragraph{UnIRe（4D 超点无监督实例分解）（2025）}
UnIRe~\cite{2504_00763} 在没有手动标注的情况下分解动态实例。它先对每帧 LiDAR 点去地面并聚类，再用自监督 scene flow 建立跨帧聚类匹配，形成 4D superpoints。时空相关聚类进一步合并或拆分超点，解决 ID 不一致和过度分解问题。每个 4D 超点初始化一个动态实例的规范 3DGS 和 deformation field，背景用静态 3DGS。2D 和 3D 空间平滑正则提升时间一致性。相对 SplatFlow 的流场，UnIRe 更强调实例级初始化。
\textbf{实现细节：} UnIRe（2025.04） 提出 4D 超点（4D superpoints）表示——在 4D 时空空间中聚类多帧 LiDAR 点，基于时空关联实现无监督实例分离。具体而言，先对各帧 LiDAR 点独立运行 DBSCAN 获取逐帧聚类，再用自监督场景流建立帧间簇对应关系（匹配场景流对齐后落入同一簇的点数最多的簇对），处理消失、出现和分裂三种变换后得到在时间上保持一致身份的 4D 超点。为解决分裂导致的过分解问题，UnIRe 进一步利用时空相似度矩阵（融合运动方向余弦相似度和空间邻近度）对 4D 超点进行二次聚类合并。最终 4D 超点为规范空间初始化和逐点变形场提供精确的时空锚点，训练动态 3DGS 无需边界框或物体标注。
\textbf{实现补充：} 实现上还要核对输入、输出、参数、特征、初始化、训练损失、正则、渲染或推理接口，以及优化中的剪枝、分裂、采样、编码或解码步骤；这些抓手决定该方法到底改了 3DGS 管线的哪一环。
\textbf{方法信息：} Huawei 参与了这篇论文。
\textbf{设计动机与边界：} 相对原始 3DGS 以独立高斯和图像重建误差为主的设定，UnIRe 把问题落在以下可实现的改动上：UnIRe框架包含两个核心组件。第一，4D超点初始化：首先对每帧LiDAR点去除地面点后用DBSCAN聚类，然后用自监督场景流（LetItFlow）估计帧间流。基于流的对应关系建立跨帧聚类匹配，形成4D超点S={C\^{}t}\_{t=1}\^{}T（同一实例在多帧中的聚类集合）。通过时空相关性聚类进一步合并/分割4D超点，解决不一致聚类ID、过度分解和不足分解三个问题。最终每个4D超点对应一个动态实例。第二，动态场景重建：利用4D超点为每个动态实例提供规范化空间和逐点变形的初始化。静态背景用标准3DGS表示，每个动态实例拥有独立的规范化3DGS和变形场。引入2D和3D平滑正则化防止过拟合并提升时间一致性。整个方法仅需RGB图像和LiDAR点云，无需边界框或物体模板标注。其设计目的不是单独增加一个模块，而是用上述表示、监督或渲染约束消除原流程在目标任务中的具体失配。提出UnIRe，一种基于3DGS的方法，仅使用RGB图像和LiDAR点云将场景分解为静态背景和独立动态实例。核心创新是4D超点表示——在4D空间中聚类多帧LiDAR点，基于时空相关性实现无监督实例分离。4D超点为分解的4D初始化提供空间和时间初始化，训练动态3DGS时无需边界框或物体模板。引入2D和3D空间平滑正则化策略提升时间一致性。适用时还应保留论文给出的边界：4D超点聚类质量受LiDAR点密度和场景流精度影响；对长时间遮挡的实例恢复能力有限；未处理实例间的交互关系建模。
% METHOD: 2504.19557
\paragraph{CE-NPBG（连接增强神经点图形）（2025）}
CE-NPBG~\cite{2504_19557} 本身是 neural point-based graphics，但也展示了如何增强 3DGS 大规模街景。它建立 RGB 图像与 3D 点之间的 connection graph，通过贪婪分组和剪枝，只检索当前视角相关点，而不是投影整个点云地图。图像合成器和多分辨率判别器联合训练，推理时去掉判别器。与 3DGS 集成时，连接性策略提升渲染质量和可扩展性。相对标准 3DGS 的视锥筛选，它更强调图像-点可见连接关系。
\textbf{实现细节：} 多个工作从不同角度丰富了街景重建工具箱。MagicRoad（2025.07） 通过语义引导颜色增强和障碍物修复实现鲁棒道路表面重建。GS-RoadPatching（2025.09） 通过 3D 空间结构匹配实现驾驶场景的替代式修复。MOGS（2025.09） 用单目深度替代 LiDAR 降低传感器依赖。LT-Gaussian（2025.08） 提出高斯地图长期更新方法。LS-Gaussian（2025.07） 实现轻量级流式渲染的算法-硬件协同设计。CE-NPBG（2025.04） 通过连接关系图改善大规模驾驶场景中点云渲染的可扩展性。MuDG（2025.03） 将多模态扩散模型与高斯泼溅结合增强极端视角变化下的渲染鲁棒性。Zheng et al.（2025.03） 提出自适应透明度机制消除运动物体。ReconDreamer（2024.11） 通过在线修复和渐进数据更新实现大偏移渲染。
\textbf{方法信息：} Huawei 参与了这篇论文。
\textbf{设计动机与边界：} 相对原始 3DGS 以独立高斯和图像重建误差为主的设定，CE-NPBG 把问题落在以下可实现的改动上：CE-NPBG利用两种模态（姿态图像和同步LiDAR点云）进行新视角合成。核心创新包括：(1) 连接关系图：建立RGB图像和3D点之间的连接关系，通过贪婪阶段（将3D点与对应图像分组）和剪枝阶段建立关联。仅从大型3D点云地图中检索与当前视角相关的点用于渲染，而非投影整个地图。(2) 联合对抗和点光栅化训练：为增强神经描述符的编码效果，将图像合成器网络与多分辨率判别器配对训练。推理时解耦判别器。(3) 与3DGS集成：将连接性增强策略集成到3DGS中，展示在渲染质量和可扩展性方面的提升。其设计目的不是单独增加一个模块，而是用上述表示、监督或渲染约束消除原流程在目标任务中的具体失配。提出CE-NPBG，通过建立外观和几何之间的连接关系图解决可见性不匹配问题，显著提升大规模自动驾驶场景的渲染质量和可扩展性，并集成到3DGS中以展示其增强效果。适用时还应保留论文给出的边界：依赖LiDAR点云作为几何输入；连接关系图的构建需要一定的预处理时间；对极端稀疏点云的效果可能受限。
% METHOD: 2506.22099
\paragraph{BézierGS（贝塞尔曲线动态轨迹）（2025）}

BézierGS~\cite{2506_22099} 用可学习 Bezier curves 表示动态高斯运动轨迹。每个动态基元的时变位置由控制点定义，曲线参数从数据中优化，可自动补偿动态物体标注误差。方法同时建模位置和速度，并加入动态物体渲染监督和曲线间一致性，保证动静分离合理。该方法提供代码，仓库关注度较高，便于复核其轨迹参数化实现。相对 SplatFlow 的连续神经流场，BezierGS 的优点是轨迹显式、低维、可解释，适合自动驾驶中相对规则的车辆运动。
\textbf{实现细节：} B\'ezierGS（2025.06） 修改了 3DGS 的动态物体运动建模环节——用可学习贝塞尔曲线轨迹替代全局变形场（如 HexPlane）的隐式运动表示——静态背景的标准 3DGS 重建保持不变。Baseline 做法与问题。此前动态街景方法对运动轨迹的建模可分为两类。第二类是自监督方法（$S^{}3$Gaussian），用 HexPlane 等时空分解网络隐式建模运动轨迹——将 4D 时空网格分解为多分辨率特征平面，通过网络预测每帧变形。这种隐式表示有两个困难：(1) 全局变形场的自由度极高，优化困难，难以保证轨迹的物理合理性；(2) 无法显式约束同一物体上不同高斯基元的运动一致性，容易出现某些高斯脱离物体整体运动的"漂浮"现象。PVG 虽尝试用周期振动拼接长轨迹，但振动模式和实际运动不匹配，且分段拼接无法利用同一物体跨时间的完整一致性。核心改动——可学习贝塞尔曲线轨迹。B\'ezierGS 的核心思路极其直观：与其让神经网络隐式学习运动，不如直接用一条参数化曲线拟合物体运动轨迹。就像用一根弹性绳穿过若干锚点来勾勒平滑路径一样，贝塞尔曲线用少量控制点就能精确描述车辆从起点到终点的完整运动，曲线的平滑性天然保证运动没有突变和抖动。
\textbf{实现补充：} 实现上还要核对输入、输出、参数、特征、初始化、训练损失、正则、渲染或推理接口，以及优化中的剪枝、分裂、采样、编码或解码步骤；这些抓手决定该方法到底改了 3DGS 管线的哪一环。
\textbf{设计动机与边界：} 相对原始 3DGS 以独立高斯和图像重建误差为主的设定，BézierGS 把问题落在以下可实现的改动上：核心方法：(1)可学习贝塞尔曲线轨迹：使用n阶贝塞尔曲线参数化动态高斯基元的运动轨迹，由n+1个控制点定义。曲线完全可学习，自动从数据中优化轨迹，能补偿动态物体的标注误差。(2)动静分离：静态高斯构建背景，动态高斯通过贝塞尔曲线定义时变位置。动态物体通过自监督学习自动分离。(3)额外监督与约束：引入动态物体渲染额外监督和曲线间一致性约束，确保合理准确的场景元素分离和重建。(4)速度建模：不仅建模位置轨迹，还建模动态物体的速度信息，支持任意时间戳的新视角合成。其设计目的不是单独增加一个模块，而是用上述表示、监督或渲染约束消除原流程在目标任务中的具体失配。提出使用可学习贝塞尔曲线建模动态物体运动轨迹，自动纠正姿态标注误差，无需高精度标注即可实现动静分离和高质量重建，在Waymo和nuPlan上超越SOTA。适用时还应保留论文给出的边界：贝塞尔曲线阶数选择影响轨迹精度；对高度非线性轨迹的表达能力有限；依赖3D边界框检测。
\textbf{方法信息：} 代码：\url{https://github.com/fudan-zvg/beziergs}；GitHub 关注度较高，约 138 stars。



% METHOD: 2506.22800
\paragraph{RGE-GS（奖励引导的扩散外扩重建）（2025.06）}

RGE-GS~\cite{2506_22800} 处理单次扫描无法完整覆盖道路结构的问题。它先用原始视图训练基础 GS，再用扩散模型从扩展视点生成伪图像；关键不在于盲目相信扩散，而是训练 reward network 为像素生成一致性打分，只把空间稳定的生成模式保留下来。差异化训练策略根据场景收敛指标自动调节 Gaussian 优化进度。相对 SGD 直接用扩散 SDS 正则，RGE-GS 加了奖励筛选层，减少扩散幻觉被写入 3DGS。
\textbf{实现细节：} RGE-GS（2025.06） 与 VEGS 同样利用扩散先验增强外推视角渲染，但采用两轮重建策略——第一轮从原始视角真值图像重建初始 GS 场 $G_o$，第二轮利用扩散生成的外推视角伪图像扩展场景覆盖。与 VEGS 的差异。VEGS 将扩散先验作为去噪得分匹配监督直接引导高斯协方差优化；RGE-GS 引入轻量奖励网络 $F_c$（3 层 U-Net）生成逐像素置信度图 $C_e \in [0,1]^{H \times W}$，以对抗验证方式动态调节扩散生成图像的监督贡献。奖励网络通过重投影误差损失 $L_{reproj} = \frac{1}{HW}\|C_e \odot (\mathrm{Proj}(P) - I_e)\|_2$（将着色 LiDAR 点云投影到外推视角提供粗监督）和二值化正则化 $L_{reg} = \sum C_e(1-C_e)/(HW)$（鼓励置信度集中在 0 或 1 避免模糊预测）联合训练。差异化训练策略。
\textbf{实现补充：} 为避免只把它读成名称或数据点，还应把实现接口明确拆开：输入是什么，输出是高斯参数、图像、轨迹、标注还是评测指标，哪些参数、网络模块、训练损失或渲染接口真正接到原始 3DGS / Scaffold-GS / 4D-GS 管线上，以及剪枝、分裂、采样、解码、优化或评测协议在哪一步发生。若该条目是survey、benchmark 或 dataset，这里的“实现”对应论文集合、任务划分、相机/轨迹/标注组织、评价指标和基线配置，而不是虚构新的 rasterizer；若是算法条目，则对应高斯位置、尺度、旋转、不透明度、特征或 MLP 解码如何进入训练和渲染流程。
\textbf{方法信息：} Alibaba 参与了这篇论文；代码：\url{https://github.com/cn-adlab/rge-gs}。



% METHOD: 2507.03886
\paragraph{ArmGS（多粒度动态外观细化）（2025）}
ArmGS~\cite{2507_03886} 不主要改几何分解，而是处理动态城市场景的外观变化。它在复合 Gaussian scene 中加入三层外观建模：局部高斯级用每帧低维嵌入和 Hash 网格位置编码学习高斯变换；全局图像级用网络校正不同相机视角曝光和色彩；动态 actor 级为每个车辆或行人学习特定外观变换。相对 Street Gaussians 的 4D SH，ArmGS 把外观变化拆成局部、全局和 actor 三个粒度。
\textbf{实现细节：} DENSER（2024.09） 引入小波动态估计 SH 基增强远距离动态物体外观，跨多帧致密化点云加速收敛。VDG（2024.06） 首次集成自监督视觉里程计实现仅 RGB 输入的无位姿动态高斯。ArmGS（2025.07） 通过多粒度外观精化（局部高斯级、全局图像级、动态角色级）提升帧间和相机间精细变化建模。2D 先验与 SDF 引导（2025.10） 结合 SDF 与 3DGS 实现无 LiDAR 城市场景渲染和场景编辑。SRIF（2024.09） 利用扩散图像变形和流估计实现语义形状注册。
\textbf{方法信息：} Huawei 参与了这篇论文。
\textbf{设计动机与边界：} 相对原始 3DGS 以独立高斯和图像重建误差为主的设定，ArmGS 把问题落在以下可实现的改动上：ArmGS构建基于3DGS的复合驾驶场景模型（背景高斯、动态actor高斯、天空高斯），并从三个粒度进行外观细化：(1) 局部复合高斯级：构建每帧的低维可学习嵌入，通过多分辨率Hash网格提取高斯位置编码，学习变换参数细化复合高斯以建模帧间场景变化；(2) 全局图像级：通过全局外观变换网络建模不同相机视角的全局外观变化（如曝光差异），学习图像级颜色校正参数；(3) 动态actor级：对每个动态actor（如车辆）单独建模外观变化，通过actor特定的变换参数调整。三个级别的外观细化参数联合优化，使重建结果在保持全局一致性的同时保留局部细节。其设计目的不是单独增加一个模块，而是用上述表示、监督或渲染约束消除原流程在目标任务中的具体失配。提出ArmGS，通过多粒度外观建模方案（局部高斯级、全局图像级、动态actor级）优化复合高斯变换参数，有效建模帧间和视角间的细粒度外观变化，在Waymo、KITTI等数据集上优于SOTA。适用时还应保留论文给出的边界：多级外观建模增加训练复杂度；对动态actor的检测和跟踪有依赖；长序列场景可能存在累积误差。
% METHOD: 2508.15376
\paragraph{DriveSplat（LOD 背景与对象中心神经高斯）（2025）}
DriveSplat~\cite{2508_15376} 将动态驾驶场景分为静态背景和动态前景。背景采用 scene-aware learnable LOD，把环境分成近、中、远三个深度范围，为不同范围分配不同尺度和密度的高斯；动态前景使用对象中心 neural Gaussian primitives，通过全局刚体变换和两阶段变形处理非刚性运动，先调 anchor 再更新高斯属性。相对 OmniRe 的场景图，DriveSplat 更关注背景多尺度和动态对象的细粒度变形。
\textbf{实现细节：} DriveSplat（2025.08） 提出场景感知可学习 LOD 建模策略显式处理近/中/远深度范围的自适应多尺度高斯分配（近距离使用密集小高斯捕获细节，远距离使用稀疏大高斯覆盖广阔区域）。与 OmniRe 为每种动态参与者设计独立高斯簇不同，DriveSplat 的核心改动是统一的 LOD 驱动分配机制——可学习的细节层次参数根据深度范围自动调节高斯密度和尺度，无需针对不同物体类型手动设计表示策略。动态角色通过全局刚体变换和两阶段变形处理非刚性动态——第一阶段调整锚点位置捕获物体整体运动，第二阶段更新每个高斯的局部偏移建模非刚性变形（如行人姿态变化）。这种分层变形策略兼顾了车辆等刚体物体的运动约束和行人的非刚性特征，在保持渲染质量的同时大幅降低远距离区域的高斯数量。
\textbf{实现补充：} 实现上还要核对输入、输出、参数、特征、初始化、训练损失、正则、渲染或推理接口，以及优化中的剪枝、分裂、采样、编码或解码步骤；这些抓手决定该方法到底改了 3DGS 管线的哪一环。
\textbf{方法信息：} 代码：\url{https://github.com/sokldjs554/DriveSplat}。
\textbf{设计动机与边界：} 相对原始 3DGS 以独立高斯和图像重建误差为主的设定，DriveSplat 把问题落在以下可实现的改动上：方法包含静态背景和动态前景两部分：(1)静态背景采用场景感知可学习LOD策略，将驾驶环境分为近、中、远三个深度范围，为每个范围分配不同尺度和密度的高斯基元，实现自适应多尺度建模；(2)动态前景使用对象中心神经高斯基元，通过全局刚体变换建模物体运动，并设计两阶段变形策略——第一阶段调整锚点位置，第二阶段更新高斯属性以处理非刚性动态；(3)引入来自预训练模型的密集深度和表面法线先验作为辅助监督正则化优化过程。其设计目的不是单独增加一个模块，而是用上述表示、监督或渲染约束消除原流程在目标任务中的具体失配。提出DriveSplat，首个统一神经高斯框架用于动态驾驶场景重建。引入场景感知可学习LOD建模策略，显式处理近、中、远三个深度范围的背景几何，实现自适应多尺度高斯分配。对动态物体采用对象中心神经高斯基元，通过全局刚体变换和两阶段变形（先调整锚点再更新高斯）处理非刚性运动。在Waymo和KITTI上达到SOTA。适用时还应保留论文给出的边界：依赖预训练深度和法线模型的质量；对极端非刚性运动（如行人复杂姿态变化）的建模仍有改进空间；计算开销相对较大。
% METHOD: 2509.06685
\paragraph{MOGS（对象锚定深度替代 LiDAR）（2025）}
MOGS~\cite{2509_06685} 解决单目大场景 3DGS 缺少稠密度量深度的问题。它不依赖 LiDAR，而是利用视觉惯性 SfM 的稀疏点作为锚点，通过对象级语义形状先验把稀疏约束扩展为 dense depth。多尺度形状共识模块把语义区域合并成 SfM 支持较好的粗对象，并拟合平面、圆柱等参数形状；跨对象深度细化再加入共面、平行、接触等几何关系。相对 VEGS 的 LiDAR 初始化，MOGS 的增量是对象先验驱动的单目度量几何。
\textbf{实现细节：} 多个工作从不同角度丰富了街景重建工具箱。MagicRoad（2025.07） 通过语义引导颜色增强和障碍物修复实现鲁棒道路表面重建。GS-RoadPatching（2025.09） 通过 3D 空间结构匹配实现驾驶场景的替代式修复。MOGS（2025.09） 用单目深度替代 LiDAR 降低传感器依赖。LT-Gaussian（2025.08） 提出高斯地图长期更新方法。LS-Gaussian（2025.07） 实现轻量级流式渲染的算法-硬件协同设计。CE-NPBG（2025.04） 通过连接关系图改善大规模驾驶场景中点云渲染的可扩展性。MuDG（2025.03） 将多模态扩散模型与高斯泼溅结合增强极端视角变化下的渲染鲁棒性。Zheng et al.（2025.03） 提出自适应透明度机制消除运动物体。ReconDreamer（2024.11） 通过在线修复和渐进数据更新实现大偏移渲染。
\textbf{实现补充：} 实现上还要核对输入、输出、参数、特征、初始化、训练损失、正则、渲染或推理接口，以及优化中的剪枝、分裂、采样、编码或解码步骤；这些抓手决定该方法到底改了 3DGS 管线的哪一环。
\textbf{方法信息：} 代码：\url{https://github.com/clarencezsk/mogs}。
\textbf{设计动机与边界：} 相对原始 3DGS 以独立高斯和图像重建误差为主的设定，MOGS 把问题落在以下可实现的改动上：MOGS框架的核心是对象锚定度量策略：(1)多尺度形状共识模块：利用图像语义将场景分割为对象区域，自适应地将小段合并为SfM最佳支持的粗对象，并拟合参数形状模型(平面、圆柱等)。解决了SfM覆盖不足(点云聚类在角点和高纹理边界)的问题。(2)跨对象深度细化模块：在组合目标(几何一致性+先验锚定+边缘感知平滑)下逐像素优化深度，解决跨对象几何不一致问题(如平行性、共面性、与地面的接触)。方法使用视觉惯性SfM前端提供稀疏但可靠的度量线索，通过对象级形状先验将稀疏锚点放大为稠密深度支持，无需昂贵的LiDAR传感器。其设计目的不是单独增加一个模块，而是用上述表示、监督或渲染约束消除原流程在目标任务中的具体失配。提出MOGS，首个用对象锚定稠密深度替代LiDAR的单目3DGS大场景框架。利用图像语义假设逐对象形状先验，用稀疏VI-SfM点作为锚点，通过多尺度形状共识和跨对象深度细化模块将稀疏约束传播为稠密深度。适用时还应保留论文给出的边界：依赖语义分割的准确性；参数形状模型可能不适用于所有对象类型；长距离低纹理区域的深度估计仍有挑战。
% METHOD: 2509.19937
\paragraph{GS-RoadPatching（直接在高斯空间修复道路）（2025）}
GS-RoadPatching~\cite{2509_19937} 不走 2D 扩散修复再重训的路线，而是直接在 3DGS 模态中做替代式 patching。它先在高斯场中嵌入特征，定义多尺度 patch 度量抽象局部上下文；再利用驾驶场景中道路和结构高度重复的特点，在 3D 空间搜索相似补丁；最后通过替换和融合优化保持视觉和几何和谐。相对 FaithFusion 依赖扩散生成，GS-RoadPatching 的改变是把 PatchMatch 式搜索搬到高斯空间，避免时空生成不一致和额外重建轮次。
\textbf{实现细节：} 多个工作从不同角度丰富了街景重建工具箱。MagicRoad（2025.07） 通过语义引导颜色增强和障碍物修复实现鲁棒道路表面重建。GS-RoadPatching（2025.09） 通过 3D 空间结构匹配实现驾驶场景的替代式修复。MOGS（2025.09） 用单目深度替代 LiDAR 降低传感器依赖。LT-Gaussian（2025.08） 提出高斯地图长期更新方法。LS-Gaussian（2025.07） 实现轻量级流式渲染的算法-硬件协同设计。CE-NPBG（2025.04） 通过连接关系图改善大规模驾驶场景中点云渲染的可扩展性。MuDG（2025.03） 将多模态扩散模型与高斯泼溅结合增强极端视角变化下的渲染鲁棒性。Zheng et al.（2025.03） 提出自适应透明度机制消除运动物体。ReconDreamer（2024.11） 通过在线修复和渐进数据更新实现大偏移渲染。
\textbf{方法信息：} Alibaba、Baidu 参与了这篇论文。
\textbf{设计动机与边界：} 相对原始 3DGS 以独立高斯和图像重建误差为主的设定，GS-RoadPatching 把问题落在以下可实现的改动上：GS-RoadPatching方法包含三个核心组件：（1）特征嵌入3DGS场景构建：在3DGS场景中嵌入特征，提出多尺度补丁度量方法抽象不同尺度的局部上下文。（2）3D结构搜索：在3D空间中有效搜索候选补丁，利用驾驶场景高度重复模式的3DGS特征空间相似性进行结构匹配。受PatchMatch启发，将2D图像修复中的搜索-度量-替换策略迁移到3D高斯空间。（3）替换-融合优化：提出简单有效的替换和融合优化策略，确保视觉和谐。核心优势是直接在紧凑的3D高斯场景上操作，避免了重新生成原始训练数据进行额外重建轮次，消除了对2D视频扩散模型时空一致性的依赖。在一般场景中也验证了适用性。其设计目的不是单独增加一个模块，而是用上述表示、监督或渲染约束消除原流程在目标任务中的具体失配。提出GS-RoadPatching，首创直接在3DGS模态中进行替代式场景修复，无需2D扩散/GAN生成和重新训练高斯。利用驾驶场景中高度重复模式的隐式3DGS特征空间相似性进行结构匹配，提出多尺度补丁度量和3D结构搜索方法，以及简单的替换-融合优化。适用时还应保留论文给出的边界：依赖驾驶场景的结构重复性；对非结构化场景可能不适用；补丁搜索在大场景中可能耗时。
% METHOD: 2511.19235
\paragraph{IDSplat（实例分解与可学习运动轨迹）（2025.11）}
IDSplat~\cite{2511_19235} 通过 explicit instance decomposition 自监督建模动态驾驶场景。每个高斯存储 RGB 特征、视图或 LiDAR 特征和 instance ID，动态高斯通过 $SE(3)$ pose 变换，背景高斯在世界坐标中定义。实例分解来自 Grounded-SAM-2 的零样本类提示视频追踪，LiDAR 点投影到图像后获得实例 ID，并通过 DBSCAN 过滤。初始轨迹由实例 LiDAR 点配准得到，协调转弯平滑处理车辆快速转向的姿态不连续。相对 UnIRe，它更多利用视觉语言分割和实例跟踪。
\textbf{实现细节：} IDSplat（2025.11） 将动态物体建模为经历 SE(3)（特殊欧氏群，描述三维空间中的刚体变换）刚体变换的连贯实例，每个实例的所有高斯共享同一变换 $T_{z,t} \in \mathrm{SE}(3)$，从规范坐标系变换到世界坐标系。实例分解方面，利用 Grounded-SAM-2（基于语言提示的通用视频分割模型）生成零样本语言接地的视频实例掩码，将 2D 掩码投影到 LiDAR 点云后经腐蚀和 DBSCAN 过滤锚定到 3D。轨迹估计方面，以实例最密集帧的包围盒中心定义规范坐标系，利用 DINOv3（DINO 系列第三代自监督视觉Transformer）特征建立帧间对应关系并通过 RANSAC+Umeyama 估计刚体变换；再通过协调转弯（Coordinated-Turn）平滑优化构建物理一致的运动轨迹，迭代剔除离群测量值提升鲁棒性。该方法在 Waymo 上无需重训练即泛化到多样序列。
\textbf{实现补充：} 实现上还要核对输入、输出、参数、特征、初始化、训练损失、正则、渲染或推理接口，以及优化中的剪枝、分裂、采样、编码或解码步骤；这些抓手决定该方法到底改了 3DGS 管线的哪一环。
\textbf{方法信息：} 代码：\url{https://github.com/zenseact/idsplat}。
\textbf{设计动机与边界：} 相对原始 3DGS 以独立高斯和图像重建误差为主的设定，IDSplat 把问题落在以下可实现的改动上：IDSplat包含三个核心步骤：(1)场景表示——每个高斯存储不透明度、均值、协方差、RGB特征、视图/激光雷达特征和实例ID，动态高斯通过SE(3)姿态进行刚体变换，背景高斯直接在世界坐标系中定义；(2)实例分解——使用Grounded-SAM-2以类提示方式生成实例掩码（零样本），将激光雷达点投射到图像平面分配实例ID，通过腐蚀掩码和DBSCAN聚类过滤异常点；(3)轨迹估计——通过配准各时间步的实例激光雷达点估计初始轨迹，引入协调转弯平滑方案处理快速转向时的姿态不连续性，使用特征对应进行姿态细化，最终联合优化物体姿态和高斯参数。其设计目的不是单独增加一个模块，而是用上述表示、监督或渲染约束消除原流程在目标任务中的具体失配。提出IDSplat自监督框架，在无需人工标注的情况下实现动态场景的显式实例分解和可学习运动轨迹。利用零样本语言引导视频追踪进行实例分解，协调转弯平滑方案获得时间物理一致的运动轨迹，在Waymo数据集上实现有竞争力的重建质量。适用时还应保留论文给出的边界：依赖Grounded-SAM-2的追踪质量；刚体假设不适用于变形物体；快速运动可能导致追踪失败
% METHOD: 2511.21113
\paragraph{FaithFusion（像素级信息增益融合重建与生成）（2025.11）}

FaithFusion~\cite{2511_21113} 也融合 3DGS 和扩散，但它用 expected information gain 决定哪里需要生成。方法通过 Laplace 近似估计 3DGS 参数后验和像素级不确定性，EIG 在生成分支中引导扩散模型修复高不确定区域，在重建分支中又作为像素权重，把生成结果选择性蒸馏回 3DGS。整体采用 render-restoration-feedback 渐进融合。相对 RGE-GS 的 reward network，FaithFusion 的关键是用重建不确定性定义“该不该相信扩散”的空间策略。
\textbf{实现细节：} MagicDrive3D（2024.05） 首先训练多视角视频生成模型再构建 3DGS，支持多条件控制。SGD（2024.03） 利用扩散先验和 LiDAR 深度增强 3DGS 的街景外推能力。StreetUnveiler（2024.05） 从拥挤观测中恢复空街道。FaithFusion（2025.11） 通过像素级期望信息增益融合 3DGS 和扩散模型。Drone-assisted Road GS（2024.08） 引入跨视图不确定性使无人机数据辅助地面视角合成。LR-SGS（2026.03） 利用 LiDAR 反射率引导显著高斯表示增强复杂光照场景。Omran et al.（2025.10） 提出高斯泼溅与扩散模型混合的两阶段新视角合成方案。Hybrid Gaussian（2025.03） 和 Hybrid Rendering（2025.03） 分别探索了混合表示和混合渲染策略。
\textbf{实现补充：} 实现上还要核对输入、输出、参数、特征、初始化、训练损失、正则、渲染或推理接口，以及优化中的剪枝、分裂、采样、编码或解码步骤；这些抓手决定该方法到底改了 3DGS 管线的哪一环。
\textbf{设计动机与边界：} 相对原始 3DGS 以独立高斯和图像重建误差为主的设定，FaithFusion 把问题落在以下可实现的改动上：FaithFusion是一个即插即用的3DGS-扩散融合框架，核心创新在于引入像素级期望信息增益(EIG)作为统一空间策略。首先，基于Laplace近似计算3DGS参数后验分布的Hessian矩阵，推导出像素级EIG估计，量化每个像素的重建不确定性。在生成分支，EIG作为空间先验引导扩散模型聚焦于高不确定性区域（如遮挡区域或大视角偏移处的不可见区域）进行修复。在重建分支，EIG作为像素级加权因子，选择性地将扩散生成的编辑结果蒸馏回3DGS，避免过度修复和几何漂移。整体采用渐进式融合方案(render-restoration-feedback)，框架可无缝集成到主流街景3DGS系统（如StreetGS、DeformableGS等）中。关键设计包括：(1)将3DGS可微渲染器与Laplace近似EIG紧密耦合；(2)双分支EIG引导（生成+重建）；(3)EIG驱动的渐进式知识整合。其设计目的不是单独增加一个模块，而是用上述表示、监督或渲染约束消除原流程在目标任务中的具体失配。提出基于像素级期望信息增益(EIG)的3DGS-扩散模型融合框架，EIG作为统一空间策略引导扩散模型合成高不确定性区域，并通过像素级加权将编辑结果蒸馏回3DGS。该框架即插即用，无需额外先验条件，在Waymo数据集上实现NTA-IoU、NTL-IoU和FID的SOTA性能，在6米车道偏移下仍保持FID 107.47。适用时还应保留论文给出的边界：当前仅在驾驶场景（Waymo）上验证，泛化到其他场景类型的能力未验证；Laplace近似对后验分布的假设可能在极端情况下不够准确；计算EIG引入额外开销；渐进式融合方案需要多次迭代。
\textbf{方法信息：} Huawei 参与了这篇论文。



% METHOD: 2601.01386
\paragraph{ParkGaussian（停车位感知鱼眼重建）（2026）}
ParkGaussian~\cite{2601_01386} 把自动驾驶 3DGS 放到停车场景。它构建 ParkRecon3D 基准，包含环视鱼眼相机、地下停车场轨迹和车位角点标注。方法基于 3DGUT 适配鱼眼投影，并引入 IPM 投影和 parking-slot-aware reconstruction：利用车位检测模型提取车位区域 mask，在训练中提高车位区域损失权重。相对 UniGaussian 的通用鱼眼模型，ParkGaussian 的增量是把下游停车位检测语义直接作为重建关注区域。
\textbf{实现细节：} ParkGaussian（2026.01） 首次将 3DGS 引入自动泊车场景重建，创建 ParkRecon3D 基准数据集（约 $220\times110$ m$^{}2$ 地下车库、430+ 车位、四路环视鱼眼相机 $1280\times960$@10Hz、超过 40K 同步帧和 60K 手动标注车位角点）。技术核心包括两方面：其一，采用基于无迹变换（Unscented Transform）的投影替代传统 EWA 泼溅的雅可比线性化，直接在鱼眼图像上训练 3DGS 而无需逐相机推导雅可比，显著提升强畸变下的几何稳定性；其二，设计可微分逆透视映射（IPM）模块将渲染鱼眼图转换为鸟瞰图，并引入车位感知重建策略——利用预训练车位检测器 DMPR-PS 和 GCN-Parking 的冻结特征图构建车位角点与边缘的混合权重图（温度 $T=0.5$、混合系数 $\alpha=0.8$），反向传播监督信号至 3DGS 优化，使重建结果在车位区域的结构保真度显著提升。
\textbf{方法信息：} Xiaomi 参与了这篇论文；代码仓库：\url{https://github.com/wm-research/ParkGaussian}。
\textbf{设计动机与边界：} 相对原始 3DGS 以独立高斯和图像重建误差为主的设定，ParkGaussian 把问题落在以下可实现的改动上：方法包含两个核心组件：(1) ParkRecon3D基准：基于AVM-SLAM数据集扩展，包含4个环视鱼眼相机数据，约220×110米地下停车场、430+车位，使用COLMAP标定外参，手动标注车位角点；(2) ParkGaussian框架：使用3DGUT扩展3DGS以适应鱼眼相机投影，引入IPM投影模块和停车位感知重建策略。停车位感知策略利用现有车位检测方法提取车位区域掩码，在3DGS训练中增加车位区域的额外损失权重，增强车位区域的重建质量。其设计目的不是单独增加一个模块，而是用上述表示、监督或渲染约束消除原流程在目标任务中的具体失配。构建首个停车场景3D重建基准ParkRecon3D（4个鱼眼相机、4万+帧、6万+车位标注），提出ParkGaussian框架整合3DGS进行停车场景重建。核心创新是停车位感知重建策略，利用现有停车感知方法增强车位区域的合成质量，提升重建与下游车位检测的对齐度。适用时还应保留论文给出的边界：基准限于地下停车场场景；依赖COLMAP外参标定；鱼眼相机畸变模型仍有改进空间。
% METHOD: 2601.07540
\paragraph{SetDiff（几何引导集合扩散增强）（2026）}
Enhancing Novel View Synthesis via Geometry Grounded Set Diffusion~\cite{2601_07540} 也处理新视角退化，但它把参考和目标视图看成无序集合。SetDiff 使用 pixel-aligned coordinate maps 和 pose-aware Pluecker ray embeddings 作为显式 3D 条件，set mixer 在 token 级做全局注意力，可同时处理可变数量参考和目标视角。选择性解码保持计算效率。相对 FaithFusion 依赖 3DGS 后验不确定性，SetDiff 更像一个多视角扩散增强器，负责把几何先验注入扩散恢复过程。

使用时，参考相机提供图像、位姿和坐标图，目标相机提供待恢复视角的射线条件；集合扩散网络在统一 token 池里交换跨视角信息，再只解码需要增强的目标视图。它不直接替换 3DGS 表示，而是给驾驶场景 NVS 提供更干净、更几何一致的观测或修复结果，因此可以作为后续高斯重建、轨迹外渲染或数据增强的前端模块。
\textbf{实现细节：} SetDiff（2026.01） 与 VEGS 同样利用扩散模型增强外推视角渲染，但 VEGS 将扩散先验嵌入 3DGS 优化过程，而 SetDiff 作为后处理增强——以 3DGS 渲染的降质图像为输入，通过扩散去噪恢复高频细节。动机。将扩散先验嵌入优化过程需要重新训练整个 3DGS 场景，代价高昂且不灵活；后处理增强方案可以利用任意预训练扩散模型，即插即用地提升任意已重建场景的外推渲染质量。显式 3D 几何先验注入。核心创新是将像素对齐坐标图（每个像素编码其对应的 3D 空间位置）和姿态感知 Pl\"ucker 射线嵌入（编码相机 6DoF 位姿的 6 维表示）作为显式 3D 先验融入集混合器（Set Mixer）。与标准扩散模型仅在潜空间操作不同，SetDiff 在全局标记级注意力中允许同一视角内的像素 token 和跨视角的像素 token 在同一注意力层中交互——参考视角和目标视角的所有像素 token 被展平为统一序列后执行全自注意力，使模型能从参考视角的可靠信息中推断目标视角缺失内容的合理外观。可变视角数量处理。
\textbf{方法信息：} Qualcomm 参与了这篇论文。
\textbf{设计动机与边界：} 相对原始 3DGS 以独立高斯和图像重建误差为主的设定，Enhancing Novel View Synthesis via Geometry Grounded Set Diffusion 把问题落在以下可实现的改动上：SetDiff框架将显式3D先验（像素对齐坐标图和姿态感知Plücker射线嵌入）集成到基于集合的扩散模型中，可联合处理可变数量的参考和目标视角。采用无序集合混合器（set mixer）执行全局token级注意力，支持可扩展的多摄像头增强，通过潜空间监督和选择性解码保持计算效率。该方法能够鲁棒处理遮挡、减少低信号条件下的幻觉，并改善视觉内容恢复的光度保真度。其设计目的不是单独增加一个模块，而是用上述表示、监督或渲染约束消除原流程在目标任务中的具体失配。提出SetDiff，一种几何引导的多视角扩散增强框架，通过像素对齐坐标图和姿态感知Plücker射线嵌入提供显式3D先验，采用无序集合混合器进行全局token级注意力，实现可变数量参考和目标视角的联合处理，显著提升3DGS在新视角合成中的感知保真度和鲁棒性。适用时还应保留论文给出的边界：扩散模型的推理延迟可能影响实时性；主要针对自动驾驶场景评估，其他场景泛化性未验证。
% METHOD: 2601.15951
\paragraph{EVolSplat4D（三分支前馈 4D 城市场景）（2026）}
EVolSplat4D~\cite{2601_15951} 把前馈动态城市场景分成三支。近距静态区域从 3D feature volume 直接预测多视图一致高斯几何，并用遮挡感知 IBR 估计外观；动态车辆在对象中心规范空间中建模，通过 motion-adjusted IBR 聚合时间特征；远场景观用像素对齐分支和跨视图注意力预测。相对 DrivingForward 的统一逐像素预测，EVolSplat4D 的增量是把近景、动态对象和远景分配给不同前馈机制。
\textbf{实现细节：} EVolSplat4D（2026.01） 提出前馈框架统一体积和像素高斯预测——近距离静态区域从 3D 特征体预测一致的 3D 高斯（通过体素化特征网格直接查询），动态角色利用对象中心规范空间和运动调整渲染模块聚合时间特征（在规范空间中建模动态物体的变形，再通过位姿变换到世界坐标），远景由高效逐像素高斯分支覆盖（避免对远处弱纹理区域构建完整 3D 表示的开销）。在 KITTI-360、KITTI、Waymo、PandaSet 上超越所有基线。
\textbf{实现补充：} 实现上还要核对输入、输出、参数、特征、初始化、训练损失、正则、渲染或推理接口，以及优化中的剪枝、分裂、采样、编码或解码步骤；这些抓手决定该方法到底改了 3DGS 管线的哪一环。
\textbf{方法信息：} Huawei 参与了这篇论文。
\textbf{设计动机与边界：} 相对原始 3DGS 以独立高斯和图像重建误差为主的设定，EVolSplat4D 把问题落在以下可实现的改动上：框架将场景分解为三个组件，各有专门的前馈分支：(1)近距体积分支：利用深度先验构建全局点云初始化，通过3D CNN解码多视图一致的3DGS几何参数，再由遮挡感知IBR模块预测外观；(2)动态车辆分支：在规范空间建模每个动态物体，提出运动调整IBR模块聚合时序特征，确保在噪声运动先验下稳定4D重建；(3)远场景观分支：训练带跨视图注意力的U-Net预测像素对齐高斯。三分支融合表示无界街景。其设计目的不是单独增加一个模块，而是用上述表示、监督或渲染约束消除原流程在目标任务中的具体失配。提出EVolSplat4D前馈框架，超越逐像素范式，统一体素基和像素基高斯预测为三个专门分支：近距静态区域从3D特征体直接预测一致3DGS几何；动态物体使用对象中心规范空间和运动调整渲染模块；远场景观用高效像素对齐分支。在KITTI-360、KITTI、Waymo和PandaSet上超越per-scene优化和SOTA前馈基线。适用时还应保留论文给出的边界：三分支系统复杂度较高；依赖LiDAR观测的3D边界框追踪；对未见类型动态物体的泛化能力待验证。
% METHOD: 2602.13549
\paragraph{Nighttime PBR-GS（夜间物理材质高斯）（2026）}
Nighttime Autonomous Driving Scene Reconstruction with Physically-Based Gaussian Splatting~\cite{2602_13549} 关注夜间驾驶。普通 3DGS 把颜色直接烘焙到高斯里，难以解释车灯、反光路面和高光。该方法把渲染方程和 BRDF 材质属性接入复合场景图：静态背景、动态参与者和天空 cubemap 分开建模；漫反射由全局光照模块预测 SH 系数，镜面反射由 anisotropic spherical Gaussians 捕捉高频高光，再经 HDR 到 LDR tone mapping 渲染。相对 Street Gaussians 的外观时间函数，它把夜间问题拆成光照与材质，而不是只拟合颜色。
\textbf{实现细节：} 特定传感器与成像条件方面，ParkGaussian 首次将 3DGS 引入自动泊车场景、用基于无迹变换的投影处理鱼眼畸变；DC-Gaussian 处理行车记录仪视频中挡风玻璃反射，将观测图像分解为透射和反射分量；Kim et al. 将基于物理的渲染引入 3DGS、增强夜间驾驶场景、联合优化 BRDF 材质——三者覆盖"非常规相机—非常规光学—非常规光照"成像异常谱。
\textbf{实现补充：} 实现上还要核对输入、输出、参数、特征、初始化、训练损失、正则、渲染或推理接口，以及优化中的剪枝、分裂、采样、编码或解码步骤；这些抓手决定该方法到底改了 3DGS 管线的哪一环。
\textbf{方法信息：} Huawei 参与了这篇论文。
\textbf{设计动机与边界：} 相对原始 3DGS 以独立高斯和图像重建误差为主的设定，Nighttime Autonomous Driving Scene Reconstruction with Physically-Based Gaussian Splatting 把问题落在以下可实现的改动上：方法核心：1)复合场景图结构：静态背景、动态参与者（刚性用轨迹更新、非刚性用SMPL/变形高斯）和天空cubemap三部分；2)BRDF建模：将渲染方程引入GS，每个高斯配备BRDF材质属性；3)漫反射分量：全局光照模块预测SH系数建模漫反射光照，不依赖环境贴图；4)镜面反射分量：各向异性球面高斯(ASGs)捕捉高频镜面效果；5)HDR到LDR色调映射：通过Reinhard映射后使用标准可微分3DGS光栅化器渲染。解耦光照和材质，更好地处理夜间复杂光照条件。其设计目的不是单独增加一个模块，而是用上述表示、监督或渲染约束消除原流程在目标任务中的具体失配。首次将物理渲染集成到复合场景GS表示中用于夜间自动驾驶场景重建。提出全局光照模块建模漫反射分量、各向异性球面高斯建模镜面反射分量，联合优化BRDF材质属性，实现夜间场景的实时高保真重建。适用时还应保留论文给出的边界：依赖LiDAR点初始化；对极端光照变化（如闪电）处理有限；行人等非刚性物体建模仍有挑战。
% METHOD: 2603.08254
\paragraph{DynamicVGGT（动态点图前馈 4D 重建）（2026）}
DynamicVGGT~\cite{2603_08254} 把 VGGT 从静态 3D 感知扩展到动态 4D。它引入 dynamic point maps，在同一规范框架中联合预测当前和未来帧点图，从点图差异隐式学习点级运动。motion-aware temporal attention 将 patch tokens 与可学习 motion tokens 拼接，显式编码时间相关性。动态 3DGS head 再在 scene flow 监督下预测高斯速度并优化动态几何。相对 EVolSplat4D 的三分支设计，DynamicVGGT 更像把基础 3D 模型改造成统一动态前馈骨干。
\textbf{实现细节：} DynamicVGGT（2026.03） 将 VGGT 从静态 3D 扩展到动态 4D 重建，联合预测当前帧和未来帧的点图（pointmap），引入运动感知时间注意力模块学习运动连续性。与 DrivingScene 的两阶段静态-动态分离策略不同，DynamicVGGT 的核心改动是在 VGGT 的交替注意力块之后插入运动感知时间注意力（MTA）模块——通过可学习运动标记在时间维度上跨帧传递信息，使网络在特征层面建立帧间运动对应关系。未来点预测头（FPH）从当前帧的增强特征直接预测未来帧的 3D 点图，通过时间一致性正则化约束预测位移与真值位移对齐，在点图层面隐式学习运动。动态 3DGS 头通过场景流监督预测每个高斯的 3D 速度向量——将 $t$ 时刻的高斯按预测速度平移到 $t+1$ 时刻，渲染后与 $t+1$ 真值比较提供显式的运动监督信号，与 FPH 的隐式点图位移监督形成互补。
\textbf{方法信息：} Huawei 参与了这篇论文。
\textbf{设计动机与边界：} 相对原始 3DGS 以独立高斯和图像重建误差为主的设定，DynamicVGGT 把问题落在以下可实现的改动上：DynamicVGGT基于VGGT骨干构建，引入三大核心组件：(1)动态点图(DPM)——在VGGT学习的规范框架中直接预测点图，通过联合预测当前帧P\_{v,t}和未来帧P\_{v,t+Δ}的点图，隐式学习点级运动Δ\_{v,t}=P\_{v,t+Δ}-P\_{v,t}，避免依赖外部指定的帧到参考变换；(2)运动感知时间注意力(MTA)——在VGGT的AA块之后，将补丁token与可学习运动token拼接，并行计算所有帧的时间相关性，显式在特征层编码运动线索引导时间注意力关注运动一致区域，避免StreamVGGT中AA块顺序堆叠导致的不稳定训练；(3)动态3DGS头部——从几何先验初始化高斯基元，通过轻量运动编码器编码运动流，使用可学习运动token在场景流监督下预测高斯速度，通过连续3DGS优化精化动态几何。其设计目的不是单独增加一个模块，而是用上述表示、监督或渲染约束消除原流程在目标任务中的具体失配。提出DynamicVGGT，首个将VGGT从静态3D感知扩展到动态4D重建的统一前馈框架。引入动态点图(DPM)机制，通过联合预测当前和未来帧的点图隐式学习动态点表示。提出运动感知时间注意力(MTA)模块和学习运动token，以及动态3DGS头部通过场景流监督显式建模点运动。在自动驾驶数据集上显著超越现有方法。适用时还应保留论文给出的边界：依赖VGGT骨干的质量；运动token容量有限可能影响长时间动态建模；对快速运动物体可能跟踪失败。
% METHOD: 2605.11594
\paragraph{PointForward（点对齐前馈驾驶重建）（2026）}
PointForward~\cite{2605_11594} 反思 DrivingForward 的逐像素高斯范式。逐像素预测容易产生多视角不一致和分层伪影；PointForward 的改变是把高斯锚到 3D points，而不是每个 2D pixel。它还用 3D scene graph 直接在三维空间建模动态实例，动态查询合并到规范空间以保证跨视角一致。相对 DrivingForward 的深度反投影，PointForward 更强调点对齐表示和实例图结构，减少像素域预测带来的几何歧义。
\textbf{实现细节：} PointForward（2026.05） 摒弃 DrivingForward 的逐像素高斯预测范式，转而在世界空间中初始化稀疏 3D 查询点（每个查询是 12 维向量：3D 坐标 + RGB 颜色 + 6 维 Pl\"ucker 射线坐标），通过时空融合聚合多视角信息，在单次前馈中强制跨视图一致性。与 DrivingForward 的差异。DrivingForward 从单帧图像预测逐像素深度，本质是 "2.5D" 表示；PointForward 利用预训练 EfficientNetV2 提取图像特征和冻结的 InfiniDepth 预测密集深度，将像素直接提升到 3D 空间形成稀疏查询集 $Q = Q_s \cup Q_d$（$V \times N_s$ 个静态查询 + $O \times N_d$ 个动态查询），避免了像素级预测中的视角冲突和冗余计算。场景图驱动的动态建模。每个 3D 查询点通过可变形注意力机制从多视角图像特征中采样信息，预测对应的高斯属性（位置 $\mu$、旋转 $r$、缩放 $s$、不透明度 $o$、高维颜色特征 $c$）。利用场景图显式组织运动实例——将落入 3D 边界框的查询标记为动态并变换到实例规范空间（以首个边界框为参考），其余标记为静态，实现实例级运动传播和时间一致动态表示。
\textbf{技术实现补充：} PointForward（2026.05） 摒弃 DrivingForward 的逐像素高斯预测范式，转而在世界空间中初始化稀疏 3D 查询点（每个查询是 12 维向量：3D 坐标 + RGB 颜色 + 6 维 Pl\"ucker 射线坐标），通过时空融合聚合多视角信息，在单次前馈中强制跨视图一致性。与 DrivingForward 的差异。DrivingForward 从单帧图像预测逐像素深度，本质是 "2.5D" 表示；PointForward 利用预训练 EfficientNetV2 提取图像特征和冻结的 InfiniDepth 预测密集深度，将像素直接提升到 3D 空间形成稀疏查询集 $Q = Q_s \cup Q_d$（$V \times N_s$ 个静态查询 + $O \times N_d$ 个动态查询），避免了像素级预测中的视角冲突和冗余计算。场景图驱动的动态建模。每个 3D 查询点通过可变形注意力机制从多视角图像特征中采样信息，预测对应的高斯属性（位置 $\mu$、旋转 $r$、缩放 $s$、不透明度 $o$、高维颜色特征 $c$）。利用场景图显式组织运动实例——将落入 3D 边界框的查询标记为动态并变换到实例规范空间（以首个边界框为参考），其余标记为静态，实现实例级运动传播和时间一致动态表示。渲染时动态查询从规范空间变换回目标时间戳的世界坐标，高维颜色特征经轻量 UNet 解码为 RGB 图像。不透明度通过时间衰减调整：$o_{t'} = o \cdot e^{-12(t'-t+1)^2}$，确保高斯仅在对应时间窗口内有效。时空融合的跨视图一致性。对每个查询，将其 3D 坐标投影到所有输入图像采样特征 $f_i$，同时计算深度差异 $d_i$ 反映几何一致性。融合权重综合考虑深度可靠性和时间邻近性：$w_i = \sigma(W_d \cdot \mathrm{MLP}([f_{s,i}; d_i]))$，其中动态查询在静态区域采样的特征被置零以避免不一致的运动特征。最终融合特征 $f = \sum_{i=1}^{V} w_i f_{t,i}$ 聚合所有视角信息，实现全局一致的 3D 表示。
\textbf{方法信息：} Xiaomi 参与了这篇论文。
\textbf{设计动机与边界：} 相对原始 3DGS 以独立高斯和图像重建误差为主的设定，PointForward 把问题落在以下可实现的改动上：两个核心改进：(1)点对齐表示：将高斯预测锚定到3D点而非2D像素，消除逐像素预测的多视角不一致和分层伪影；(2)3D场景图：直接在3D空间建模动态实例（车辆、行人等），动态查询合并到规范空间确保跨视角一致性。场景图节点对应实例，边表示空间关系。其设计目的不是单独增加一个模块，而是用上述表示、监督或渲染约束消除原流程在目标任务中的具体失配。提出PointForward，用点对齐表示替代逐像素高斯预测范式，解决多视角不一致和分层伪影问题。通过3D场景图直接在3D空间建模动态实例，避免密集光流预测的局限。适用时还应保留论文给出的边界：对复杂动态场景的场景图构建有挑战。

\section{仿真与数据生成}

仿真方向把 3DGS 从重建结果变成数据机器：它可以闭环渲染自车未来观测、生成 V2X 视角、编辑交通参与者、扩增目标检测数据，或把神经重建伪影交给扩散增强器修复。这里的关键不是单帧 PSNR，而是是否能在控制、评估、传感器和数据增强中稳定使用。
% METHOD: 2410.13571
\paragraph{DriveDreamer4D（世界模型作为 4DGS 数据机）（2024.10）}
DriveDreamer4D~\cite{2410_13571} 用世界模型改进 4D driving scene representation。Novel Trajectory Generation Module 生成结构化交通条件，并把前景和背景运动分开控制；这些条件随车辆操作投影到新视角，保证时空约束。Cousin Data Training Strategy 将真实和合成数据时间对齐后共同训练 4DGS，并用正则保持感知一致。相对 DreamDrive，它更强调命令对齐的新颖轨迹和 synthetic-real 数据混合训练，把视频世界模型当作 4DGS 的数据发动机。
\textbf{设计动机与边界：} 相对原始 3DGS 以独立高斯和图像重建误差为主的设定，DriveDreamer4D 把问题落在以下可实现的改动上：方法分为两个核心模块：(1)新颖轨迹生成模块(NTGM)：生成多样化结构化交通条件，独立调控复杂驾驶环境中前景和背景元素的运动动态，条件经过与车辆操作同步的视图投影确保时空约束一致性；(2)Cousin Data Training Strategy(CDTS)：将时间对齐的真实和合成数据融合训练4DGS，引入正则化损失确保感知一致性。世界模型基于DriveDreamer2，利用扩散模型生成多样化、命令对齐的视频视角。4DGS模型通过时间场模块将规范空间的3D高斯体变形到各时间帧。为了消除原始 3DGS 在该目标任务中的具体失配，方法把上述表示、监督或渲染约束作为可学习和可执行的改动，而不是单独堆叠模块。首次利用视频生成模型改进驾驶场景的4D重建，提出结构化交通条件独立调控前景/背景运动，设计CDTS策略融合真实与合成数据优化4DGS适用时还应保留论文给出的边界：依赖世界模型生成质量；合成数据与真实数据之间可能存在域差距；尚未进行闭环端到端驾驶评估
\textbf{端到端流程：} DriveDreamer4D 的输入是真实驾驶序列、车辆操作条件和世界模型生成的命令对齐视频。首先 NTGM 分别控制前景与背景的运动并用视图投影保持时空一致；随后 CDTS 将时间对齐的真实/合成 cousin data 混合，用感知一致性正则训练 4DGS。时间场把规范空间高斯变形到每一帧，渲染器再按当前相机投影、排序和 alpha 混合输出动态驾驶视图；因此其输出既是可渲染 4D 场，也保留了合成数据域差距这一训练边界。
\textbf{方法信息：} 引用很多，Semantic Scholar 约 123 次。
% METHOD: 2412.01718
\paragraph{HUGSIM（闭环照片级自动驾驶仿真器）（2024.12）}
HUGSIM~\cite{2412_01718} 是基于 3DGS 的闭环实时自动驾驶仿真器。它从采集 RGB 图像提升到 3D 高斯空间，并解决闭环中的三类挑战：自车偏离原轨迹的外推视点，演员车辆的 360 度渲染，以及离散观测动态物体到连续轨迹的生成。仿真循环根据控制命令更新自车和 actor 状态，再渲染新的传感器观测。相对重建方法，HUGSIM 的锚点是闭环：规划和控制会改变未来视角，渲染系统必须能响应。
\textbf{实现细节：} HUGSIM（2024.12） 是首个基于 3DGS 的完整闭环驾驶仿真平台。其核心突破在于将 2D RGB 图像重建为可交互的 3D 高斯场景，并动态更新自车和其他参与者的状态与观测。直觉上，这相当于把 2D 视频游戏升级为 3D 可交互世界——不仅能看到场景，还能在其中自由驾驶并观察所有参与者的实时响应。Baseline 与问题。闭环仿真要求仿真器根据自动驾驶系统的控制指令（转向、加速、制动）实时更新场景状态，生成新的传感器观测，形成完整的感知-规划-控制循环。传统 NeRF 仿真器渲染速度太慢（秒级/帧），无法支持实时闭环测试。CARLA 等游戏引擎仿真器虽然实时，但网格资产的建模成本高昂，渲染保真度与真实传感器数据存在显著 sim-to-real 差距。亟需一种兼具照片级保真度和实时渲染速度的场景表示。改了什么。HUGSIM 将场景表示从 NeRF 的隐式辐射场或游戏引擎的网格资产，替换为 3DGS 的显式高斯场景图——静态背景、动态物体和天空分别建模，统一泼溅渲染。场景表示与渲染。HUGSIM 将城市场景建模为包含静态背景、动态物体和天空的 3D 高斯场景图。静态背景用标准 3D 高斯表示，动态车辆在物体局部坐标系中建模后通过位姿变换到世界坐标，天空用环境贴图表示。利用 3DGS 的实时渲染能力实现闭环仿真中的连续帧生成，同时解决视角外推（自车偏离训练轨迹时）和 360 度车辆渲染（环绕视角观测动态车辆）两大挑战。设计动机。3DGS 的显式点云表示使场景编辑天然可行——动态物体可独立编辑（移动、删除、插入）而不影响背景，且泼溅渲染管线对场景修改实时响应。NeRF 的隐式表示无法实现这一点：修改场景需要重新训练整个辐射场。同时，高斯基元的世界坐标与物理引擎兼容，可直接从高斯位置提取碰撞检测所需的几何信息。性能与基准。
\textbf{实现补充：} 实现上还要核对输入、输出、参数、特征、初始化、训练损失、正则、渲染或推理接口，以及优化中的剪枝、分裂、采样、编码或解码步骤；这些抓手决定该方法到底改了 3DGS 管线的哪一环。
\textbf{方法信息：} Huawei 参与了这篇论文。



% METHOD: 2501.00601
\paragraph{DreamDrive（2D-3D-4D 渐进生成驾驶场景）（2025.01）}
DreamDrive~\cite{2501_00601} 从街景图像生成可控 4D driving scenes。它先利用视频扩散模型生成 2D 视觉参考序列和潜特征，再通过高斯初始化提升到 3D；随后用扩散早期层潜特征捕获动态，用聚类在 4D 时空中分解静态和动态区域。静态结构用 3DGS，动态物体用 deformable 4DGS。相对 MagicDrive3D 的多视角视频到 3DGS，DreamDrive 更强调从野外驾驶数据自监督生成 4D 场景。
\textbf{实现细节：} DreamDrive（2025.01） 修改了 3DGS 的场景表示和初始化管线——用混合高斯表示（静态时间无关 + 动态时间相关）替代单一时间相关高斯，高斯泼溅渲染函数 $F_\mathrm{splat}$ 保持不变。核心思路是将视频扩散模型的生成能力与 3DGS 的几何一致性渲染相结合，沿着 2D $\to$ 3D $\to$ 4D 的渐进路径构建驾驶场景。Baseline 与问题。标准 3DGS 场景重建依赖精心标注的驾驶日志（标定相机位姿、物体边界框、3D 点云），限制了在无标注野外数据上的扩展性。另一条路线——纯视频扩散模型——虽然可以从任意驾驶视频生成动态视频，但因缺乏 3D 几何约束，跨帧几何一致性差，生成的视频无法可靠用于自动驾驶下游任务。Deformable GS 等方法将全部高斯建模为时间相关，对生成图像中固有的 3D 不一致性过拟合，导致静态背景在新视角下出现"虚假变形"（fake dynamics）。2D $\to$ 3D：视频扩散先验与高斯初始化。DreamDrive 的第一个关键步骤是用视频扩散模型（如 Vista、MagicDrive）从输入控制信号（单张图像 $I_\mathrm{ctrl}$ 或地图 $M_\mathrm{ctrl}$）生成 2D 参考图像序列 $I_\mathrm{ref}$ 和潜特征 $Z_\mathrm{ref}$。
\textbf{方法信息：} NVIDIA 参与了这篇论文。
\textbf{设计动机与边界：} 相对原始 3DGS 以独立高斯和图像重建误差为主的设定，DreamDrive 把问题落在以下可实现的改动上：DreamDrive遵循2D-3D-4D渐进生成流程：1）视频扩散先验：利用视频扩散模型生成2D视觉参考序列和潜在特征；2）3D提升：通过高斯初始化将2D参考提升到3D空间；3）自监督场景分解：利用扩散模型早期层的潜在特征捕获视觉动态，通过聚类分组策略在4D时空中分解静态和动态区域；4）混合高斯表示：用不同的高斯表示建模静态结构（静态3DGS）和动态物体（可变形4DGS），支持从任意驾驶轨迹渲染3D一致的视频。方法不依赖昂贵的物体标注，可直接从野外驾驶数据学习。其设计目的不是单独增加一个模块，而是用上述表示、监督或渲染约束消除原流程在目标任务中的具体失配。提出结合视频扩散生成能力和3DGS几何一致性渲染的4D场景生成方法，引入自监督混合高斯表示和聚类分组策略分解静态/动态区域，实现从野外驾驶数据生成可控、泛化的4D驾驶场景。适用时还应保留论文给出的边界：视频扩散模型的3D不一致性仍可能影响最终质量；对复杂动态场景的分解可能不够完美。
% METHOD: 2503.11731
\paragraph{Industrial-Grade Sensor Simulation（模块化工业级 GS 仿真）（2025.03）}
Industrial-Grade Sensor Simulation via Gaussian Splatting~\cite{2503_11731} 从系统角度重构自动驾驶传感器仿真。它选择 2D neural Gaussians 进行物理一致场景和传感器建模，支持相机与 LiDAR；显式高斯基元支持拖放式场景编辑和资产库增强；可控扩散模型负责场景扩展与协调。相对 HUGSIM 的闭环仿真器，这篇更强调工业流水线：几何、光度、场景编辑、扩散补全和全栈验证模块化集成。
\textbf{实现细节：} 多个工作从工业应用和特殊角度推进仿真技术。Nouri et al.（2026.05） 系统分析 3DGS 重建在安全相关场景中的保真度退化。Tan et al.（2026.05） 提出正交投影梯度（OPG）解决闭环仿真中空间-时间参数的信用分配困境。3DArticCyclists（2024.10） 使用 3DGS 生成可控关节自行车骑行者数据。RenderWorld（2024.09） 用自监督高斯 Img2Occ 模块生成 3D 占用标签，实现仅视觉端到端驾驶。Gaussian Splatting as Data Generator（2025.04） 证明 3DGS 重建场景可直接用作 3D 检测的高效数据增强源。Zeng et al.（2025.03） 提出工业级传感器仿真模块化框架。NeRF2GS（2025.03） 利用 NeRF 的泛化能力蒸馏到 3DGS 实现混合渲染。Sensor2Sensor（2026.05） 通过 4DGS 重建实现行车记录仪到多模态传感器套件的生成转换。ReCamDriving（2025.12） 提出纯视觉相机控制新轨迹视频生成框架，利用 3DGS 渲染提供密集几何引导实现精确相机可控生成，构建 ParaDrive 数据集包含超过 110K 并行轨迹视频对。
\textbf{技术实现补充：} 多个工作从工业应用和特殊角度推进仿真技术。Nouri et al.（2026.05） 系统分析 3DGS 重建在安全相关场景中的保真度退化。Tan et al.（2026.05） 提出正交投影梯度（OPG）解决闭环仿真中空间-时间参数的信用分配困境。3DArticCyclists（2024.10） 使用 3DGS 生成可控关节自行车骑行者数据。RenderWorld（2024.09） 用自监督高斯 Img2Occ 模块生成 3D 占用标签，实现仅视觉端到端驾驶。Gaussian Splatting as Data Generator（2025.04） 证明 3DGS 重建场景可直接用作 3D 检测的高效数据增强源。Zeng et al.（2025.03） 提出工业级传感器仿真模块化框架。NeRF2GS（2025.03） 利用 NeRF 的泛化能力蒸馏到 3DGS 实现混合渲染。Sensor2Sensor（2026.05） 通过 4DGS 重建实现行车记录仪到多模态传感器套件的生成转换。ReCamDriving（2025.12） 提出纯视觉相机控制新轨迹视频生成框架，利用 3DGS 渲染提供密集几何引导实现精确相机可控生成，构建 ParaDrive 数据集包含超过 110K 并行轨迹视频对。 %===================================================================== 近期方法脉络：可编辑仿真与Sim-to-Real、车辆资产生成。这一组近期工作延续 editable 3DGS driving simulation、feed-forward 3D Gaussians，主要把改动集中在增量改进、应用等环节。下面只保留相对既有路线的变化，便于和前文方法直接对照。
\textbf{方法信息：} Alibaba 参与了这篇论文。
\textbf{设计动机与边界：} 相对原始 3DGS 以独立高斯和图像重建误差为主的设定，Industrial-Grade Sensor Simulation via Gaussian Splatting 把问题落在以下可实现的改动上：框架通过GS重构三个关键组件：(1) 选择2D神经高斯表示（2DGS）进行物理一致的场景和传感器建模，支持相机和LiDAR传感器仿真；(2) 场景编辑管道，利用高斯基元的显式特性支持拖放操作，结合高斯基元库（如3DRealCar）进行数据增强；(3) 耦合可控扩散模型进行场景扩展和协调，补充不完整的重建场景。通过消融研究展示了GS相比NeRF在帧级仿真延迟、几何和光度一致性、以及可解释显式编辑方面的优势。其设计目的不是单独增加一个模块，而是用上述表示、监督或渲染约束消除原流程在目标任务中的具体失配。首个工业级GS传感器仿真系统，通过模块化重构三大核心组件（2D神经高斯、场景编辑管道、扩散模型场景扩展），实现端到端自动驾驶算法的全栈测试。适用时还应保留论文给出的边界：专有数据集限制可复现性；扩散模型场景扩展的质量控制有挑战；多传感器标定精度要求高。
% METHOD: 2504.16740
\paragraph{3DGS Data Generator（目标检测数据增强）（2025.04）}
Gaussian Splatting is an Effective Data Generator for 3D Object Detection~\cite{2504_16740} 基于 OmniRe 重建的实例感知 3DGS，在三维空间中直接放置外部车辆或物体资产。这样生成图像同时拥有精确 3D 姿态标注和物理合理位置。实验发现对 3D 检测而言，几何多样性比外观多样性更重要，且这种 3DGS 数据增强优于只做 2D 扩散生成。相对 CRUISE 的 V2X 场景编辑，它更聚焦相机 3D 目标检测训练样本。
\textbf{实现细节：} 多个工作从工业应用和特殊角度推进仿真技术。Nouri et al.（2026.05） 系统分析 3DGS 重建在安全相关场景中的保真度退化。Tan et al.（2026.05） 提出正交投影梯度（OPG）解决闭环仿真中空间-时间参数的信用分配困境。3DArticCyclists（2024.10） 使用 3DGS 生成可控关节自行车骑行者数据。RenderWorld（2024.09） 用自监督高斯 Img2Occ 模块生成 3D 占用标签，实现仅视觉端到端驾驶。Gaussian Splatting as Data Generator（2025.04） 证明 3DGS 重建场景可直接用作 3D 检测的高效数据增强源。Zeng et al.（2025.03） 提出工业级传感器仿真模块化框架。NeRF2GS（2025.03） 利用 NeRF 的泛化能力蒸馏到 3DGS 实现混合渲染。Sensor2Sensor（2026.05） 通过 4DGS 重建实现行车记录仪到多模态传感器套件的生成转换。ReCamDriving（2025.12） 提出纯视觉相机控制新轨迹视频生成框架，利用 3DGS 渲染提供密集几何引导实现精确相机可控生成，构建 ParaDrive 数据集包含超过 110K 并行轨迹视频对。
\textbf{方法信息：} Qualcomm 参与了这篇论文。
\textbf{设计动机与边界：} 相对原始 3DGS 以独立高斯和图像重建误差为主的设定，Gaussian Splatting is an Effective Data Generator for 3D Object Detection 把问题落在以下可实现的改动上：方法基于OmniRe构建场景的3D高斯表示，将场景分解为静态节点(背景)、刚性节点(前景车辆)和可变形节点(行人)。通过实例感知的3DGS重建，在3D空间中直接应用几何变换放置新的3D物体资产，确保物理合理的位置和精确的3D姿态标注。渲染过程使用alpha混合聚合深度排序的高斯基元。方法还探索了困难样本生成策略（最大化检测损失或施加高视觉遮挡），发现这些策略对基于相机的3D检测效果有限。其设计目的不是单独增加一个模块，而是用上述表示、监督或渲染约束消除原流程在目标任务中的具体失配。首次提出利用3DGS（OmniRe）重建驾驶场景作为数据增强工具，通过在重建的3D空间中直接放置外部3D物体，确保物理合理性和精确的3D标注。实验证明几何多样性比外观多样性更重要，且该方法优于基于扩散的方法。适用时还应保留论文给出的边界：外观多样性受限于3D资产库大小；场景重建质量依赖OmniRe的性能；对动态场景的处理可能有限。
% METHOD: 2506.04218
\paragraph{Pseudo-Simulation（预生成观测评估规划）（2025.06）}
Pseudo-Simulation for Autonomous Driving~\cite{2506_04218} 提出介于开环和闭环之间的评估范式。系统先在真实数据上训练 3DGS，再围绕每个评估位置预生成不同位置、朝向和速度的合成观测。规划器预测轨迹后，评估用基于距离的权重选择接近该轨迹的合成观测，从而近似未来闭环状态。相对 HUGSIM 真正运行闭环，Pseudo-Simulation 用预生成高斯观测换取可扩展性，同时比纯开环更接近因果评估。
\textbf{实现细节：} 仿真范式方面，REAP 提出基于非对称混合强化学习的端到端自动泊车方法，actor 从环视图像学习策略、critic 使用仿真器提供的无伪影 BEV 特权信息评估价值函数；ParkingWorld 提出纠错回路样本高效强化学习框架，多级回放缓冲区显式配对"策略失败片段"与"人工纠正片段"提升样本效率——两者将 3DGS 仿真器与 RL 训练循环深度耦合。SceneCrafter 和 Pseudo-Simulation 走"伪仿真"路线，用 3DGS 预生成多样化合成观测建立公共排行榜，Pseudo-Simulation 进一步定量验证伪仿真与闭环仿真的高相关性。
\textbf{实现补充：} 实现上还要核对输入、输出、参数、特征、初始化、训练损失、正则、渲染或推理接口，以及优化中的剪枝、分裂、采样、编码或解码步骤；这些抓手决定该方法到底改了 3DGS 管线的哪一环。
\textbf{方法信息：} NVIDIA 参与了这篇论文。
\textbf{设计动机与边界：} 相对原始 3DGS 以独立高斯和图像重建误差为主的设定，Pseudo-Simulation for Autonomous Driving 把问题落在以下可实现的改动上：核心思想是用3DGS近似AV可能遇到的未来状态。首先在真实数据集上训练3DGS模型，然后在每个评估位置周围预生成大量合成观测（变化位置、朝向、速度）。评估时，将AV预测轨迹与合成观测对比，使用基于距离的加权方案：越接近AV预测位置的合成观测权重越高。这样可以在不进行序列交互仿真的情况下，评估误差恢复和因果混淆缓解能力。支持nuscenes等数据集。其设计目的不是单独增加一个模块，而是用上述表示、监督或渲染约束消除原流程在目标任务中的具体失配。提出伪仿真评估范式，在真实数据集上使用3DGS预生成多位置/朝向/速度的合成观测，通过基于距离的加权方案对AV可能行为进行评估。与闭环仿真的相关性R\^{}2=0.8，优于最佳开环方法的R\^{}2=0.7，同时建立了公开排行榜。适用时还应保留论文给出的边界：合成观测的多样性受限于3DGS渲染质量；无法完全模拟交互式场景变化；生成合成观测需要预计算开销。
% METHOD: 2507.18473
\paragraph{CRUISE（V2X 协同重建与编辑）（2025.07）}

CRUISE~\cite{2507_18473} 面向 V2X 场景。它采用分解式 Gaussian Splatting 重建真实场景，将动态交通参与者拆成可编辑高斯表示，支持从自车和基础设施视角渲染。框架分为数据、重建、编辑、渲染和应用五层，既能替换动态 actor，也能增强 V2X 数据集。相对单车视角 DrivingGaussian，CRUISE 的增量是协同视角和基础设施相机，重建结果直接服务 V2X 感知训练。
\textbf{实现细节：} CRUISE（2025.07） 与 MultiEditor 的单车单模态编辑不同，CRUISE 面向车路协同（V2X）场景——同时从自车和基础设施视角重建、编辑和渲染，生成大规模 V2X 数据增强训练集。V2X 专用重建适配。直接将 Street Gaussians 应用于 V2X 数据会遇到独特问题：基础设施视角的标注包含自车边界框，在自车视角训练时该框覆盖大部分可见场景，导致静态背景被错误重建为动态高斯。CRUISE 移除自车框并引入自车掩码遮蔽自车区域，实现动静正确分离。多帧 LiDAR 点云融合（自车 + 基础设施）提供稳健初始化，跟踪框聚合多帧点云增强物体高斯密度。训练损失包含 8 项——颜色、深度、法向、天空、语义、尺度、比率比和浮动正则化——联合约束几何和外观。GPT 驱动的交通流编辑。收集互联网多视角车辆图像和尺寸信息，通过 TRELLIS 生成高质量 3D 高斯车辆资产。编辑时简化向量地图（BFS 扩展交叉口相关车道），将车道中心线的起止位置和方向提供给 GPT-4o，结合自车轨迹采样生成合理的多车插入轨迹。生成的高斯车辆资产按轨迹定位后渲染来自两个视角的同步图像和 3D 标注，构建完全标注的合成 V2X 数据集。
\textbf{设计动机与边界：} 相对原始 3DGS 以独立高斯和图像重建误差为主的设定，CRUISE 把问题落在以下可实现的改动上：采用分层结构的分解式高斯泼溅表示，将动态交通参与者分解为可编辑的高斯表示，支持无缝修改和增强驾驶场景。框架分为五个层次：数据层将原始传感器数据转换为统一格式；重建层用分解式GS重建场景；编辑层支持动态参与者替换和场景修改；渲染层从自车和基础设施双视角渲染图像；应用层用于训练和评估V2X系统。其设计目的不是单独增加一个模块，而是用上述表示、监督或渲染约束消除原流程在目标任务中的具体失配。提出首个面向V2X驾驶环境的完整重建-合成框架CRUISE，采用分解式高斯泼溅重建真实场景并支持灵活编辑，可从自车和基础设施视角渲染图像，实现大规模V2X数据集增强。适用时还应保留论文给出的边界：V2X通信延迟和带宽对实时渲染的影响；极端天气条件下的重建质量；仿真数据与真实数据的域差距。
\textbf{方法信息：} Baidu 参与了这篇论文；代码：\url{https://github.com/sainingzhang/cruise}。



% METHOD: 2508.18788
\paragraph{PseudoMapTrainer（无 HD 地图的伪标签训练）（2025.08）}
PseudoMapTrainer~\cite{2508_18788} 用高斯 surfels 从多视角相机时序重建道路表面，并把几何、颜色和语义属性渲染成 BEV 语义地图，再导出矢量化伪地图标签。训练在线地图模型时，引入 mask-aware assignment 和损失处理遮挡和缺失伪标签。相对 ParkGaussian 强化停车位区域重建，PseudoMapTrainer 把高斯道路表面变成地图学习标签生成器。
\textbf{实现细节：} PseudoMapTrainer（2025.08） 利用高斯泼溅从多相机图像重建道路表面，结合预训练 2D 分割网络的语义生成伪标签，首次实现无需任何真值高精地图的在线建图模型训练。与 MapGS 用 3DGS 渲染新传感器配置图像的思路不同，PseudoMapTrainer 利用 3DGS 的核心能力在于"精确 3D 表面重建"——从多相机图像重建场景后，沿道路表面均匀采样 3D 点并投影到各相机视角，用预训练 2D 语义分割网络预测各投影点的语义标签，再将多视角语义投票聚合回 3D 空间生成伪高精地图标注。这种范式将 3DGS 作为几何桥梁，将 2D 语义信息提升到 3D 鸟瞰图空间，绕过了高精地图真值标注的成本瓶颈。
\textbf{方法信息：} 代码：\url{https://github.com/boschresearch/pseudomaptrainer}。
\textbf{设计动机与边界：} 相对原始 3DGS 以独立高斯和图像重建误差为主的设定，PseudoMapTrainer 把问题落在以下可实现的改动上：方法包含三个核心组件：(1)伪标签生成：使用高斯面元(Gaussian surfels)从时序多视角相机图像重建道路表面为网格，面元编码几何、颜色和语义属性（来自预训练2D分割网络），直接渲染语义BEV地图并导出矢量化伪地图标签；(2)掩码感知训练：引入掩码感知分配算法和损失函数，处理伪标签中不可避免的遮挡和数据缺失问题，实现从不完整标签鲁棒学习；(3)半监督学习：在大规模无标注众包数据上预训练，再在有限的真值标签上微调。其设计目的不是单独增加一个模块，而是用上述表示、监督或渲染约束消除原流程在目标任务中的具体失配。首次提出无需任何真值HD地图即可训练在线地图模型的方法。利用高斯泼溅重建道路表面生成伪标签：将道路表面建模为高斯面元(surfels)网格，结合语义分割网络渲染语义BEV地图并导出矢量化伪地图标签。引入掩码感知分配算法和损失函数处理部分遮挡的伪标签。支持半监督预训练利用大规模无标注众包数据。适用时还应保留论文给出的边界：伪标签质量受限于高斯泼溅重建质量；对遮挡严重的区域伪标签不完整；众包数据的多样性和质量控制是挑战。
% METHOD: 2602.24096
\paragraph{DiffusionHarmonizer（在线单步扩散增强仿真）（2026.02）}
DiffusionHarmonizer~\cite{2602_24096} 不是重建新高斯，而是把 3DGS 或 NeRF 渲染的伪影帧实时增强为高保真输出。它把多步扩散改成确定性单步增强器，直接输入干净 latent，固定时间步和空文本条件，并加入短历史时间上下文保持时序稳定。数据策展管线构建合成-真实配对，解决外观协调、伪影校正和光照真实感。相对 FaithFusion 把扩散结果写回 3DGS，DiffusionHarmonizer 更像仿真渲染后的在线后处理器。
\textbf{实现细节：} DiffusionHarmonizer（2026.02） 是首个将预训练图像扩散模型转化为确定性单步在线增强器的工作，用于在仿真循环中实时修复 3DGS 重建渲染的缺陷。核心创新——单步确定性增强。标准扩散模型需要多步去噪（通常 20--50 步），无法满足在线仿真的实时需求。DiffusionHarmonizer 将预训练扩散模型改造为确定性单步增强器：直接将清洁潜变量 $E_\theta(I_t)$ 输入网络（不注入噪声），时间步和文本条件标记固定为常量"空"值，产生从输入潜变量到增强潜变量的稳定确定性映射。为保持时间一致性，骨干网络 $F_\theta$ 扩展了时间注意力层，以当前帧和 $K=4$ 帧历史增强帧作为上下文输入。定制数据管线。由于缺乏公开的合成-真实配对数据，设计了五个互补的数据合成组件：（1）新视角伪影校正——通过稀疏重建、循环重建和模型欠拟合生成降质渲染；（2）ISP 修改——随机化色调映射、曝光、白平衡等参数模拟前景-背景外观不匹配；（3）重光照——用扩散重光照模型在随机光照条件下重新生成选定区域；（4）基于物理的阴影仿真——用物理渲染器在可控光照下合成精确阴影；（5）资产重新插入——将动态前景重新插入 3DGUT 重建的背景中但不生成正确阴影。
\textbf{方法信息：} NVIDIA 参与了这篇论文。
\textbf{设计动机与边界：} 相对原始 3DGS 以独立高斯和图像重建误差为主的设定，DiffusionHarmonizer 把问题落在以下可实现的改动上：方法包含两个核心组件：(1) 在线帧到帧增强器：将预训练扩散模型改编为确定性单步增强器，直接输入干净潜码而非加噪潜码，固定时间步和文本条件为常量null值。引入时间上下文（K=4帧历史）通过时间注意力层保持时间稳定性；(2) 数据策展管线：构建合成-真实配对数据集，覆盖外观和谐化、伪影校正和光照真实感等多因素。训练时采用噪声-轨迹匹配策略解决多步预训练与单步推理间的差异。其设计目的不是单独增加一个模块，而是用上述表示、监督或渲染约束消除原流程在目标任务中的具体失配。提出首个面向在线仿真的扩散增强框架，核心是将多步扩散模型转换为单步时间条件增强器，可在单GPU上实时运行。关键创新在于解决噪声-轨迹不匹配问题（多步预训练与单步推理间的差异），以及设计可扩展的数据策展管线构建合成-真实配对数据。适用时还应保留论文给出的边界：时间上下文长度有限可能导致长序列漂移；对极端光照变化的处理能力受限；依赖高质量合成-真实配对数据。
% METHOD: 2604.18468
\paragraph{Asset Harvester（驾驶日志到可仿真 3D 资产）（2026.04）\cite{2604_18468}}
Asset Harvester 补的是自动驾驶仿真中一个与整场景重建不同的缺口。DrivingGaussian、Street Gaussians、HUGSIM 或 NuRec 这类系统可以把一段驾驶日志重建成可渲染场景，但闭环仿真还需要可复用、可操控、可插入的新资产：车辆、行人、骑行者或路边物体应该能脱离原始日志，在别的场景、别的视角和别的交通配置中使用。直接从单张车载图像做 image-to-3D 往往看不见背面和遮挡区域；直接使用日志中的稀疏多视角又会遇到视角夹角小、目标低分辨率、夜间和模糊等真实 AV 数据问题。Asset Harvester 的改变，是把真实驾驶日志整理成对象级稀疏观测，再用多视图生成和 3D Gaussian lifting 组合成 simulation-ready assets。

具体流程可以分成三层。第一层是数据和几何预处理：系统从 AV logs 中检测并裁剪对象，利用相机标定、目标 mask、对象位姿和跨帧观测构造 object-centric training tuples；这些样本不是干净的电商物体图，而是带遮挡、运动、有限视角和传感器差异的真实驾驶对象。第二层是 SparseViewDiT，它以一个或多个稀疏视角为条件生成更完整的多视角外观，专门处理 limited-angle views；训练中还结合混合数据策展、增强和 self-distillation，让模型学会从少量真实观测补全对象背面和未见区域。第三层是 3D Gaussian lifting：生成出的多视图对象不只是作为 2D 图像结果保存，而是被提升成可渲染的高斯资产，使该对象可以插回 NuRec/CARLA 式仿真环境，并与 DiffusionHarmonizer 等后处理或协调模块配合减少插入伪影。

因此它与 MagicDrive3D、DreamDrive 或 DriveDreamer4D 的关系不是“再生成一个整段驾驶视频”，而是从日志中收割可复用 object assets；与 3DGS Data Generator 的关系也不同，后者强调把已有外部资产放进 3DGS 场景增强检测数据，Asset Harvester 则负责从真实驾驶数据中生成资产库本身。NVIDIA 官方仓库 \url{https://github.com/NVIDIA/asset-harvester} 已公开代码，约 219 stars，并提供 Hugging Face 模型/演示链接；这使它在可信书中应作为自动驾驶仿真资产生成路线的重要方法，而不是被生成章节或普通前馈重建章节吞掉。
\textbf{实现细节：} DENSER（2024.09） 引入小波动态估计 SH 基增强远距离动态物体外观，跨多帧致密化点云加速收敛。VDG（2024.06） 首次集成自监督视觉里程计实现仅 RGB 输入的无位姿动态高斯。ArmGS（2025.07） 通过多粒度外观精化（局部高斯级、全局图像级、动态角色级）提升帧间和相机间精细变化建模。2D 先验与 SDF 引导（2025.10） 结合 SDF 与 3DGS 实现无 LiDAR 城市场景渲染和场景编辑。SRIF（2024.09） 利用扩散图像变形和流估计实现语义形状注册。
\textbf{实现补充：} 进一步拆开来看，输入侧对应图像、点云、位姿、语义提示或已训练高斯资产；输出侧对应新的高斯参数、渲染图像、编辑结果、地图状态或评测指标。表示层需要说明位置、尺度、旋转、不透明度、特征、材质或 MLP 解码器怎样改；训练层需要说明损失、正则、采样、剪枝、分裂或压缩何时发生；推理层则看渲染、查询、规划或系统接口如何接到原始 3DGS / Scaffold-GS / 4DGS 管线上。
\textbf{方法信息：} NVIDIA 提出并开源；代码：\url{https://github.com/NVIDIA/asset-harvester}；GitHub 星标较高，约 219 stars；Hugging Face 模型页：\url{https://huggingface.co/nvidia/asset-harvester}。
\section{传感器融合与定位}

相机 3DGS 只需渲染 RGB，而自动驾驶仿真还需要 LiDAR 点云、强度、ray-drop、雷达回波和复杂相机模型。传感器融合方法的增量主要是改渲染器：LiDAR 常用 range-view、球面或 ray-splat 交叉；雷达需要天线增益、多径和噪声；鱼眼和滚动快门则要求相机模型本身进入可微渲染。


% METHOD: 2410.05111
\paragraph{LiDAR-GS（range-view 激光束 splatting）（2024.10）}
LiDAR-GS~\cite{2410_05111} 设计可微 laser beam splatting。LiDAR 与相机不同，点应组织在 range-view 中，高斯应投影到与光束对齐的微平面，而不是标准 NDC 空间。方法用 neural Gaussian representation 学习 LiDAR 强度与入射方向、距离的关系，并为每个高斯学习 ray-drop 概率，同时渲染深度、强度和 ray-drop。相对 GS-LiDAR 的 2D 周期高斯，LiDAR-GS 更强调激光束几何和可微重仿真。
\textbf{实现细节：} LiDAR-GS（2024.10） 是首个将 3DGS 扩展到 LiDAR 点云仿真的工作，核心创新是可微分激光束泼溅（differentiable laser beam splatting）。直觉上，这将相机的"拍照"模式转换为激光的"测距"模式——原本投影到 2D 照片的操作变为沿激光射线计算精确距离，每个激光脉冲击中高斯椭球后返回的距离值构成距离图。Baseline 与问题。3DGS 的渲染管线基于针孔相机的局部仿射投影——将 3D 高斯投影到 2D 图像平面进行 $\alpha$-混合。但 LiDAR 的成像原理完全不同：主动发射激光脉冲，通过飞行时间（ToF）测量距离，生成距离图（range image）而非 RGB 图像。直接将 3DGS 的针孔投影用于 LiDAR 会产生局部仿射近似伪影，因为激光射线的锥形展开与针孔相机的线性投影在数学上不同。核心改动——可微分激光束泼溅。LiDAR-GS 使用距离图表示，将激光射线投射到 3D 高斯的微截面（micro-facet）上实现精确表面泼溅。每个高斯基元除了标准的位置、协方差、不透明度外，还附加神经高斯表示——一个小型 MLP 集成视角依赖线索来建模 LiDAR 的关键属性：激光入射方向和外部因素（天气、表面材质）对回波的影响。三通道渲染。
\textbf{方法信息：} Alibaba 参与了这篇论文。
\textbf{设计动机与边界：} 相对原始 3DGS 以独立高斯和图像重建误差为主的设定，LiDAR-GS 把问题落在以下可实现的改动上：LiDAR-GS针对LiDAR与相机的四个关键差异进行设计：(1)采用距离视图（range-view）表示组织激光束，将3D高斯投影到与光束对齐的微平面而非统一NDC空间，消除仿射近似误差；(2)提出神经高斯表示（Neural Gaussian Representation）替代球谐函数，通过MLP学习LiDAR强度和入射方向/距离的关系；(3)为每个高斯基元分配可学习的ray-drop概率，通过前向渲染判断激光束有效性；(4)引入空间尺度正则化和紧凑包围盒策略约束高斯尺度，结合动态实例分解机制处理动态场景。整体流程：先通过COLMAP初始化点云，然后联合优化静态背景和动态实例的高斯参数。其设计目的不是单独增加一个模块，而是用上述表示、监督或渲染约束消除原流程在目标任务中的具体失配。首次提出基于3DGS的LiDAR实时重仿真方法LiDAR-GS，设计了可微分激光束溅射方法消除局部仿射近似伪影，引入神经高斯表示替代球谐函数来建模受入射方向影响的LiDAR属性，同时支持深度、强度和ray-drop通道的联合重仿真适用时还应保留论文给出的边界：论文未明确提及具体局限性，但方法依赖于动态实例的分解标注，且神经高斯表示增加了额外的MLP计算开销
% METHOD: 2411.16816
\paragraph{SplatAD（统一相机与 LiDAR 渲染锚点）（2024.11）}
SplatAD~\cite{2411_16816} 是多传感器 3DGS 仿真的代表性锚点。Street Gaussians、PVG 和 OmniRe 主要把 3DGS 用作相机渲染器，LiDAR 只作为初始化或深度监督；但自动驾驶仿真真正要重放的是传感器套件，包括相机、LiDAR 强度、ray drop 和滚动快门。SplatAD 的改变是让同一组 3D Gaussians 同时服务于相机和 LiDAR 渲染，并为两种传感器分别设计投影、tiling、rolling-shutter 校正和解码头。

场景表示仍基于 3DGS，但不用固定球谐函数作为唯一外观参数，而是给每个高斯学习基色和特征向量，并为不同传感器学习 embedding。动态部分沿用场景图分解：静态背景和各 actor 分开建模，actor 高斯定义在物体局部坐标中，渲染时按 SE(3) 位姿变换到世界坐标；位姿和速度还带可学习偏移以修正跟踪误差。相机渲染保留 tile-based splatting，但把滚动快门近似到 2D 图像空间：根据相机和 actor 的相对速度移动高斯投影位置，并扩大 tile 覆盖范围。LiDAR 渲染则把高斯投影到球坐标中的 azimuth/elevation/range，按真实扫描线和非均匀 diode 排布分块；每条 LiDAR ray 渲染 range、intensity 和 ray-drop probability，并用 line-of-sight 损失惩罚真值量测前方的错误不透明度。公开仓库主要开放了基于 gsplat 扩展的相机/激光雷达 rasterization、rolling-shutter 示例和测试，完整模型与数据加载接口说明会转到 neurad-studio 维护。相对 NeuRAD 一类 NeRF 多传感器仿真，SplatAD 的优势是 rasterization 速度；相对相机 3DGS，它把主动传感器的成像物理纳入了高斯渲染。

\textbf{锚点细化：} SplatAD：针对相机和 LiDAR 成像物理过程不同（相机被动接收光线、LiDAR 主动发射激光）、需同时精确建模两种传感器特定现象（滚动快门、LiDAR 强度和光线丢弃）的局限，用\textbf{统一的可微泼溅管线}替代分离的相机/LiDAR 仿真器，核心思想是\textbf{将场景分解为静态背景、天空环境和动态物体（含刚性和可变形物体）统一建模}——动态物体通过 4D 傅里叶球谐建模时变外观，滚动快门通过逐行独立的相机外参建模，LiDAR 模拟通过可微泼溅同时渲染深度、强度和光线丢弃，渲染速度比 NeRF 方法快一个数量级，是首个同时实现动态场景相机和 LiDAR 实时真实感渲染的 3DGS 方法。
\textbf{方法信息：} Zenseact 参与了这篇论文；代码已公开：\url{https://github.com/carlinds/splatad}，仓库约 404 stars，不过 README 说明完整 SplatAD 模型和数据加载等组件会通过 \url{https://github.com/georghess/neurad-studio} 发布；Semantic Scholar 约 70 次引用。



% METHOD: 2501.13971
\paragraph{GS-LiDAR（全景 LiDAR 高斯）（2025.01）}

GS-LiDAR~\cite{2501_13971} 专门生成真实 LiDAR 点云。它使用带周期振动属性的 2D Gaussians，更适合道路和薄表面；全景渲染通过显式 ray-splat intersection 适配 LiDAR 的柱面或球面投影。每个高斯还含强度和 ray-drop SH 系数，模拟材质反射和激光束丢失。该方法提供代码，GitHub 关注度较高。相对 SplatAD 的统一多传感器框架，GS-LiDAR 更深入地为 LiDAR 几何、强度和丢点设计高斯表示。
\textbf{实现细节：} 与 LiDAR-GS 使用 3D 高斯基元不同，GS-LiDAR（2025.01） 提出基于 2D 高斯基元的全景 LiDAR 仿真框架，核心创新在于周期振动 2D 高斯表示和显式射线-泼溅求交的统一渲染。直觉上，2D 高斯圆盘如同无数个跳动的圆形薄膜拼接成运动物体，每个薄膜独立振动，接力完成复杂运动轨迹的表示。周期振动 2D 高斯。驾驶场景中车辆和行人持续运动，GS-LiDAR 使用 2D 高斯圆盘（而非 3D 椭球）作为基元，每个圆盘定义在局部切平面空间中（中心点、两个主切向量 $\mathbf{t}_u, \mathbf{t}_v$、缩放向量 $(s_u, s_v)$）。关键创新是为每个 2D 高斯引入周期振动属性——振动方向 $\mathbf{d}$、生命峰值 $p$ 和时间衰减率 $\gamma$： \begin{equation*} \boldsymbol{\mu}(t) = \boldsymbol{\mu} + \mathbf{d} \sin(2\pi(t - p)/l), \quad o(t) = o \cdot \exp(-\frac{(t - p)^2}{2\gamma^2}) \end{equation*} 每个基元执行简单周期运动，多个基元以接力方式组合即可表示动态元素的复杂运动。显式射线-泼溅求交。LiDAR 的激光射线在全景坐标系中参数化，射线由两个正交齐次平面定义。射线与 2D 高斯圆盘的交点通过求解线性方程组直接计算 $(u, v)$ 坐标，然后从交点坐标精确计算距离 $r$。这避免了 3DGS 中局部仿射近似在 LiDAR 射线几何上的误差。三通道 LiDAR 渲染。
\textbf{实现补充：} 实现上还要核对输入、输出、参数、特征、初始化、训练损失、正则、渲染或推理接口，以及优化中的剪枝、分裂、采样、编码或解码步骤；这些抓手决定该方法到底改了 3DGS 管线的哪一环。
\textbf{设计动机与边界：} 相对原始 3DGS 以独立高斯和图像重建误差为主的设定，GS-LiDAR 把问题落在以下可实现的改动上：GS-LiDAR框架包含：(1) 周期振动2D高斯：使用2D高斯基元（而非3D高斯）配合周期振动属性，准确捕获驾驶场景中静态和动态元素的表面行为。2D高斯更适合建模薄表面和道路等平面结构。(2) 全景渲染：提出创新的基于显式射线-溅射交叉的全景渲染技术，配合全景LiDAR监督。不同于相机的透视投影，LiDAR使用柱面/球面投影。(3) LiDAR建模：在高斯基元中引入强度和射线丢失(ray-drop)球谐(SH)系数，增强渲染点云的真实感，模拟不同物体材质的反射特性和激光束丢失现象。(4) 静态-动态分离建模。其设计目的不是单独增加一个模块，而是用上述表示、监督或渲染约束消除原流程在目标任务中的具体失配。提出GS-LiDAR框架，使用具有周期振动特性的2D高斯基元和全景高斯溅射生成真实LiDAR点云。引入全景渲染技术配合显式射线-溅射交叉，以及强度和射线丢失球谐系数增强点云真实感。在KITTI-360和nuScenes上实现训练和渲染效率的大幅提升。适用时还应保留论文给出的边界：2D高斯对厚实体的建模能力有限。周期振动属性增加了参数量。全景渲染的分辨率受限于高斯基元的密度。
\textbf{方法信息：} 代码：\url{https://github.com/fudan-zvg/gs-lidar}；GitHub 关注度较高，约 156 stars。



% METHOD: 2503.08317
\paragraph{Uni-Gaussians（相机光栅化与 LiDAR 光追统一）（2025.03）}
Uni-Gaussians~\cite{2503_08317} 认为相机和 LiDAR 应共用高斯场景图，但渲染策略不必相同。相机 RGB 继续使用 3DGS 光栅化，因为质量和速度更合适；LiDAR 则用 Gaussian ray tracing，因为主动激光束发射和返回不能用普通光栅化准确模拟。动态驾驶场景被分成静态背景和前景对象，包括刚性车辆、SMPL 行人和可变形节点。相对 SplatAD 的统一 rasterization 倾向，Uni-Gaussians 明确按传感器物理选择不同渲染后端。
\textbf{实现细节：} Uni-Gaussians（2025.03） 提出混合方法——光栅化渲染图像数据，高斯光线追踪渲染 LiDAR 数据——统一动态驾驶场景的多传感器仿真。与 SplatAD 通过可微泼溅统一相机和 LiDAR 的渲染管线不同，Uni-Gaussians 的核心洞察是相机和 LiDAR 的成像原理本质不同（相机是被动接收光线的光栅化过程，LiDAR 是主动发射激光的光线追踪过程），强行统一到单一渲染模式会损害各自的仿真质量。因此该框架在场景表示层面统一（所有动态参与者用 2D 高斯基元建模，附加 SH 强度和 SH 光线丢弃概率属性），但在渲染层面分流：相机渲染使用标准 2D-to-2D 投影和 $\alpha$-混合光栅化；LiDAR 渲染将 2D 高斯转换为三角面片构建 BVH，通过 OptiX 光线追踪计算射线-高斯交点，逐射线前向后积分深度、强度和光线丢弃概率。场景图建模沿用 OmniRe 的静态/刚性/非刚性分解，同时解决了几何表面重建中 3D 高斯的体积表示与薄表面之间的冲突。
\textbf{实现补充：} 实现上还要核对输入、输出、参数、特征、初始化、训练损失、正则、渲染或推理接口，以及优化中的剪枝、分裂、采样、编码或解码步骤；这些抓手决定该方法到底改了 3DGS 管线的哪一环。
\textbf{方法信息：} Xiaomi 参与了这篇论文。
\textbf{设计动机与边界：} 相对原始 3DGS 以独立高斯和图像重建误差为主的设定，Uni-Gaussians 把问题落在以下可实现的改动上：方法包含三个核心部分：(1) 场景图建模：受OmniRe启发，将动态驾驶场景分解为静态背景和前景对象（刚性车辆+非刚性行人），分别用不同节点类型建模（刚性节点、SMPL节点、可变形节点）。通过图变换在时间t将前景高斯变形到世界坐标系。(2) 相机数据渲染：使用标准3DGS光栅化渲染RGB图像，因为研究表明光线追踪虽然质量可比但计算时间显著增加。(3) LiDAR数据渲染：提出高斯光线追踪方法，因为光栅化无法生成LiDAR的柱面距离图像，也不符合LiDAR激光束发射和接收的主动传感机制。光线追踪能准确模拟LiDAR的物理成像过程，避免点位置混乱。为增强几何表示，还引入表面法线约束。其设计目的不是单独增加一个模块，而是用上述表示、监督或渲染约束消除原流程在目标任务中的具体失配。首次提出统一高斯框架同时仿真相机和LiDAR数据，针对不同传感器特性选择不同渲染策略（光栅化vs光线追踪），解决传统3DGS光栅化无法准确建模LiDAR主动传感机制的问题。适用时还应保留论文给出的边界：光线追踪渲染LiDAR数据速度慢于光栅化；场景图建模增加了系统复杂度；对动态行人的非刚性建模仍有提升空间。
% METHOD: 2506.01379
\paragraph{RadarSplat（雷达高斯场与多径噪声）（2025.06）}
RadarSplat~\cite{2506_01379} 较早把 3D Gaussian field 推到自动驾驶雷达模态。它基于雷达物理方程，引入方位角和仰角天线增益，并用频谱泄漏建模模糊距离测量。噪声方面，多径 ghost targets 通过 FFT 检测并周期建模，接收器饱和用频域常数项检测，视角依赖噪声概率进入渲染。相对 LiDAR-GS，RadarSplat 的重点不是射线点云，而是雷达回波、多径和噪声的逆渲染。
\textbf{实现细节：} SplatAD（2024.11） 是首个同时实现动态场景相机和 LiDAR 实时真实感渲染的 3DGS 方法，精确建模滚动快门、LiDAR 强度和光线丢弃等传感器特定现象。方法将场景分解为静态背景、天空环境和动态物体（支持刚性和可变形物体），动态物体通过 4D 傅里叶球谐建模时变外观，滚动快门通过逐行独立的相机外参建模，LiDAR 模拟通过可微泼溅同时渲染深度、强度和光线丢弃。
\textbf{实现补充：} 进一步拆开来看，输入侧对应图像、点云、位姿、语义提示或已训练高斯资产；输出侧对应新的高斯参数、渲染图像、编辑结果、地图状态或评测指标。表示层需要说明位置、尺度、旋转、不透明度、特征、材质或 MLP 解码器怎样改；训练层需要说明损失、正则、采样、剪枝、分裂或压缩何时发生；推理层则看渲染、查询、规划或系统接口如何接到原始 3DGS / Scaffold-GS / 4DGS 管线上。
\textbf{方法信息：} 代码：\url{https://github.com/umautobots/radarsplat}。
\textbf{设计动机与边界：} 相对原始 3DGS 以独立高斯和图像重建误差为主的设定，RadarSplat 把问题落在以下可实现的改动上：核心贡献：(1) 雷达高斯场渲染公式：基于雷达物理方程，集成方位角和仰角天线增益，考虑频谱泄漏建模模糊的雷达距离测量。(2) 雷达噪声建模：对多径效应（ghost targets）使用基于FFT的检测方法并建模其周期性；对接收器饱和使用频域常数项检测；引入视角依赖的噪声概率模型。(3) 雷达逆渲染：从含噪雷达图像中分解真实目标、噪声和多径效应，支持新视角合成中恢复多径效果。(4) 雷达去噪方法：基于噪声检测的去噪策略，生成无噪占用图作为监督信号。整体框架将3DGS扩展到雷达模态，填补了该领域空白。其设计目的不是单独增加一个模块，而是用上述表示、监督或渲染约束消除原流程在目标任务中的具体失配。首次将3D高斯场应用于雷达模态，提出包含多径效应和噪声建模的雷达渲染公式，实现逼真雷达数据合成和增强3D重建，PSNR提升3.4dB、RMSE降低40%。适用时还应保留论文给出的边界：雷达分辨率和精度受限于硬件特性；对复杂城区密集多径场景效果可能受限；仅处理扫描雷达图像。
% METHOD: 2510.12901
\paragraph{SimULi（任意相机与旋转 LiDAR 实时仿真）（2025.10）}
SimULi~\cite{2510_12901} 基于 3DGUT 扩展复杂相机和 LiDAR。它把相机和 LiDAR 属性编码在独立半透明 3D Gaussian particle 集合中，并共享几何信息以缓解跨传感器标定不一致。LiDAR 渲染使用自动分块策略适配任意旋转模型，并用 Unscented Transform 估计每个高斯对光线的贡献。相对 SplatAD 的球坐标 rasterization，SimULi 更强调任意相机模型、任意旋转 LiDAR 和实时速度。
\textbf{实现细节：} SimULi（2025.10） 与 Uni-Gaussians 的混合光栅化-光线追踪策略不同，SimULi 基于 3DGUT 的无迹变换（Unscented Transform）统一处理任意相机模型和 LiDAR 的实时渲染——两者共享同一套粒子表示和泼溅管线。Baseline 与问题。现有方法要么为不同传感器维护独立表示（Uni-Gaussians 分别用光栅化和光线追踪），要么强行统一到单一渲染模式（SplatAD 统一用泼溅）。前者增加存储和同步开销，后者牺牲 LiDAR 的仿真精度。同时，相机数据（稠密 RGB）和 LiDAR 数据（稀疏深度+强度）的分布特性差异巨大，联合优化时一种模态的质量往往被另一种模态的不一致性拖累。因子化 3D 高斯表示。相机属性集 $G_c$（3 阶球谐编码视角相关颜色）和 LiDAR 属性集 $G_l$（3 阶球谐编码视角相关强度和光线丢弃概率）使用独立的粒子集但共享位置 $\mu_i$、不透明度 $o_i$ 和协方差 $\Sigma_i$。这种因子化解耦使每个传感器模态的重建质量不受另一模态的传感器不一致性影响——避免了 UniSplat 等方法中跨传感器数据不一致时必须优先某一模态质量的权衡。LiDAR 自动分块与光线剔除。
\textbf{方法信息：} NVIDIA 参与了这篇论文。
\textbf{设计动机与边界：} 相对原始 3DGS 以独立高斯和图像重建误差为主的设定，SimULi 把问题落在以下可实现的改动上：SimULi基于3DGUT构建（支持复杂相机和时间依赖效应）。核心设计包括：(1)粒子表示：相机和LiDAR属性编码在独立的半透明3D高斯粒子集合G\_c和G\_l中，各有位置、不透明度、协方差矩阵和三阶球谐系数。LiDAR高斯还编码光束强度和射线丢弃概率。(2)场景表示：采用图分解，粒子分配给动态物体或静态背景，每个物体关联3D边界框和SE(3)位姿序列（带可学习偏移）。渲染时物体高斯从局部坐标变换到世界坐标。(3)LiDAR渲染：通过自动化分块策略适配任意旋转LiDAR模型，基于光线剔除优化计算。使用Unscented Transform确定每个高斯对各光线的贡献。(4)分解锚定策略：通过分解相机和LiDAR高斯属性并共享几何信息，解决跨传感器标定不一致性，降低误差最多40%。渲染使用体渲染+alpha合成+双线性网格处理光照变化。其设计目的不是单独增加一个模块，而是用上述表示、监督或渲染约束消除原流程在目标任务中的具体失配。提出SimULi首个能实时渲染任意相机模型和LiDAR数据的方法。扩展3DGUT支持复杂相机模型和LiDAR，通过自动化分块策略和基于光线的剔除适配任意旋转LiDAR模型。设计分解式3D高斯表示和锚定策略解决跨传感器不一致性，将相机和深度误差降低最多40%。渲染速度比光线追踪快10-20倍，比先前光栅化方法快1.5-10倍。适用时还应保留论文给出的边界：分解式表示增加了存储开销；对极高质量要求的场景可能仍需后处理；LiDAR射线丢弃模型对天气条件的建模有限。
% METHOD: 2603.14763
\paragraph{LiDAR-EVS（外推视角 LiDAR 伪监督）（2026.03）}
LiDAR-EVS~\cite{2603_14763} 专门增强外推视角 LiDAR 合成。它利用世界坐标系几何不变性，从单遍扫描通过多帧 LiDAR 融合、视角变换、遮挡 warping 和强度调整生成伪外推视角点云监督。空间约束 dropout 正则进一步提升轨迹变化鲁棒性。相对 GS-LiDAR 的传感器表示，LiDAR-EVS 是一个即插即用训练补丁，目标是让 LiDAR 3DGS 在偏离原轨迹时仍可靠。
\textbf{实现细节：} LiDAR-EVS（2026.03） 提出轻量级框架增强沿未驾驶轨迹的 LiDAR 仿真。伪外推视角点云监督通过四步生成：多帧 LiDAR 融合（聚合不同时刻的点云扩展空间覆盖）、视图变换（将点云投影到外推视角的坐标系）、遮挡卷曲（处理外推视角中新暴露区域的缺失点云）、强度调整（根据新视角的激光入射角重新计算反射强度）。空间约束 dropout 正则化在外推区域随机丢弃点云，防止模型过拟合到不完整的几何。
\textbf{实现补充：} 实现上还要核对输入、输出、参数、特征、初始化、训练损失、正则、渲染或推理接口，以及优化中的剪枝、分裂、采样、编码或解码步骤；这些抓手决定该方法到底改了 3DGS 管线的哪一环。
\textbf{技术实现补充：} LiDAR-EVS（2026.03） 提出轻量级框架增强沿未驾驶轨迹的 LiDAR 仿真。伪外推视角点云监督通过四步生成：多帧 LiDAR 融合（聚合不同时刻的点云扩展空间覆盖）、视图变换（将点云投影到外推视角的坐标系）、遮挡卷曲（处理外推视角中新暴露区域的缺失点云）、强度调整（根据新视角的激光入射角重新计算反射强度）。空间约束 dropout 正则化在外推区域随机丢弃点云，防止模型过拟合到不完整的几何。
\textbf{方法信息：} NVIDIA 参与了这篇论文。
\textbf{设计动机与边界：} 相对原始 3DGS 以独立高斯和图像重建误差为主的设定，LiDAR-EVS 把问题落在以下可实现的改动上：LiDAR-EVS是轻量级即插即用框架，包含两个关键组件：(1)伪外推视角点云监督：利用世界坐标系中的几何不变性，通过多帧LiDAR融合、视角变换、遮挡卷曲和强度调整，从单遍扫描数据生成外推视角的伪监督信号；(2)空间约束dropout正则化：促进对真实驾驶中各种轨迹变化的鲁棒性。该框架可轻松扩展到多种LiDAR传感器（如16/32/64线）和神经渲染基线，最小化修改即可集成。其设计目的不是单独增加一个模块，而是用上述表示、监督或渲染约束消除原流程在目标任务中的具体失配。提出即插即用的LiDAR-EVS框架，通过伪外推视角点云监督（多帧融合+视角变换+遮挡卷曲+强度调整）和空间约束dropout正则化，无需外部多遍数据即可实现鲁棒的外推视角LiDAR合成，在三个数据集上达到SOTA。适用时还应保留论文给出的边界：仅处理静态场景元素；伪监督质量受原始LiDAR密度限制；对极端外推角度效果可能下降。

\section{3D 占用估计与感知}

占用和 BEV 感知方法不一定把 3DGS 当作最终渲染资产，而是把高斯当成 2D 到 3D 的可微桥梁：图像特征可以 splat 到 3D，高斯可以投影回图像作为训练监督，也可以转换为 voxel occupancy。
% METHOD: 2506.05563
\paragraph{VoxelSplat（训练时高斯正则占用与场景流）（2025.06）}
VoxelSplat~\cite{2506_05563} 把 3DGS 作为训练正则，而不是推理表示。从体素特征中采样高斯中心，解码动态 3DGS；语义 logits 替代 SH 系数，通过 alpha blending 渲染到相机视图，用 2D 标签监督 3D 语义。场景流学习中，用预测 flow 变换高斯位置到下一帧，再用相邻帧 2D 标签自监督。推理时不增加额外开销。相对 ST-GS，VoxelSplat 的重点是“高斯渲染损失”而不是最终高斯表示。
\textbf{实现细节：} VoxelSplat（2025.06） 利用 3DGS 作为正则化框架增强占用和场景流预测，核心思路是将占用预测网络解码出的稀疏语义 3D 高斯投影到 2D 相机视角，与输入图像的语义分割结果比较提供额外语义监督。同时使用预测的场景流建模高斯运动（将 $t$ 时刻的高斯按场景流偏移到 $t+1$ 时刻，渲染后与 $t+1$ 的观测比较），实现自监督场景流学习。整个 3DGS 正则化模块仅在训练时使用，可无缝集成到现有占用模型不增加推理时间。
\textbf{方法信息：} Baidu 参与了这篇论文。
\textbf{设计动机与边界：} 相对原始 3DGS 以独立高斯和图像重建误差为主的设定，VoxelSplat 把问题落在以下可实现的改动上：核心思想是将3DGS作为训练时的正则化工具：(1)从体素特征中通过加权点采样选择高斯中心，解码出动态3DGS；(2)语义渲染：将语义logits替代SH系数，通过alpha混合渲染到相机视图，利用2D标签监督3D语义学习；(3)场景流学习：用预测的场景流变换高斯位置预测下一帧语义场，利用相邻帧2D标签自监督。分解动态/静态物体分别处理。所有GS操作仅在训练时作为额外渲染头和损失。其设计目的不是单独增加一个模块，而是用上述表示、监督或渲染约束消除原流程在目标任务中的具体失配。提出VoxelSplat正则化框架，利用3DGS增强占用预测：(1)将稀疏语义3DGS投影到2D视图提供额外语义监督信号；(2)用预测场景流驱动高斯运动，利用相邻帧标签自监督学习场景流。训练时额外开销，推理时无额外成本。显著提升语义占用和场景流估计精度。适用时还应保留论文给出的边界：训练时增加渲染计算开销；高斯稀疏性可能影响细粒度语义；依赖场景流预测质量。
% METHOD: 2506.09417
\paragraph{ODG（双高斯占用预测）（2025.06）}
ODG~\cite{2506_09417} 把驾驶场景分为静态背景和动态物体，用 dual Gaussian queries 分别建模。每组查询包含高斯原语和特征向量，动态高斯额外建模时间相关位置偏移。层次化 Gaussian Transformer 逐步细化占据体素中心和语义类别，并用 3DGS 渲染监督把可用深度和语义标注注入像素级对齐。相对普通 BEV query，ODG 的增量是让 occupancy query 具有显式高斯几何属性。
\textbf{实现细节：} ODG（2025.06） 与 GaussianOcc 使用统一高斯查询不同，ODG 基于驾驶场景可分解为静态和动态部分的观察，设计双高斯查询更好地建模多样场景物体。与 GaussianOcc 的差异。GaussianOcc 将体素顶点直接视为高斯位置，使用固定缩放和旋转；ODG 定义两组独立查询——静态查询 $G^s$ 在固定 BEV 网格位置均匀初始化（适合捕获道路、建筑等静态结构），动态查询 $G^d$ 在检测到的 3D 边界框中心初始化（适合捕获车辆、行人的局部细节）。关键扩展是为动态查询增加 3D 检测框属性 $\mathbf{b} = [l, w, h, \theta, v_x, v_y, v_z]$（长宽高、朝向、速度），使动态高斯同时承担占用预测和物体检测的双重任务。层次粗到细预测。不同于前述方法使用固定数量的高斯，ODG 通过 $L$ 层变换器逐层增加每查询的高斯数量（$K_1 < K_2 < \cdots < K_L$），每层预测高斯中心、语义类别和（仅动态查询的）检测框属性。动静交叉注意力（DaS Attention）将两组查询的特征拼接后执行全自注意力，实现动静信息的双向增强——静态查询为动态查询提供场景上下文，动态查询为静态查询补充物体级语义。渲染监督。
\textbf{方法信息：} Qualcomm 参与了这篇论文。
\textbf{设计动机与边界：} 相对原始 3DGS 以独立高斯和图像重建误差为主的设定，ODG 把问题落在以下可实现的改动上：ODG基于对驾驶场景可分解为静态背景（道路、建筑）和动态物体（车辆、行人）的观察，设计双高斯查询机制。每组高斯查询包含一组高斯原语（位置、尺度、四元数、不透明度）和特征向量。静态高斯查询数S个、动态高斯查询数D个，分别初始化。通过注意力机制实现双查询间的特征交互。动态高斯额外建模物体运动，包含时间相关的位置偏移。层次化高斯Transformer逐步细化占据体素中心和语义类别预测。利用3DGS的实时渲染能力，引入渲染监督，结合可用深度和语义标注注入像素级对齐来增强占据学习。整体保持低推理开销。其设计目的不是单独增加一个模块，而是用上述表示、监督或渲染约束消除原流程在目标任务中的具体失配。提出分层双稀疏高斯表示方法，将驾驶场景分解为静态背景和动态物体分别建模，结合层次化高斯Transformer和3DGS渲染监督，在Occ3D-nuScenes和Occ3D-Waymo上取得SOTA结果。适用时还应保留论文给出的边界：动态高斯的运动建模依赖历史帧，对突然出现的物体处理有限；高斯数量仍需手动设定；在极端遮挡场景下性能受限。
% METHOD: 2509.16552
\paragraph{ST-GS（时空语义占用高斯）（2025.09）}
ST-GS~\cite{2509_16552} 增强视觉中心 3D semantic occupancy。Guidance-Informed Spatial Aggregation 分为 Gaussian-Guided Attention 和 View-Guided Attention，前者用高斯 3D 属性精炼局部采样，后者沿相机射线利用多视图语义连续性；Geometry-Aware Temporal Fusion 通过自运动对齐跨帧高斯嵌入，再门控融合历史信息。最后 Gaussian-to-voxel splatting 生成密集语义占用。相对 ODG 的双查询，ST-GS 更强调空间视图交互和时间一致性。
\textbf{实现细节：} ST-GS（2025.09） 与 GraphGSOcc 从图结构增强空间交互不同，ST-GS 从注意力机制的参考点设计出发，提出引导知情空间聚合（GISA）策略——通过双模式注意力（高斯引导注意力 GGA + 视图引导注意力 VGA）提供互补的 3D 空间采样偏移。双模式空间聚合。高斯引导注意力（GGA）利用每个高斯固有的均值和协方差属性生成采样偏移——在局部高斯坐标系中初始化 $M$ 个结构化 3D 网格点作为基础偏移 $P_L^G$，通过可学习缩放因子 $s^G$ 调节采样半径，再经旋转 $R^G$ 和缩放 $S^G$ 变换到感知坐标系，使偏移方向与高斯椭球的主轴对齐。视图引导注意力（VGA）沿相机视线方向生成偏移——在局部视线坐标系的 y-z 平面初始化 2D 网格偏移，根据高斯中心的方位角 $\phi$ 构建旋转矩阵 $R^V(\phi)$ 将偏移变换到感知坐标系，利用多视角重叠区域的空间-语义连续性增强跨视角信息聚合。门控空间特征聚合（GSFA）通过注意力驱动的门控 $\sigma_S \in [0,1]^{K \times M}$ 动态平衡两种偏移的贡献。几何感知时间融合。
\textbf{方法信息：} DJI 参与了这篇论文；代码：\url{https://github.com/sadrasafa/StereoGS}。
\textbf{设计动机与边界：} 相对原始 3DGS 以独立高斯和图像重建误差为主的设定，ST-GS 把问题落在以下可实现的改动上：ST-GS框架增强高斯式语义占据预测流程：(1) Guidance-Informed Spatial Aggregation (GISA)：设计双模态注意力机制，包含Gaussian-Guided Attention (GGA)利用高斯内在3D属性精炼局部特征采样，和View-Guided Attention (VGA)利用多视图的空间-语义连续性沿相机射线自适应采样。通过Gated Spatial Feature Aggregation (GSFA)模块高效融合两个分支的参考点。(2) Geometry-Aware Temporal Fusion (GATF)：利用自运动变换显式对齐跨帧高斯嵌入，通过Gated Temporal Feature Fusion (GTFF)模块选择性整合相关历史信息到当前关键帧表示中。整体流程：图像编码器提取多视角2D特征→GISA桥接2D特征与3D高斯嵌入→GATF提升时间连贯性→GS头解码为高斯原语→Gaussian-to-voxel splatting生成密集语义占据体素。其设计目的不是单独增加一个模块，而是用上述表示、监督或渲染约束消除原流程在目标任务中的具体失配。提出ST-GS框架，通过引导式空间聚合策略（GISA）和几何感知时间融合方案（GATF），分别增强高斯表示中的多视角空间交互和多帧时间一致性，在大规模nuScenes基准上实现了SOTA性能。适用时还应保留论文给出的边界：高斯原语数量可能影响计算效率；对动态场景中快速运动物体的时间建模仍有提升空间。
% METHOD: 2602.08784
\paragraph{GaussianCaR（相机-雷达 BEV 视图变换器）（2026.02）}
GaussianCaR~\cite{2602_08784} 将 Gaussian Splatting 重新定位为相机-雷达融合的通用 view transformer。Pixels-to-Gaussians 编码器从图像特征估计高斯位置、尺寸、朝向、不透明度和特征；Points-to-Gaussians 编码器把稀疏雷达点映射为 BEV 高斯。两种模态高斯进入统一 BEV 空间，经多尺度融合和 Transformer 解码输出 BEV 分割。相对 RadarSplat 的物理雷达渲染，GaussianCaR 更偏端到端感知网络中的模态对齐。
\textbf{实现细节：} GaussianCaR（2026.02） 将高斯泼溅重新用作高效通用视图变换器，将图像像素和雷达点映射到公共 BEV 表示。图像分支为每个像素预测深度分布并提升为 3D 高斯，雷达分支直接利用雷达的距离信息初始化高斯位置，两种模态的高斯在 BEV 平面上通过可微分光栅化泼溅融合为统一特征图。在 nuScenes 上 BEV 分割与当时方法持平但推理速度快 3.2 倍。
\textbf{技术实现补充：} GaussianCaR（2026.02） 将高斯泼溅重新用作高效通用视图变换器，将图像像素和雷达点映射到公共 BEV 表示。图像分支为每个像素预测深度分布并提升为 3D 高斯，雷达分支直接利用雷达的距离信息初始化高斯位置，两种模态的高斯在 BEV 平面上通过可微分光栅化泼溅融合为统一特征图。在 nuScenes 上 BEV 分割与当时方法持平但推理速度快 3.2 倍。 %=====================================================================。
\textbf{方法信息：} 代码：\url{https://github.com/santimontiel/gaussiancar}。
\textbf{设计动机与边界：} 相对原始 3DGS 以独立高斯和图像重建误差为主的设定，GaussianCaR 把问题落在以下可实现的改动上：GaussianCaR包含两个模态特定编码器：Pixels-to-Gaussians编码器使用EfficientViT-L2提取图像特征，通过粗到细策略估计每个高斯的位置（先通过深度分类对数得到粗位置，再通过偏移头精修），同时预测高斯的尺寸、朝向、不透明度和特征；Points-to-Gaussians编码器将稀疏雷达点映射为BEV高斯表示。两种模态的高斯被映射到统一BEV空间后，通过多尺度融合模块和Transformer解码器提取BEV特征，最终输出BEV分割图。该方法的核心创新在于将GS从场景重建工具重新定位为视图变换桥梁，解决了多模态输入的视图差异问题。其设计目的不是单独增加一个模块，而是用上述表示、监督或渲染约束消除原流程在目标任务中的具体失配。提出GaussianCaR，首个将Gaussian Splatting重新定位为通用视图变换器用于相机-雷达融合的端到端网络。通过Pixels-to-Gaussians和Points-to-Gaussians两个模态特定编码器，将图像像素和雷达点映射到统一的BEV表示，结合多尺度融合和Transformer解码器实现高效BEV分割。适用时还应保留论文给出的边界：方法依赖多相机系统和雷达输入，对传感器配置有一定要求；BEV分割分辨率受限于高斯基元数量和BEV网格大小。

\section{单图通用重建边界方法}

这一节放在自动驾驶章末，是为了标出“通用单图重建”和“驾驶专用场景建模”的边界。单图前馈方法可以在 KITTI 等数据上展示迁移能力，但它们通常没有显式车辆轨迹、动态参与者、地图先验或多传感器约束。读这类方法时，应把它们看成驾驶场景可借用的几何初始化或泛化基线，而不是完整自动驾驶 3DGS 系统。
% METHOD: 2401.08045
\paragraph{Forging Vision Foundation Models for Autonomous Driving（自动驾驶场景）（2024.01）}
Forging Vision Foundation Models for Autonomous Driving~\cite{2401_08045} 关注的问题是：把 3DGS 用于自动驾驶场景、仿真或感知 Forging Vision Foundation Models for Autonomous Driving 是宽召回补全中发现的论文，主要根据标题、摘要和 arXiv 类别判定其与 3DGS 的关系。
相对原始 3DGS，Forging Vision Foundation Models for Autonomous Driving 的变化不是只替换损失函数，而是把表示、初始化、约束或网络预测路径重新组织到任务流程中：该工作被召回是因为论文文本包含 Gaussian/3D Gaussian 相关表述；当前分类先确定其任务域和与 3DGS 的相关性，后续证据流水线再补代码、企业和引用信息。
\textbf{实现细节：} 具体实现可以按输入、核心模块和渲染/优化路径来读。输入端首先取得任务所需的图像、相机、文本、时间序列、语义或几何线索，并把这些线索转换成可以约束高斯中心、尺度、旋转、不透明度、颜色或附加特征的条件。核心结构中，方法使用的主要机制是：从摘要层面看，Forging Vision Foundation Models for Autonomous Driving 围绕“把 3DGS 用于自动驾驶场景、仿真或感知”展开。若其正文或代码后续可读，应进一步核对它具体改动的是表示、初始化、渲染、优化、损失、动态建模、压缩或应用接口中的哪一环。 训练或推理时，这些模块通常先给出初始化或中间特征，再通过可微 Gaussian rasterization、重建损失、感知/语义约束、几何正则或任务损失回传到高斯参数；如果方法包含前馈网络，则网络输出的是高斯属性、变形场、压缩码、语义特征或控制变量，而不是只输出最终 2D 图像。渲染路径仍以显式高斯投影、深度排序和 alpha 合成为基础，只是在采样、属性预测、局部编辑、动态变形、压缩编码或下游任务接口处加入额外约束。这样写入本章时，它与前后方法的关系是：相对原始 3DGS 和本章已有锚点方法，它可以理解为 3DGS extends\_or\_uses 宽召回分类阶段的关系判断，后续需要按论文方法部分复核。
\textbf{证据与边界：} 本轮召回分类尚未逐项核对实验表格；后续证据补全阶段需要读取论文源码/PDF或项目页确认数据集、指标和对比方法。 该条来自宽召回批次；源码下载受 arXiv e-print 网络限制，部分论文暂时只能依据标题、摘要和类别分类，需要在入可信书前做人工/LLM精读复核。 本书只采用本地 paper pack、数据库分类、源码/引用登记和论文摘要中可核验的信息；没有本地证据支撑的速度、排行榜或开源完整性不在正文中强化。\textbf{方法信息：} 公开链接：\url{https://github.com/zhanghm1995/forge\_vfm4ad}，完整性按本地证据保守处理；代码关注：zhanghm1995/forge\_vfm4ad has 272 stars (9.1/month, domain-relative cutoff 13.0/month)。
\textbf{设计动机与边界：} 相对原始 3DGS 以独立高斯和图像重建误差为主的设定，Forging Vision Foundation Models for Autonomous Driving 把问题落在以下可实现的改动上：从摘要层面看，Forging Vision Foundation Models for Autonomous Driving 围绕“把 3DGS 用于自动驾驶场景、仿真或感知”展开。若其正文或代码后续可读，应进一步核对它具体改动的是表示、初始化、渲染、优化、损失、动态建模、压缩或应用接口中的哪一环。其设计目的不是单独增加一个模块，而是用上述表示、监督或渲染约束消除原流程在目标任务中的具体失配。该工作被召回是因为论文文本包含 Gaussian/3D Gaussian 相关表述；当前分类先确定其任务域和与 3DGS 的相关性，后续证据流水线再补代码、企业和引用信息。适用时还应保留论文给出的边界：该条来自宽召回批次；源码下载受 arXiv e-print 网络限制，部分论文暂时只能依据标题、摘要和类别分类，需要在入可信书前做人工/LLM精读复核。
% METHOD: 2406.04343
\paragraph{Flash3D（单图分层高斯前馈重建）（2024.06）}
Flash3D~\cite{2406_04343} 不是自动驾驶专用方法，但在 KITTI 等驾驶数据上体现出单图泛化价值。它从 UniDepth 这类基础深度模型出发，第一层高斯放在预测深度处，额外高斯层通过学习空间偏移处理遮挡和截断区域。U-Net 与 ResNet 预测每像素高斯参数，训练用输入图、目标图和相对位姿的渲染损失。该方法提供代码，并且引用较多。相对 DrivingForward 的多帧环视输入，Flash3D 是更通用的单图前馈边界：质量受单图歧义限制，但结构简单、迁移性强。

\textbf{实现细节：} Flash3D从单目深度估计基础模型（UniDepth）出发，扩展为完整3D场景重建器。核心是前馈高斯泼溅架构：(1) 分层高斯预测：第一层高斯放置在预测深度处，额外K-1层高斯通过学习空间偏移量处理遮挡和截断区域，每像素预测 $C=12+3(L+1)^{}2$ 个参数。(2) 深度先验：利用预训练深度估计器提供强几何先验。(3) 网络架构：基于U-Net+ResNet-50预测每像素高斯参数。训练使用三元组（输入图、目标图、相对位姿）最小化渲染损失。关键创新是多层高斯表示，解决场景重建中缺乏背景像素池的问题。 提出Flash3D方法，基于基础深度模型扩展为完整的3D形状和外观重建器。核心创新是分层高斯预测：在预测深度处放置第一层高斯，再添加空间偏移的额外高斯层处理遮挡和截断区域。单GPU一天可完成训练，在RealEstate10k上达到SOTA，迁移到NYU和KITTI上大幅超越竞争对手。 单目重建对极端遮挡区域仍有限制；多层高斯数量增加内存消耗；对相机位姿估计有一定依赖。
\textbf{设计动机与边界：} 相对原始 3DGS 以独立高斯和图像重建误差为主的设定，Flash3D 把问题落在以下可实现的改动上：Flash3D从单目深度估计基础模型（UniDepth）出发，扩展为完整3D场景重建器。核心是前馈高斯泼溅架构：(1) 分层高斯预测：第一层高斯放置在预测深度处，额外K-1层高斯通过学习空间偏移量处理遮挡和截断区域，每像素预测C=12+3(L+1)\^{}2个参数。(2) 深度先验：利用预训练深度估计器提供强几何先验。(3) 网络架构：基于U-Net+ResNet-50预测每像素高斯参数。训练使用三元组（输入图、目标图、相对位姿）最小化渲染损失。关键创新是多层高斯表示，解决场景重建中缺乏背景像素池的问题。其设计目的不是单独增加一个模块，而是用上述表示、监督或渲染约束消除原流程在目标任务中的具体失配。提出Flash3D方法，基于基础深度模型扩展为完整的3D形状和外观重建器。核心创新是分层高斯预测：在预测深度处放置第一层高斯，再添加空间偏移的额外高斯层处理遮挡和截断区域。单GPU一天可完成训练，在RealEstate10k上达到SOTA，迁移到NYU和KITTI上大幅超越竞争对手。适用时还应保留论文给出的边界：单目重建对极端遮挡区域仍有限制；多层高斯数量增加内存消耗；对相机位姿估计有一定依赖。
\textbf{方法信息：} 代码：\url{https://github.com/flashvision/flash3d}；引用很多，Semantic Scholar 约 128 次。

% METHOD: 2411.12452
\paragraph{GaussianPretrain（自动驾驶场景）（2024.11）}
GaussianPretrain~\cite{2411_12452} 关注的问题是：把 3DGS 用于自动驾驶场景、仿真或感知 GaussianPretrain 是宽召回补全中发现的论文，主要根据标题、摘要和 arXiv 类别判定其与 3DGS 的关系。
相对原始 3DGS，GaussianPretrain 的变化不是只替换损失函数，而是把表示、初始化、约束或网络预测路径重新组织到任务流程中：该工作被召回是因为论文文本包含 Gaussian/3D Gaussian 相关表述；当前分类先确定其任务域和与 3DGS 的相关性，后续证据流水线再补代码、企业和引用信息。
\textbf{实现细节：} 具体实现可以按输入、核心模块和渲染/优化路径来读。输入端首先取得任务所需的图像、相机、文本、时间序列、语义或几何线索，并把这些线索转换成可以约束高斯中心、尺度、旋转、不透明度、颜色或附加特征的条件。核心结构中，方法使用的主要机制是：从摘要层面看，GaussianPretrain 围绕“把 3DGS 用于自动驾驶场景、仿真或感知”展开。若其正文或代码后续可读，应进一步核对它具体改动的是表示、初始化、渲染、优化、损失、动态建模、压缩或应用接口中的哪一环。 训练或推理时，这些模块通常先给出初始化或中间特征，再通过可微 Gaussian rasterization、重建损失、感知/语义约束、几何正则或任务损失回传到高斯参数；如果方法包含前馈网络，则网络输出的是高斯属性、变形场、压缩码、语义特征或控制变量，而不是只输出最终 2D 图像。渲染路径仍以显式高斯投影、深度排序和 alpha 合成为基础，只是在采样、属性预测、局部编辑、动态变形、压缩编码或下游任务接口处加入额外约束。这样写入本章时，它与前后方法的关系是：相对原始 3DGS 和本章已有锚点方法，它可以理解为 3DGS extends\_or\_uses 宽召回分类阶段的关系判断，后续需要按论文方法部分复核。
\textbf{证据与边界：} 本轮召回分类尚未逐项核对实验表格；后续证据补全阶段需要读取论文源码/PDF或项目页确认数据集、指标和对比方法。 该条来自宽召回批次；源码下载受 arXiv e-print 网络限制，部分论文暂时只能依据标题、摘要和类别分类，需要在入可信书前做人工/LLM精读复核。 本书只采用本地 paper pack、数据库分类、源码/引用登记和论文摘要中可核验的信息；没有本地证据支撑的速度、排行榜或开源完整性不在正文中强化。\textbf{方法信息：} 代码：\url{https://github.com/public-bots/gaussianpretrain}；代码关注：public-bots/gaussianpretrain has 112 stars (5.7/month, domain-relative cutoff 13.0/month)。
\textbf{设计动机与边界：} 相对原始 3DGS 以独立高斯和图像重建误差为主的设定，GaussianPretrain 把问题落在以下可实现的改动上：从摘要层面看，GaussianPretrain 围绕“把 3DGS 用于自动驾驶场景、仿真或感知”展开。若其正文或代码后续可读，应进一步核对它具体改动的是表示、初始化、渲染、优化、损失、动态建模、压缩或应用接口中的哪一环。其设计目的不是单独增加一个模块，而是用上述表示、监督或渲染约束消除原流程在目标任务中的具体失配。该工作被召回是因为论文文本包含 Gaussian/3D Gaussian 相关表述；当前分类先确定其任务域和与 3DGS 的相关性，后续证据流水线再补代码、企业和引用信息。适用时还应保留论文给出的边界：该条来自宽召回批次；源码下载受 arXiv e-print 网络限制，部分论文暂时只能依据标题、摘要和类别分类，需要在入可信书前做人工/LLM精读复核。
% METHOD: 2412.10371
\paragraph{GaussianAD（自动驾驶场景）（2024.12）}
GaussianAD~\cite{2412_10371} 关注的问题是：把 3DGS 用于自动驾驶场景、仿真或感知 GaussianAD 是宽召回补全中发现的论文，主要根据标题、摘要和 arXiv 类别判定其与 3DGS 的关系。
相对原始 3DGS，GaussianAD 的变化不是只替换损失函数，而是把表示、初始化、约束或网络预测路径重新组织到任务流程中：该工作被召回是因为论文文本包含 Gaussian/3D Gaussian 相关表述；当前分类先确定其任务域和与 3DGS 的相关性，后续证据流水线再补代码、企业和引用信息。
\textbf{实现细节：} 具体实现可以按输入、核心模块和渲染/优化路径来读。输入端首先取得任务所需的图像、相机、文本、时间序列、语义或几何线索，并把这些线索转换成可以约束高斯中心、尺度、旋转、不透明度、颜色或附加特征的条件。核心结构中，方法使用的主要机制是：从摘要层面看，GaussianAD 围绕“把 3DGS 用于自动驾驶场景、仿真或感知”展开。若其正文或代码后续可读，应进一步核对它具体改动的是表示、初始化、渲染、优化、损失、动态建模、压缩或应用接口中的哪一环。 训练或推理时，这些模块通常先给出初始化或中间特征，再通过可微 Gaussian rasterization、重建损失、感知/语义约束、几何正则或任务损失回传到高斯参数；如果方法包含前馈网络，则网络输出的是高斯属性、变形场、压缩码、语义特征或控制变量，而不是只输出最终 2D 图像。渲染路径仍以显式高斯投影、深度排序和 alpha 合成为基础，只是在采样、属性预测、局部编辑、动态变形、压缩编码或下游任务接口处加入额外约束。这样写入本章时，它与前后方法的关系是：相对原始 3DGS 和本章已有锚点方法，它可以理解为 3DGS extends\_or\_uses 宽召回分类阶段的关系判断，后续需要按论文方法部分复核。
\textbf{证据与边界：} 本轮召回分类尚未逐项核对实验表格；后续证据补全阶段需要读取论文源码/PDF或项目页确认数据集、指标和对比方法。 该条来自宽召回批次；源码下载受 arXiv e-print 网络限制，部分论文暂时只能依据标题、摘要和类别分类，需要在入可信书前做人工/LLM精读复核。 本书只采用本地 paper pack、数据库分类、源码/引用登记和论文摘要中可核验的信息；没有本地证据支撑的速度、排行榜或开源完整性不在正文中强化。\textbf{方法信息：} 代码：\url{https://github.com/wzzheng/gaussianad}；代码关注：wzzheng/gaussianad has 128 stars (6.8/month, domain-relative cutoff 13.0/month)。
\textbf{设计动机与边界：} 相对原始 3DGS 以独立高斯和图像重建误差为主的设定，GaussianAD 把问题落在以下可实现的改动上：从摘要层面看，GaussianAD 围绕“把 3DGS 用于自动驾驶场景、仿真或感知”展开。若其正文或代码后续可读，应进一步核对它具体改动的是表示、初始化、渲染、优化、损失、动态建模、压缩或应用接口中的哪一环。其设计目的不是单独增加一个模块，而是用上述表示、监督或渲染约束消除原流程在目标任务中的具体失配。该工作被召回是因为论文文本包含 Gaussian/3D Gaussian 相关表述；当前分类先确定其任务域和与 3DGS 的相关性，后续证据流水线再补代码、企业和引用信息。适用时还应保留论文给出的边界：该条来自宽召回批次；源码下载受 arXiv e-print 网络限制，部分论文暂时只能依据标题、摘要和类别分类，需要在入可信书前做人工/LLM精读复核。
% METHOD: 2412.10373
\paragraph{GaussianWorld（自动驾驶场景）（2024.12）}
GaussianWorld~\cite{2412_10373} 关注的问题是：把 3DGS 用于自动驾驶场景、仿真或感知 GaussianWorld 是宽召回补全中发现的论文，主要根据标题、摘要和 arXiv 类别判定其与 3DGS 的关系。
相对原始 3DGS，GaussianWorld 的变化不是只替换损失函数，而是把表示、初始化、约束或网络预测路径重新组织到任务流程中：该工作被召回是因为论文文本包含 Gaussian/3D Gaussian 相关表述；当前分类先确定其任务域和与 3DGS 的相关性，后续证据流水线再补代码、企业和引用信息。
\textbf{实现细节：} 具体实现可以按输入、核心模块和渲染/优化路径来读。输入端首先取得任务所需的图像、相机、文本、时间序列、语义或几何线索，并把这些线索转换成可以约束高斯中心、尺度、旋转、不透明度、颜色或附加特征的条件。核心结构中，方法使用的主要机制是：从摘要层面看，GaussianWorld 围绕“把 3DGS 用于自动驾驶场景、仿真或感知”展开。若其正文或代码后续可读，应进一步核对它具体改动的是表示、初始化、渲染、优化、损失、动态建模、压缩或应用接口中的哪一环。 训练或推理时，这些模块通常先给出初始化或中间特征，再通过可微 Gaussian rasterization、重建损失、感知/语义约束、几何正则或任务损失回传到高斯参数；如果方法包含前馈网络，则网络输出的是高斯属性、变形场、压缩码、语义特征或控制变量，而不是只输出最终 2D 图像。渲染路径仍以显式高斯投影、深度排序和 alpha 合成为基础，只是在采样、属性预测、局部编辑、动态变形、压缩编码或下游任务接口处加入额外约束。这样写入本章时，它与前后方法的关系是：相对原始 3DGS 和本章已有锚点方法，它可以理解为 3DGS extends\_or\_uses 宽召回分类阶段的关系判断，后续需要按论文方法部分复核。
\textbf{证据与边界：} 本轮召回分类尚未逐项核对实验表格；后续证据补全阶段需要读取论文源码/PDF或项目页确认数据集、指标和对比方法。 该条来自宽召回批次；源码下载受 arXiv e-print 网络限制，部分论文暂时只能依据标题、摘要和类别分类，需要在入可信书前做人工/LLM精读复核。 本书只采用本地 paper pack、数据库分类、源码/引用登记和论文摘要中可核验的信息；没有本地证据支撑的速度、排行榜或开源完整性不在正文中强化。\textbf{方法信息：} 代码：\url{https://github.com/zuosc19/gaussianworld}；引用信号：75 citations (4.0/month, domain-relative cutoff 4.0/month, semantic\_scholar)；代码关注：zuosc19/gaussianworld has 167 stars (8.8/month, domain-relative cutoff 13.0/month)。
\textbf{设计动机与边界：} 相对原始 3DGS 以独立高斯和图像重建误差为主的设定，GaussianWorld 把问题落在以下可实现的改动上：从摘要层面看，GaussianWorld 围绕“把 3DGS 用于自动驾驶场景、仿真或感知”展开。若其正文或代码后续可读，应进一步核对它具体改动的是表示、初始化、渲染、优化、损失、动态建模、压缩或应用接口中的哪一环。其设计目的不是单独增加一个模块，而是用上述表示、监督或渲染约束消除原流程在目标任务中的具体失配。该工作被召回是因为论文文本包含 Gaussian/3D Gaussian 相关表述；当前分类先确定其任务域和与 3DGS 的相关性，后续证据流水线再补代码、企业和引用信息。适用时还应保留论文给出的边界：该条来自宽召回批次；源码下载受 arXiv e-print 网络限制，部分论文暂时只能依据标题、摘要和类别分类，需要在入可信书前做人工/LLM精读复核。
% METHOD: 2412.13193
\paragraph{GaussTR（自动驾驶场景）（2024.12）}
GaussTR~\cite{2412_13193} 关注的问题是：把 3DGS 用于自动驾驶场景、仿真或感知 GaussTR 是宽召回补全中发现的论文，主要根据标题、摘要和 arXiv 类别判定其与 3DGS 的关系。
作为 supplementary 类条目，GaussTR 不应被写成单纯的新渲染算子。它更像是把 3DGS 放进 自动驾驶场景 的协议、系统或应用边界里：该工作被召回是因为论文文本包含 Gaussian/3D Gaussian 相关表述；当前分类先确定其任务域和与 3DGS 的相关性，后续证据流水线再补代码、企业和引用信息。
\textbf{实现细节：} 具体实现可以按输入、核心模块和渲染/优化路径来读。输入端首先取得任务所需的图像、相机、文本、时间序列、语义或几何线索，并把这些线索转换成可以约束高斯中心、尺度、旋转、不透明度、颜色或附加特征的条件。核心结构中，方法使用的主要机制是：从摘要层面看，GaussTR 围绕“把 3DGS 用于自动驾驶场景、仿真或感知”展开。若其正文或代码后续可读，应进一步核对它具体改动的是表示、初始化、渲染、优化、损失、动态建模、压缩或应用接口中的哪一环。 训练或推理时，这些模块通常先给出初始化或中间特征，再通过可微 Gaussian rasterization、重建损失、感知/语义约束、几何正则或任务损失回传到高斯参数；如果方法包含前馈网络，则网络输出的是高斯属性、变形场、压缩码、语义特征或控制变量，而不是只输出最终 2D 图像。渲染路径仍以显式高斯投影、深度排序和 alpha 合成为基础，只是在采样、属性预测、局部编辑、动态变形、压缩编码或下游任务接口处加入额外约束。这样写入本章时，它与前后方法的关系是：相对原始 3DGS 和本章已有锚点方法，它可以理解为 3DGS extends\_or\_uses 宽召回分类阶段的关系判断，后续需要按论文方法部分复核。
\textbf{证据与边界：} 本轮召回分类尚未逐项核对实验表格；后续证据补全阶段需要读取论文源码/PDF或项目页确认数据集、指标和对比方法。 该条来自宽召回批次；源码下载受 arXiv e-print 网络限制，部分论文暂时只能依据标题、摘要和类别分类，需要在入可信书前做人工/LLM精读复核。 本书只采用本地 paper pack、数据库分类、源码/引用登记和论文摘要中可核验的信息；没有本地证据支撑的速度、排行榜或开源完整性不在正文中强化。\textbf{方法信息：} 代码：\url{https://github.com/hustvl/gausstr}；引用信号：51 citations (2.7/month, domain-relative cutoff 4.0/month, semantic\_scholar)；代码关注：hustvl/gausstr has 217 stars (11.6/month, domain-relative cutoff 13.0/month)。
\textbf{设计动机与边界：} 相对原始 3DGS 以独立高斯和图像重建误差为主的设定，GaussTR 把问题落在以下可实现的改动上：从摘要层面看，GaussTR 围绕“把 3DGS 用于自动驾驶场景、仿真或感知”展开。若其正文或代码后续可读，应进一步核对它具体改动的是表示、初始化、渲染、优化、损失、动态建模、压缩或应用接口中的哪一环。其设计目的不是单独增加一个模块，而是用上述表示、监督或渲染约束消除原流程在目标任务中的具体失配。该工作被召回是因为论文文本包含 Gaussian/3D Gaussian 相关表述；当前分类先确定其任务域和与 3DGS 的相关性，后续证据流水线再补代码、企业和引用信息。适用时还应保留论文给出的边界：该条来自宽召回批次；源码下载受 arXiv e-print 网络限制，部分论文暂时只能依据标题、摘要和类别分类，需要在入可信书前做人工/LLM精读复核。
% METHOD: 2512.23180
\paragraph{GaussianDWM（自动驾驶场景）（2025.12）}
GaussianDWM~\cite{2512_23180} 关注的问题是：把 3DGS 用于自动驾驶场景、仿真或感知 GaussianDWM 是宽召回补全中发现的论文，主要根据标题、摘要和 arXiv 类别判定其与 3DGS 的关系。
相对原始 3DGS，GaussianDWM 的变化不是只替换损失函数，而是把表示、初始化、约束或网络预测路径重新组织到任务流程中：该工作被召回是因为论文文本包含 Gaussian/3D Gaussian 相关表述；当前分类先确定其任务域和与 3DGS 的相关性，后续证据流水线再补代码、企业和引用信息。
\textbf{实现细节：} 具体实现可以按输入、核心模块和渲染/优化路径来读。输入端首先取得任务所需的图像、相机、文本、时间序列、语义或几何线索，并把这些线索转换成可以约束高斯中心、尺度、旋转、不透明度、颜色或附加特征的条件。核心结构中，方法使用的主要机制是：从摘要层面看，GaussianDWM 围绕“把 3DGS 用于自动驾驶场景、仿真或感知”展开。若其正文或代码后续可读，应进一步核对它具体改动的是表示、初始化、渲染、优化、损失、动态建模、压缩或应用接口中的哪一环。 训练或推理时，这些模块通常先给出初始化或中间特征，再通过可微 Gaussian rasterization、重建损失、感知/语义约束、几何正则或任务损失回传到高斯参数；如果方法包含前馈网络，则网络输出的是高斯属性、变形场、压缩码、语义特征或控制变量，而不是只输出最终 2D 图像。渲染路径仍以显式高斯投影、深度排序和 alpha 合成为基础，只是在采样、属性预测、局部编辑、动态变形、压缩编码或下游任务接口处加入额外约束。这样写入本章时，它与前后方法的关系是：相对原始 3DGS 和本章已有锚点方法，它可以理解为 3DGS extends\_or\_uses 宽召回分类阶段的关系判断，后续需要按论文方法部分复核。
\textbf{证据与边界：} 本轮召回分类尚未逐项核对实验表格；后续证据补全阶段需要读取论文源码/PDF或项目页确认数据集、指标和对比方法。 该条来自宽召回批次；源码下载受 arXiv e-print 网络限制，部分论文暂时只能依据标题、摘要和类别分类，需要在入可信书前做人工/LLM精读复核。 本书只采用本地 paper pack、数据库分类、源码/引用登记和论文摘要中可核验的信息；没有本地证据支撑的速度、排行榜或开源完整性不在正文中强化。\textbf{方法信息：} 代码：\url{https://github.com/dtc111111/gaussiandwm}。


自动驾驶 3DGS 的主线可以概括为五层。第一层是复合场景表示：DrivingGaussian、Street Gaussians、AutoSplat 和 OmniRe 把背景、动态 actor、天空和语义结构拆开。第二层是外推视角：VEGS、SGD、ReconDreamer、FaithFusion、RGE-GS 和 SetDiff 都在回答偏离行车轨迹后如何不崩。第三层是动态建模：DeSiRe-GS、S3Gaussian、SplatFlow、BezierGS、UnIRe 和 IDSplat 在减少 3D 标注依赖的同时恢复实例和运动。第四层是仿真与传感器：HUGSIM、CRUISE、SplatAD、GS-LiDAR、LiDAR-GS、RadarSplat 和 SimULi 让高斯场输出相机、LiDAR、雷达和闭环观测。第五层是感知接口：ODG、ST-GS、VoxelSplat 和 GaussianCaR 把高斯从渲染表示变成 occupancy、BEV 或伪标签训练工具。因而，自动驾驶中的 3DGS 不是单一重建算法，而是一套可编辑、可仿真、可外推、可监督感知模型的显式三维底座。

\textbf{设计动机与边界：} 相对原始 3DGS 以独立高斯和图像重建误差为主的设定，GaussianDWM 把问题落在以下可实现的改动上：从摘要层面看，GaussianDWM 围绕“把 3DGS 用于自动驾驶场景、仿真或感知”展开。若其正文或代码后续可读，应进一步核对它具体改动的是表示、初始化、渲染、优化、损失、动态建模、压缩或应用接口中的哪一环。其设计目的不是单独增加一个模块，而是用上述表示、监督或渲染约束消除原流程在目标任务中的具体失配。该工作被召回是因为论文文本包含 Gaussian/3D Gaussian 相关表述；当前分类先确定其任务域和与 3DGS 的相关性，后续证据流水线再补代码、企业和引用信息。适用时还应保留论文给出的边界：该条来自宽召回批次；源码下载受 arXiv e-print 网络限制，部分论文暂时只能依据标题、摘要和类别分类，需要在入可信书前做人工/LLM精读复核。
\textbf{尾项实现澄清：} GaussianDWM 的动机是常规驾驶世界模型缺乏三维理解，点云或 BEV 特征又难同语言精确对齐。它给每个高斯写入语言特征完成早期视觉-语言对齐，再以任务感知的语言引导采样去除冗余高斯并将紧凑三维 token 输入 LLM；生成侧以高层语言条件和低层图像条件共同驱动多模态生成。高斯表示因此既是理解接口也是生成条件，而不仅是渲染资产；摘要未公开具体损失配比，可信版仅保留该双条件和采样流程。
